% Written by a model for Adit. Not Adit's words.
% Build: content\teach\build.cmd (xelatex from TinyTeX, two passes).
\documentclass[11pt,a4paper]{article}
\usepackage{fontspec}
\setmainfont{DejaVu Serif}
\setsansfont{DejaVu Sans}
\setmonofont{DejaVu Sans Mono}
\usepackage[margin=2.4cm]{geometry}
\usepackage{amsmath,amssymb,amsthm}
\usepackage{booktabs,array,longtable}
\usepackage{enumitem}
\usepackage{titlesec}
\usepackage{fancyhdr}
\usepackage[hidelinks]{hyperref}
\setlength{\parskip}{0.45em}
\setlength{\parindent}{0pt}
\titleformat{\section}{\large\bfseries}{\thesection}{0.6em}{}
\titleformat{\subsection}{\normalsize\bfseries}{\thesubsection}{0.6em}{}
\pagestyle{fancy}\fancyhf{}
\rhead{\small Putnam. Written by a model for Adit; not Adit's words.}
\cfoot{\small\thepage}
\theoremstyle{plain}
\newtheorem{theorem}{Theorem}[section]
\newtheorem{lemma}[theorem]{Lemma}
\newtheorem{proposition}[theorem]{Proposition}
\theoremstyle{definition}
\newtheorem{definition}[theorem]{Definition}
\newtheorem{example}[theorem]{Example}
\theoremstyle{remark}
\newtheorem*{trap}{Trap}
\newtheorem*{firstmoves}{First moves}
\newtheorem*{recognise}{Recognise it}
\newtheorem*{hides}{Where it hides}
\newenvironment{ladder}{\paragraph{Ladder.}\begin{enumerate}[nosep,leftmargin=*]}{\end{enumerate}}
\newcommand{\putnam}[1]{\textsf{#1}}
\newcommand{\seminar}[1]{\textsf{#1}}
\newcommand{\Z}{\mathbb{Z}}
\newcommand{\R}{\mathbb{R}}
\newcommand{\Q}{\mathbb{Q}}
\newcommand{\C}{\mathbb{C}}
\newcommand{\N}{\mathbb{N}}
\newcommand{\E}{\mathbb{E}}
\DeclareMathOperator{\ord}{ord}
\DeclareMathOperator{\rank}{rank}
\DeclareMathOperator{\tr}{tr}
\DeclareMathOperator{\lcm}{lcm}

\title{Putnam\\[0.3em]\large Techniques, ladders and the exam}
\author{Written by a model for Adit. Not Adit's words.}
\date{Second iteration, 2026-11-07, for the sitting on 2026-12-05}

\begin{document}
\maketitle

\begin{sloppypar}
The summary of this document is \texttt{content/teach/putnam.md} in the vault; marks go there, section by section, and the next teach pass rewrites both files. No competition problem is reproduced anywhere in this document. Problems are named by source and number: \putnam{2021 A1} is a Putnam problem, in Kiran Kedlaya's archive at \url{https://kskedlaya.org/putnam-archive/} (paper \texttt{2021.pdf}, solutions \texttt{2021s.pdf}); \seminar{PS 9 \#14} is problem 14 of set 9 of MIT 18.A34, Fall 2018, by Yufei Zhao, at \url{https://ocw.mit.edu/courses/18-a34-mathematical-problem-solving-putnam-seminar-fall-2018/pages/assignments/} (CC BY-NC-SA 4.0); \seminar{Pigeonhole}, \seminar{Limits}, \seminar{Roots} and \seminar{Floor} are the four sets from that seminar's 2007 edition. Every example in this document is invented for it. The topic tags on Putnam problems are a model's reading of the statements, not of the solutions.
\end{sloppypar}

\tableofcontents
\newpage

\section{What the exam is and where the points are}

The Putnam is one Saturday, and from 2026 it has a new shape. Four sessions of 90 minutes, A, B, C and D, with three problems each: A1 to A3, B1 to B3, C1 to C3 and D1 to D3. Until 2025 it was two sessions of three hours with six problems each, A1 to A6 and B1 to B6, and every past paper named in this document has that older shape. The total is unchanged: twelve problems, each marked out of 10, so the paper is still out of 120. Every site sits each session at the same moment. In the Eastern time zone the sessions start at 11:00, 12:45, 16:00 and 17:45, so there are fifteen minutes between A and B, an hour and three quarters between B and C, fifteen minutes between C and D, and the day ends at 19:15. The source is the MAA's policy document for the 2026 competition, in the version updated 2026-08-31.

What the policy does not say is how difficulty is spread. In the old papers the problems of a session ran roughly from most approachable to hardest, with A1 and B1 meant for most contestants and A6 and B6 meant to be solved by very few people in the country. That each new session runs from an approachable first problem to a hard third is likely and not stated, and nobody outside the committee has seen a paper in the new shape. Treat the order as a hint and the reading of all three problems as the evidence.

\subsection*{How a problem is scored}

The scoring is close to binary. Almost every score given is 0, 1, 9 or 10. The reasons are worth knowing, because they decide how to spend a session. The 2026 policy says only that each problem is scored from 0 to 10 by at least two graders working blind, and that scores are final. The pattern below is the long practice of the old format, and nothing in the policy suggests it has changed.

\begin{itemize}[nosep,leftmargin=*]
\item 10: a complete and correct argument.
\item 9: a complete argument with a small, repairable gap. A missed equality case, a base case not written, a symbol not defined but clear from context.
\item 1 or 2: a correct idea with no proof, or a proof of a special case, or the right answer with a wrong or missing reason.
\item 0: everything else, including a page of true statements that do not reach the conclusion.
\end{itemize}

There is nearly nothing in the middle. A half-finished solution is worth about as much as a blank page. That single fact drives the strategy in this document: finishing one problem beats starting three.

\subsection*{What a score means}

The median total most years is a single digit, often 0, 1 or 2. In a typical year, one fully solved problem (10 points, or 9) puts a contestant above the median. Two fully solved problems, around 20 points, are usually in the top quarter of the roughly four thousand people who sit it. Three are near the top ten percent. These thresholds move a little from year to year with the paper's difficulty, but not by much. They come from the two-session papers. The 2026 paper is the first in the new shape and its thresholds are not known in advance. The argument does not depend on them: a finished problem is worth ten started ones under either shape.

So the realistic target for a first sitting is two problems solved and written cleanly across the four sessions, with a third attempted wherever a session leaves time. Under the old shape that target was reached almost entirely in the first two slots of each session.

\subsection*{Where the points are}

Across 2016 to 2025 there were forty problems in the slots A1, A2, B1 and B2 of the old papers, the four approachable slots. The data file behind this document tags each by the kind of object and method its statement points to; most problems carry two tags, so the forty problems carry seventy tags. The counts are in Table~\ref{tab:tags}.

\begin{table}[ht]
\centering
\small
\begin{tabular}{lr}
\toprule
Tag & Problems among A1, A2, B1, B2, 2016 to 2025 \\
\midrule
Congruences and number theory & 13 \\
Sums, integrals and limits & 10 \\
Counting and configurations & 10 \\
Polynomials and roots & 9 \\
Geometry & 6 \\
Inequalities & 5 \\
Linear algebra & 3 \\
Recurrences & 3 \\
Pigeonhole and extremal & 3 \\
Invariants and monovariants & 3 \\
Hidden symmetry & 2 \\
Induction & 1 \\
Games & 1 \\
Probability & 1 \\
Generating functions & 0 \\
Abstract algebra & 0 \\
Functional equations & 0 \\
\bottomrule
\end{tabular}
\caption{Tag counts in the four easy slots of the old papers. Forty problems, seventy tags. The tags are a model's reading of the statements, not of the solutions.}
\label{tab:tags}
\end{table}

Four topics carry more than half the tags: number theory, sums and limits, counting, and polynomials. Those four are the four weekly ladders in \secref{sec:ladders}{the four ladders}, in that order. Geometry and inequalities together take another eleven; they are second rails rather than ladders, because the seminar has thin material on geometry and because inequalities tend to appear inside a limit or a polynomial problem rather than alone.

The techniques at the bottom of the table are not useless. Pigeonhole, invariants and induction appear under other tags far more often than their own counts suggest, which is why each technique section ends with a note on where it hides. But they are not where the first half of a session goes.

The new paper has four sessions where the old one had four approachable slots, and the natural guess is that each session opens with a problem of that kind. That is a guess; the policy says nothing about it. What carries over without guessing is the table: ten years of what the committee chooses for the problems most contestants can reach.

\subsection*{What this means for the month}

Four weeks, four ladders, in the order of the table. Each ladder climbs one technique through its easy slots and then into one harder problem, so that the easy slots feel easy on the day. Two timed half-papers, on 2026-11-28 and 2026-12-02, each cut into two 90-minute sittings of three problems, rehearse the new shape. The target on 2026-12-05 is two clean solutions across the four sessions, and the habit of reading all three problems of a session before touching any of them.


\section{Reading a problem, and the call}
\label{sec:reading}

A Putnam problem is an analogy puzzle at scale. Find the relevant information, induce the rule, apply it. The techniques in the later sections are the rules; this section is about the finding, and about the one habit that turns a month of problems into evidence: the call.

\subsection*{Three questions before any technique}

Read the statement twice. Then answer three questions in writing, on the page where the solution will go.

\begin{enumerate}[nosep,leftmargin=*]
\item \emph{What is the object?} An integer, a polynomial, a function, a sequence, a matrix, a set of points, a game position. Problems that look alike often have different objects, and the object is what the technique acts on. A problem about digits is about an integer; a problem about a sum of consecutive values is often about a polynomial; a problem about moves and boards is about a position.
\item \emph{What is asked?} The verbs are few. Show one exists. Show it holds for all. Find all. Compute a value. Decide yes or no. Each verb has a shape of answer. ``Find all'' needs a list and a proof that the list is complete, which is two arguments, not one. ``Decide'' needs a yes or no and a proof of whichever it is. ``Show one exists'' may be met by a construction or by a counting argument with no construction at all.
\item \emph{What is the smallest case?} $n$ equal to 1, 2 and 3. The zero polynomial and the constant polynomial. The empty set. A $2 \times 2$ matrix. Written out by hand, not imagined. The smallest cases name the technique more often than the statement does, because the statement is written to hide it and the small case is not. They are also the cheapest work available, and the case $n = 2$ is often the base case of the eventual proof.
\end{enumerate}

Ten minutes on these three questions is well spent. A technique chosen before them is a guess; after them it is a reading.

\subsection*{The call}

Before working, write one line in \texttt{calls.md} under the domain \texttt{math}:

\begin{quote}\small
\texttt{math: <problem>, <technique>, <first assumption>,}\\
\texttt{\phantom{math: }<time>. <confidence>}
\end{quote}

The technique is the one you expect to carry the problem, named in the vocabulary of this document's section titles. The first assumption is the first thing you are taking for granted: that the answer is the obvious bound, that the pattern in the small cases continues, that the function is a polynomial. The time is the box in minutes. The confidence is a number from 0 to 100 for how likely you think it is that the named technique carries it.

The line is the point. It is a prediction made before the evidence, and a prediction can be scored. Whether or not the problem gets solved, the call is settled.

\subsection*{Settling}

When the box runs out or the problem is done, open the solution in the archive and settle the call in the same sitting, with one word.

\begin{itemize}[nosep,leftmargin=*]
\item \emph{held}: the technique you named carried the problem, in your solution or in the published one.
\item \emph{missed}: a different technique from this document's list carried it. You knew the tool and did not pick it.
\item \emph{blind}: the solution uses something you did not consider at all. The tool was not in your hand.
\end{itemize}

A solved problem is settled as honestly as an unsolved one. If you solved it by a route you did not name, that is a miss, and recording it as a miss is what makes the readout trustworthy. When a later problem is carried by a technique you called on an earlier one, write the earlier call's id in the new call, as in ``as in c12''. That citation is how the readout sees a technique transferring from the problem where you learned it to the problem where you used it, and transfer is the thing the month is trying to build.

\subsection*{Time boxes}

Twenty-five minutes for an A1, A2, B1 or B2, or for a seminar problem. Forty-five for an A3 or B3, or for anything the ladder marks as harder. The box is not a measure of patience. It is the point after which the published solution teaches more than the struggle does, and the call line has already recorded what the struggle was for. Stopping at the box is part of the exercise, not a failure of it.

\subsection*{A worked reading}

The problem below is invented for this document.

\begin{example}
Let $n$ be a positive integer. Show that the number of ordered pairs $(a, b)$ of positive integers with $\lcm(a, b) = n$ is odd.
\end{example}

\emph{Object.} Pairs of positive integers tied to a fixed integer $n$ through its prime factorisation. The object is really $n$, and the pairs are a count attached to it.

\emph{Asked.} Show a count is odd for all $n$. ``Odd'' is a parity statement, and parity statements about counts are carried by two things: an explicit formula that is visibly odd, or a pairing of the counted objects that leaves exactly one unpaired.

\emph{Smallest cases.} For $n = 1$ the only pair is $(1, 1)$: one pair. For $n = 2$ the pairs are $(1, 2)$, $(2, 1)$ and $(2, 2)$: three. For $n = 4$: $(1, 4)$, $(4, 1)$, $(2, 4)$, $(4, 2)$, $(4, 4)$: five. For $n = 6$: $(1, 6)$, $(6, 1)$, $(2, 3)$, $(3, 2)$, $(2, 6)$, $(6, 2)$, $(3, 6)$, $(6, 3)$, $(6, 6)$: nine. The counts $1, 3, 5, 9$ are $2e + 1$ for a prime power $p^e$, and $9 = 3 \times 3$ for $6 = 2 \times 3$. The count looks multiplicative. Also, in every list, each pair $(a, b)$ with $a \ne b$ sits next to its reverse $(b, a)$, and only $(n, n)$ has no partner.

\emph{The call.} Two techniques present themselves, and the call names one:

\begin{quote}\small
\texttt{math: lcm pairs, multiplicative count over prime powers,}\\
\texttt{\phantom{math: }assume the count for a prime power is 2e+1, 15 min. 70}
\end{quote}

\emph{Working.} Write $n = \prod p_i^{e_i}$. A pair $(a, b)$ has $\lcm(a, b) = n$ exactly when, for each $i$, the exponents of $p_i$ in $a$ and $b$ are at most $e_i$ and at least one of them equals $e_i$. For one prime the number of exponent pairs $(\alpha, \beta)$ with $0 \le \alpha, \beta \le e$ and $\max(\alpha, \beta) = e$ is $(e+1)^2 - e^2 = 2e + 1$. The choices for different primes are independent, so the count is $\prod (2 e_i + 1)$, a product of odd numbers, hence odd.

\emph{Settling.} The published solution, had this been an archive problem, might instead pair $(a, b)$ with $(b, a)$: the map is an involution on the set of pairs, its only fixed point is $(n, n)$, so the set has odd size. That argument is two lines and uses no factorisation. If you had called the multiplicative count and it carried, the call is \emph{held}, even though a shorter route existed; the readout records which tool you reached for, not which was shortest. If the pairing had been in the small cases in front of you, as it was here, and you did not name it, the lesson goes in the margin, not in the call.

The reading took longer to write than to do. On the day it is five minutes, and it is the five minutes that decide whether the next forty are spent on the right thing.


\section{The four ladders}
\label{sec:ladders}

Six core problems a week, in order, each rung a little harder than the last on the same technique. The sixth rung is a problem filed under a different tag that the week's technique nevertheless carries; it is where the transfer citation (``as in c12'') is expected. Two optional rungs for weeks with room. The mix is three problems a week from the ladder's Putnam rungs and three from the seminar rungs; if time is short, the Putnam rungs come first. Do them in order. The ladder is the point: the easy rung teaches the move that the hard rung needs.

Each week also has a second rail, a technique that keeps appearing alongside the main one. It is named in the table heading and gets no separate ladder.


\subsection*{Week of 2026-11-09: congruences and number theory}

Second rail: pigeonhole.

\begin{center}\small
\begin{tabular}{@{}c >{\raggedright\arraybackslash}p{3.1cm} >{\raggedright\arraybackslash}p{4.8cm} >{\raggedright\arraybackslash}p{4.8cm}@{}}
\toprule
Rung & Problem & Object and ask & Move \\
\midrule
1 & \putnam{2018 A1} & An equation in positive integers; find all solutions. & Warm: factor and bound. \\
2 & \putnam{2017 B2} & Sums of consecutive integers; a smallest value. & Parity and divisibility of a sum. \\
3 & \seminar{PS 9 \#14} and \seminar{\#15} & Two classic non-divisibility statements. & The order of an element modulo a prime, and the smallest prime factor trick. \\
4 & \putnam{2020 A1} & Integers of a digit pattern; count those divisible by a given number. & Factor the pattern, then orders again. \\
5 & \putnam{2024 A1} & An exponential equation in integers; which exponents. & Reduce modulo small numbers until the cases die. \\
6 & \putnam{2021 A1} & Fixed-length hops in the plane; the minimum number to reach a point. & Tagged pigeonhole and extremal; the number theory of the hop lengths carries it. Cite the rung where you first used a modular obstruction. \\
\midrule
opt & \putnam{2023 B2} & The binary digits of multiples; a minimum. & \\
opt & \putnam{2018 B3} & Integers with linked divisibility conditions. & \\
\bottomrule
\end{tabular}
\end{center}

Seminar warm-ups: \seminar{PS 9 \#1, \#4, \#12, \#31}; \seminar{Pigeonhole \#1, \#6, \#7, \#8, \#9}.


\subsection*{Week of 2026-11-16: polynomials and their roots}

Second rail: recurrences and induction.

\begin{center}\small
\begin{tabular}{@{}c >{\raggedright\arraybackslash}p{3.1cm} >{\raggedright\arraybackslash}p{4.8cm} >{\raggedright\arraybackslash}p{4.8cm}@{}}
\toprule
Rung & Problem & Object and ask & Move \\
\midrule
1 & \putnam{2019 A1} & A symmetric expression in three integers; which values it takes. & Factor the symmetric form. \\
2 & \putnam{2023 A2} & An even-degree polynomial with prescribed values; the remaining solutions. & Build the polynomial from its known values. \\
3 & \seminar{Roots \#6} and \seminar{Roots \#7} & Two Putnam problems from 1968 and 1981 on real roots. & Rolle and sign changes. \\
4 & \putnam{2016 A1} & Derivatives of an integer polynomial; the smallest order that is always divisible by something. & Polynomials meet congruences, as in week 1. \\
5 & \putnam{2022 B1} & A power series built from an integer polynomial; its coefficients are nonzero. & Work modulo 2, which is the week 1 move inside a polynomial. \\
6 & \putnam{2017 A2} & A rational recurrence whose terms turn out to be integer polynomials. & Tagged recurrence; the polynomial identity carries it. Cite the rung where you first wrote a polynomial by its roots or values rather than its coefficients. \\
\midrule
opt & \putnam{2024 A2} & Iterates of a polynomial and divisibility; which polynomials. & \\
opt & \putnam{2019 B5} & A polynomial interpolating a famous recurrence; extrapolate. & \\
\bottomrule
\end{tabular}
\end{center}

Seminar: \seminar{Roots \#1, \#8, \#11a}; \seminar{PS 4 \#1, \#3, \#5, \#9}.


\subsection*{Week of 2026-11-23: sums, integrals and limits}

Second rail: inequalities. The first timed half paper is on Saturday 2026-11-28.

\begin{center}\small
\begin{tabular}{@{}c >{\raggedright\arraybackslash}p{3.1cm} >{\raggedright\arraybackslash}p{4.8cm} >{\raggedright\arraybackslash}p{4.8cm}@{}}
\toprule
Rung & Problem & Object and ask & Move \\
\midrule
1 & \seminar{Limits \#15} and \seminar{Limits \#12} & A Riemann sum and a product limit from older Putnams. & Recognise an integral in a sum, and logarithms to turn a product into a sum. \\
2 & \putnam{2016 B1} & A sequence defined through logarithms; sum a series from it. & Telescoping. \\
3 & \putnam{2022 A1} & A line and a logarithmic curve meeting exactly once; the parameters. & Convexity and tangency. \\
4 & \putnam{2019 B2} & A trigonometric sum; its asymptotic limit. & A sum that is a Riemann sum in disguise, with a correction term. \\
5 & \putnam{2021 A2} & A function defined by a limit; an asymptotic ratio. & Take logarithms, expand, keep the first surviving term. \\
6 & \putnam{2025 A2} & A sine bounded by a quadratic; the best constants. & Tagged inequality; the limit at zero and the equality case carry it. Cite the rung where you first expanded to the first surviving term. \\
\midrule
opt & \putnam{2020 A2} & A binomial sum; closed form. & \\
opt & \putnam{2024 B3} & Gaps between fixed points of a tangent-type function; a two-sided bound. & \\
opt & \putnam{2023 A1} & A product of cosines; its second derivative exceeds a threshold. & \\
\bottomrule
\end{tabular}
\end{center}

Seminar: \seminar{PS 2 \#1, \#5, \#8, \#15, \#16}; \seminar{PS 5 \#3, \#4, \#13}; \seminar{Limits \#1} to \seminar{\#4}, \seminar{\#11}, \seminar{\#13}, \seminar{\#14}, \seminar{\#16}.

\textbf{Mock 1, Saturday 2026-11-28.} The 2015 A paper as two sittings of 90 minutes with fifteen minutes between them, the shape of sessions A and B: A1, A3 and A5 in the first sitting, A2, A4 and A6 in the second. The split gives each sitting an approachable, a middle and a hard problem; that is a guess at how the new sessions are built, not something the policy says. Start at 11:00 and 12:45 if the day allows, the Eastern start times of sessions A and B. Room rules as in \secref{sec:examday}{exam day}: no phone or watch, one side of the paper, one problem per sheet, pens down at 90 minutes. Three call lines go on the scratch paper in the first eight minutes of each sitting. In the break, before looking at the scratch sheet again, write the three lines from memory, then compare; on the day the scratch paper is handed in, so memory is all that leaves the room. Type the six lines in that evening with the date. Settle every one, including those not attempted; \emph{moot} is the honest verb for those, and seeing how many there are is part of the point.


\subsection*{Week of 2026-11-30: counting and configurations}

Second rail: invariants, games and hidden symmetry. The second timed half paper is on Wednesday 2026-12-02. Thursday and Friday: no new problems; read the Misses and Known sections of \texttt{readout.md} once each, and the Weaknesses section of \texttt{profile.md}.

\begin{center}\small
\begin{tabular}{@{}c >{\raggedright\arraybackslash}p{3.1cm} >{\raggedright\arraybackslash}p{4.8cm} >{\raggedright\arraybackslash}p{4.8cm}@{}}
\toprule
Rung & Problem & Object and ask & Move \\
\midrule
1 & \seminar{PS 10 \#1, \#4} and \seminar{\#7} & Three exercise-like configuration problems, as the set itself says. & Colourings and sums. \\
2 & \putnam{2018 B1} & A set of lattice vectors; remove one so the rest split into two equal-sum halves. & Hidden symmetry: pair things up. \\
3 & \putnam{2019 B1} & Lattice points on circles; count squares. & Count by a bijection with something you can list. \\
4 & \putnam{2023 B1} & Coins sliding on a grid; which configurations are reachable. & An invariant of the moves. \\
5 & \putnam{2024 B1} & A grid with a prescribed set of values along a transversal; feasibility. & Invariants again, with a construction for the other direction. \\
6 & \putnam{2020 B1} & A signed sum over integers by the parity of their binary digits; a residue. & Tagged number theory; the hidden symmetry of the sign pattern carries it. Cite the rung where you first paired terms off. \\
\midrule
opt & \putnam{2017 A4} & Quiz scores; split a class into two halves with equal averages. & \\
opt & \putnam{2020 B2} & A peg-moving game; when the first player wins. & \\
\bottomrule
\end{tabular}
\end{center}

Seminar: \seminar{PS 10 \#3, \#5, \#16}; \seminar{PS 1 \#1, \#2, \#3}.

\textbf{Mock 2, Wednesday 2026-12-02.} The 2014 B paper cut the same way: B1, B3 and B5, a fifteen-minute break, then B2, B4 and B6. Same protocol.


\subsection*{Rhythm}

One problem a day minimum, two on two of the days, which makes six. The call before, the settle after, every time. Sunday evening: \texttt{/profile}. Any section of this document that misjudged you during the week, too easy or too hard or a tag that turned out wrong, gets a mark, and \texttt{/teach} runs before Monday so that the next week's reading matches where you are.


\section{Exam day, 2026-12-05}

Everything in the month points at two sessions of three hours. This section is the protocol for those six hours, and the checklist for the pages that get handed in.

\subsection*{Before the morning session}

Sleep. Eat something with protein. Bring several pens, a watch that is not a phone, and water. Do not read problems on the way in. The reading habit from \secref{sec:reading}{reading a problem} is the whole preparation; nothing learned in the last hour will be available under pressure.

\subsection*{The first fifteen minutes}

Read all six problems before writing anything. Then, on scratch paper, one call line for each, with a rank from 1 to 6:

\begin{quote}
\texttt{A1: <technique>, <first assumption>, rank 1. 70}
\end{quote}

No times, because the box on the day is set by the rank. This is the one habit from the month that pays directly. Six call lines take fifteen minutes and produce a plan for the remaining 165, and they stop the most common failure, which is spending the first hour on A3 because its statement was interesting.

The ranking is not by slot. A1 is usually the easiest, but not always, and a B2 that matches a ladder you climbed last week outranks a B1 that does not. Rank by confidence that you can finish, not by how much you like the problem.

\subsection*{The remaining 165 minutes}

The top-ranked problem until it is solved and written up. Then the second. Then the third if there is time, which there usually is not.

A problem that is not falling after 45 minutes gets dropped for the next on the rank, without regret. The score distribution in Section~1 is the argument: a finished second problem is worth 10 and a half-finished first is worth 1, so the move from a stuck problem to a fresh one is a gain of about nine points in expectation, even if the fresh one is harder. Write ``dropped at <time>'' on the scratch call line so the evening ledger is honest.

If the top-ranked problem falls quickly, write it up fully before moving on. The temptation is to leave the write-up for later and chase the second problem while the first is fresh; the result is often two 1s instead of a 10 and a 1.

Fifteen minutes before the end, stop solving. Check what is written against the checklist below. A 9 recovered to a 10 in those fifteen minutes is worth more than any progress on a new problem.

\subsection*{Writing up}

Graders give 9 and 10 to arguments they can follow without goodwill, and 0 to ideas. The write-up is a separate skill from the solving, and this is the checklist.

\begin{enumerate}[nosep,leftmargin=*]
\item State what you are proving at the top. One sentence. If the problem says ``find all'', state the answer there too.
\item Define every symbol you introduce, at its first use. ``Let $d$ be the smallest prime factor of $n$'' before $d$ appears in a formula.
\item Make every case explicit, including the one that is obvious. If the argument splits on $n$ even or odd, both branches are on the page, even if one is a line.
\item The equality case of an inequality, the base case of an induction, the $n = 1$ case of a sum: write them. These are the gaps that turn 10 into 9.
\item If the answer is a value or a set, box it, and say why no other value works. A correct value with no uniqueness argument is a 1.
\item Any fact used from outside the problem is named: Fermat's little theorem, Rolle, the pigeonhole principle. Named facts can be checked; unnamed ones look like gaps.
\item A construction is checked against every condition in the statement, in writing. ``This works because'' followed by the verification, not followed by nothing.
\item Nothing in the final write-up refers to the scratch work. Cross out what is not part of the argument. A grader who has to decide which of two attempts is the real one will pick neither.
\item Read the whole solution once as a stranger would, looking for a step that uses something not yet shown.
\end{enumerate}

Write in full sentences. A line of symbols with no verb is not an argument. ``Since $p$ does not divide $a$, Fermat gives $a^{p-1} \equiv 1 \pmod p$'' is a sentence; the same with the words removed is a hope.

\subsection*{Between sessions}

Eat. Do not discuss the morning problems with anyone, and do not look anything up. Whatever happened in the morning has no bearing on the afternoon, and the afternoon starts with the same fifteen minutes of reading.

\subsection*{That evening}

Type the twelve call lines into \texttt{calls.md}, with the date and the domain \texttt{putnam}, exactly as they were written on paper, including the ranks and the ones marked dropped. The unattempted ones are \emph{moot}. Settle them all when the archive posts the solutions, usually within the week. Run \texttt{/profile} the weekend after; the \texttt{putnam} row under Domains, and its lines under Misses, are the evidence the month was for.


\newpage
\section{Hidden independence and uniformity}

\begin{recognise}
The problem has a symmetry it does not mention. Points are chosen at random and the question is about their order or their arrangement. A sum runs over all permutations, all subsets, or all residues, and the summand is built from one coordinate. A quantity is asked for that, once computed, turns out not to depend on a parameter the problem made a fuss about. A sign pattern or a weight pairs terms off against each other. The answer to the smallest case is suspiciously clean: a probability of $1/2$ or $1/n$, an integer where a fraction was expected, a sum that is zero.
\end{recognise}

The model has three pieces. Uniformity: when a random object is drawn from a set that a bijection maps to itself, the image is drawn from the same distribution, so anything invariant under the bijection can be computed in the most convenient representative. Linearity: an expectation of a sum is a sum of expectations, whether or not the parts are independent, so a complicated count becomes a sum of probabilities of simple events. Independence proper: when a joint distribution factors, a probability of a conjunction is a product, and one hard event becomes several easy ones.

\begin{definition}
A random variable $X$ taking values in a finite set $S$ is uniform if $\Pr(X=s)=1/|S|$ for every $s\in S$. Events $A_1,\dots,A_n$ are independent if $\Pr(A_{i_1}\cap\dots\cap A_{i_k})=\Pr(A_{i_1})\cdots\Pr(A_{i_k})$ for every choice of distinct indices; pairwise independence is the case $k=2$ only.
\end{definition}

\begin{lemma}[Transport of uniformity]
Let $X$ be uniform on a finite set $S$ and let $\varphi\colon S\to S$ be a bijection. Then $\varphi(X)$ is uniform on $S$, and for every function $g$ on $S$, $\E\, g(X)=\E\, g(\varphi(X))$.
\end{lemma}

\begin{proof}
$\Pr(\varphi(X)=s)=\Pr(X=\varphi^{-1}(s))=1/|S|$. The second claim is the first one applied to the distribution of $g$.
\end{proof}

\begin{lemma}[Equal fibres]
Let a finite group $G$ act on a finite set $X$, let $A$ be a finite abelian group, and let $f\colon X\to A$ and a surjective homomorphism $\chi\colon G\to A$ satisfy $f(g\cdot x)=f(x)+\chi(g)$ for all $g$, $x$. Then every fibre $f^{-1}(a)$ has the same size $|X|/|A|$.
\end{lemma}

\begin{proof}
For $a,b\in A$ pick $g$ with $\chi(g)=b-a$. Then $x\mapsto g\cdot x$ maps $f^{-1}(a)$ into $f^{-1}(b)$, and $x\mapsto g^{-1}\cdot x$ maps it back, so the two fibres are in bijection.
\end{proof}

The lemma is the formal version of ``shift everything by one and the sum shifts by a unit''. It is why the sum of $n$ fair dice is uniform modulo $6$, why a residue class of a symmetric count is as large as any other, and why so many answers are $1/p$.

\begin{theorem}[Linearity of expectation]
For random variables $X_1,\dots,X_n$ with finite expectations, $\E(X_1+\dots+X_n)=\E X_1+\dots+\E X_n$. In particular, if $N$ counts how many of the events $A_1,\dots,A_n$ occur, then $\E N=\sum_i \Pr(A_i)$.
\end{theorem}

\begin{proof}
Expectation is a sum (or integral) weighted by the probability measure, and sums commute with sums. For the second statement write $N=\sum_i \mathbf 1_{A_i}$, where the indicator $\mathbf 1_{A}$ is $1$ on $A$ and $0$ elsewhere, and use $\E\,\mathbf 1_{A}=\Pr(A)$.
\end{proof}

\begin{proposition}[Exchangeability]
Let $X_1,\dots,X_n$ be independent with the same continuous distribution. Then every one of the $n!$ orderings of their values is equally likely, so $\Pr(X_k \text{ is the largest})=1/n$ for each $k$, and more generally any event defined by the relative order of the $X_i$ has probability equal to its share of the $n!$ orderings.
\end{proposition}

\begin{proof}
Ties have probability zero. Permuting the coordinates is a bijection of the sample space that preserves the product measure, and it maps the event ``the order is $\sigma$'' to ``the order is $\tau\sigma$''; so all orderings have the same probability, which must be $1/n!$.
\end{proof}

\begin{lemma}[Sign-reversing involution]
Let $S$ be a finite set, $f\colon S\to\R$, and $\sigma\colon S\to S$ an involution with $f(\sigma(s))=-f(s)$ whenever $\sigma(s)\neq s$. Then $\sum_{s\in S} f(s)=\sum_{s:\ \sigma(s)=s} f(s)$.
\end{lemma}

\begin{proof}
The non-fixed points split into two-element orbits $\{s,\sigma(s)\}$, each contributing $f(s)+f(\sigma(s))=0$.
\end{proof}

\begin{firstmoves}
Pair the terms: $k$ with $n-k$, a subset with its complement, a permutation with its inverse or with its composite with a fixed transposition, a point with its reflection. Ask whether the answer depends on the parameter at all; if the problem gives a specific large number, compute the case with $2$ or $3$ in its place and see whether the number was ever used. Write the quantity as a sum of indicators and compute one probability. Fix the thing that was random and randomise the thing that was fixed; the probability does not change, and often the second version is the easy one.
\end{firstmoves}

\begin{example}
A permutation $\sigma$ of $\{1,\dots,n\}$, $n\ge 2$, is chosen uniformly at random. Call a position $i$ with $1\le i\le n-1$ a step if $|\sigma(i+1)-\sigma(i)|=1$. Find the expected number of steps.
\end{example}

\begin{proof}[Solution]
Let $N$ be the number of steps and $A_i$ the event that position $i$ is a step. By linearity $\E N=\sum_{i=1}^{n-1}\Pr(A_i)$. For a fixed $i$ the ordered pair $(\sigma(i),\sigma(i+1))$ is uniform on the $n(n-1)$ ordered pairs of distinct values, because every permutation of the values maps the set of permutations bijectively to itself (transport of uniformity). Exactly $2(n-1)$ ordered pairs of distinct values differ by $1$: the pairs $(k,k+1)$ and $(k+1,k)$ for $1\le k\le n-1$. So $\Pr(A_i)=2(n-1)/(n(n-1))=2/n$ for every $i$, and
\[
\E N=(n-1)\cdot\frac{2}{n}=\frac{2(n-1)}{n}.
\]
Check at $n=2$: both permutations have one step, and the formula gives $1$. At $n=3$: the six permutations have step counts $2,1,1,1,1,2$, average $4/3$, as the formula says. The events $A_i$ are far from independent, and nothing above needed them to be.
\end{proof}

\begin{example}
For $n\ge 1$ let $\operatorname{fix}(\sigma)$ be the number of fixed points of a permutation $\sigma$ of $\{1,\dots,n\}$ and $\operatorname{sgn}(\sigma)$ its sign. Compute
\[
T_n=\sum_{\sigma\in S_n}\operatorname{sgn}(\sigma)\operatorname{fix}(\sigma).
\]
\end{example}

\begin{proof}[Solution]
Write $\operatorname{fix}(\sigma)=\sum_{i=1}^n \mathbf 1[\sigma(i)=i]$ and swap the sums:
\[
T_n=\sum_{i=1}^n\ \sum_{\sigma:\ \sigma(i)=i}\operatorname{sgn}(\sigma).
\]
Permutations fixing $i$ are permutations of the other $n-1$ elements with the same sign, so each inner sum is $\sum_{\tau\in S_{n-1}}\operatorname{sgn}(\tau)$. For $n-1\ge 2$ this is zero: composing with a fixed transposition $(a\,b)$ of two of the $n-1$ elements is a sign-reversing involution with no fixed points. For $n-1=1$ the inner sum is $1$, and for $n-1=0$ it is $1$ as well (the empty permutation has sign $1$). Hence $T_1=1$, $T_2=2$, and $T_n=0$ for all $n\ge 3$.

Check at $n=2$: the identity contributes $+2$, the transposition contributes $0$. At $n=3$: the identity gives $+3$, the three transpositions give $-1$ each, the two $3$-cycles give $0$; total $0$.

The lesson is in the small cases. The involution that kills the sum needs two elements to swap, and when there are not two the sum survives. A pairing argument always needs its fixed points listed.
\end{proof}

\begin{trap}
An independence that looks obvious and is false. Whether the first and second largest of a random sample land in particular positions are not independent events; whether two given positions of a random permutation are fixed points is not independent. The smallest case usually exposes it, so list all outcomes for $n=3$ before multiplying anything. The opposite trap: a uniformity that is true but was never used. The sum of $n$ fair dice is uniform modulo $6$ by the equal-fibres lemma, but modulo $7$ it is not (one die gives probability $0$ of a multiple of $7$), because $7$ is not the size of any coordinate's range. A pairing that leaves a fixed point unpaired, as in the second example, gives the wrong answer for exactly the cases where the fixed points exist. Finally, when a quantity does not depend on a parameter, say so and prove it; graders do not accept ``by symmetry'' without the bijection that realises it.
\end{trap}

\begin{ladder}
\item \seminar{PS 1 \#1}. Independent draws; the chance they come out in a given order.
\item \seminar{PS 1 \#2}. Ways of writing an integer as an ordered sum; a count by a bijection.
\item \seminar{PS 1 \#3}. Zero-one matrices with a parity condition on lines; how many.
\item \putnam{2018 B1}. A set of lattice vectors; remove one so the rest split into two equal-sum halves.
\item \putnam{2020 B1}. A signed sum over integers by the parity of their binary digits; a residue.
\item \seminar{PS 1 \#12}. Passengers taking seats with one rule broken; a probability about the last one.
\item \seminar{PS 1 \#20}. Particles on a segment that swap on collision; a time to clear.
\item \putnam{2019 A4}. A function with vanishing averages over all spheres of one size; must it be zero.
\item \putnam{2021 A4}. A double integral over a growing disc; the limit.
\item \putnam{2023 B3}. A random sequence and a longest subsequence of alternating type; its expectation.
\item \putnam{2017 A6}. Colourings of the edges of a polyhedron with a local condition; a count.
\end{ladder}

\begin{hides}
Inside number theory, where a residue count is the same in every class because translation moves the sum by a unit (\putnam{2020 B1}, \putnam{2017 B6}). Inside probability, where ``expected'' hides a sum of indicators (\putnam{2023 B3}, \putnam{2016 B4}). Inside analysis, where an integral over a symmetric region vanishes by an odd reflection (\putnam{2019 A4}, \putnam{2021 A4}). Inside counting, where a set is split into two equal halves by an involution (\putnam{2018 B1}). Inside recurrences, where a sequence returns to a value because the map that defines it is a bijection of a finite set (\putnam{2018 B4}).
\end{hides}

\section{Pigeonhole and extremal arguments}
\label{sec:pigeonhole}

Two arguments that prove existence without a construction. Pigeonhole: more objects than boxes forces two objects into one box. Extremal: the largest, smallest, first or last object has a property the others need not, and that property is enough.

\begin{recognise}
A problem asks to show that two things coincide, that some difference is divisible by something, that a pattern repeats, that a sum of some of the given numbers hits a target, or that an object with a property exists when no construction is in sight. The hypothesis has a count in it (``among any $n+1$ integers'', ``a set of $2n$ points'') and the conclusion has no count, only an existence. For extremal arguments: a finite set, a process that should stop, a configuration that should have a special element, or a statement to be proved for all $n$ where a smallest counterexample would be convenient.
\end{recognise}

\subsection*{The model}

\begin{theorem}[Pigeonhole principle]
If $n + 1$ or more objects are placed in $n$ boxes, some box holds at least two objects.
\end{theorem}

\begin{proof}
If every box held at most one object, the boxes would hold at most $n$ objects in total, fewer than were placed.
\end{proof}

\begin{theorem}[Strong form]
If $m$ objects are placed in $n$ boxes, some box holds at least $\lceil m/n \rceil$ objects, and some box holds at most $\lfloor m/n \rfloor$.
\end{theorem}

\begin{proof}
If every box held at most $\lceil m/n \rceil - 1 < m/n$ objects, the total would be less than $n \cdot m/n = m$. The other half is the same argument with the inequalities reversed.
\end{proof}

\begin{lemma}[Averaging]
If real numbers $x_1, \dots, x_n$ have mean $\mu$, then some $x_i \ge \mu$ and some $x_j \le \mu$.
\end{lemma}

\begin{proof}
If all $x_i < \mu$ then $\sum x_i < n\mu$, contradicting the definition of $\mu$. Likewise for the other side.
\end{proof}

The averaging lemma is the pigeonhole principle for real-valued quantities, and it is the form that appears in analysis and probability: an integral that equals zero means the integrand is zero somewhere or changes sign; an expectation equal to $\mu$ means some outcome is at least $\mu$.

\begin{theorem}[Extremal principle]
Every nonempty finite set of real numbers has a largest and a smallest element. Every nonempty set of positive integers has a smallest element.
\end{theorem}

The second sentence is the well-ordering of $\N$, and it is equivalent to induction. Its working form is the \emph{minimal counterexample}: to prove a statement for all positive integers, suppose it fails, take the smallest $n$ for which it fails, and derive a contradiction, usually by producing a smaller failure. The first sentence is used by looking at the extreme object directly: the longest edge, the point with the largest coordinate, the pair at minimal distance, the first move of a process, the last element of an ordering.

The whole skill in pigeonhole is choosing the boxes. The standard choices, each worth having in hand:

\begin{itemize}[nosep,leftmargin=*]
\item Residues modulo $n$: $n$ boxes for integers; two in a box have a difference divisible by $n$.
\item Partial sums $a_1, a_1 + a_2, \dots$ taken modulo $n$: two equal partial sums give a consecutive block with sum divisible by $n$; a zero partial sum gives one directly. Together, $n$ integers always have a nonempty consecutive block with sum divisible by $n$.
\item Intervals of length $1/n$ in $[0, 1)$: $n$ boxes for fractional parts; two in a box differ by less than $1/n$. This is how Dirichlet's approximation theorem is proved.
\item Subsets by size, or by the largest odd divisor: every positive integer is $2^k \cdot m$ with $m$ odd, and the numbers $1$ to $2n$ have only $n$ possible values of $m$.
\item Pairs that sum to a fixed value, pairs of opposite points, pairs of complementary subsets: $n$ boxes for $2n$ objects, when the conclusion wants two objects that ``match''.
\item Pieces of a region: a square cut into $n^2$ small squares, a disc cut into sectors, a segment cut into pieces; two points in a piece are close.
\end{itemize}

\begin{firstmoves}
Try the obvious boxes before anything cleverer: residues if the conclusion mentions divisibility, intervals if it mentions closeness, pairs if it mentions two things matching. Count the boxes and the objects explicitly and check that there really are more objects. For an existence statement with a bound, ask what the extreme object looks like: the largest element of the set, the configuration with the fewest crossings, the pair at smallest distance. Write down what being extreme forbids. For ``for all $n$'' statements that resist induction, take the smallest counterexample and ask what smaller object it contains.
\end{firstmoves}

\subsection*{Worked examples}

\begin{example}
A clock face has sixty marks, one at each minute. Thirty-one of the marks are painted. Show that two painted marks are diametrically opposite. Show also that thirty painted marks need not have this property.
\end{example}

\begin{proof}[Solution]
The sixty marks fall into thirty pairs of diametrically opposite marks: mark $k$ is paired with mark $k + 30$ for $k = 0, 1, \dots, 29$. These pairs are the boxes. Thirty-one painted marks placed in thirty boxes put two painted marks in the same box, and two marks in the same box are opposite.

For the second part, paint the marks $0, 1, \dots, 29$. Each box contains exactly one of these marks, so no two painted marks are opposite. Thirty-one is therefore the smallest number that forces an opposite pair.
\end{proof}

The two parts are the two halves of any good pigeonhole problem. The principle proves that the threshold is enough; a construction proves that it is sharp. A problem that says ``find the smallest $k$ such that'' always needs both, and the construction is usually a listing of one element from each box.

\begin{example}
A finite set $S$ of positive integers has the property that whenever $a$ and $b$ are in $S$ with $a < b$, the difference $b - a$ is also in $S$. Show that $S$ consists exactly of the multiples of its smallest element that do not exceed its largest element.
\end{example}

\begin{proof}[Solution]
Let $m$ be the smallest element of $S$, which exists because $S$ is finite and nonempty, and let $M$ be the largest.

First, every element of $S$ is a multiple of $m$. Take $b \in S$. If $b = m$ there is nothing to show. Otherwise $b > m$, so $b - m \in S$. Repeating, $b - m, b - 2m, \dots$ stay in $S$ as long as the previous term exceeds $m$. The sequence decreases, so it stops: some $b - km$ lies in $S$ and does not exceed $m$. Because $m$ is the smallest element of $S$, $b - km \ge m$, so $b - km = m$ and $b = (k+1) m$.

Second, every multiple of $m$ up to $M$ is in $S$. Write $M = Km$. Then $M - m = (K-1) m \in S$ when $K \ge 2$, and $(K-1)m - m = (K-2)m \in S$ when $K \ge 3$, and so on down to $m$. So $m, 2m, \dots, Km$ all lie in $S$, and by the first part nothing else does.
\end{proof}

The extremal element here is the minimum, and what being minimal forbids is any element strictly between $0$ and $m$. That single prohibition pins down the whole set. The argument is the Euclidean algorithm in disguise, which is the usual shape of an extremal argument about divisibility.

\begin{trap}
Counting the boxes wrong by one. ``Residues modulo $n$'' is $n$ boxes, not $n - 1$; ``the integers from $1$ to $2n$'' is $2n$ objects, not $2n - 1$; a sequence of $n$ partial sums together with the empty sum is $n + 1$ objects. Write the counts down.

An extremal object that need not exist. An infinite set of positive integers has a smallest element but need not have a largest; a set of real numbers in $(0, 1)$ has neither. Before taking the maximum, say why it exists.

A pigeonhole conclusion that is weaker than what was asked. Two objects in the same box share the box's property, nothing more. If the problem asks for three, the strong form is needed, and the count must be large enough for $\lceil m/n \rceil \ge 3$.

Finally, pigeonhole proves existence and says nothing about which objects. A problem that asks to find or construct them wants a different argument.
\end{trap}

\begin{ladder}
\item \seminar{Pigeonhole \#1}. A set of integers drawn from a range; show two of them have a prescribed sum. The pairs are the boxes.
\item \seminar{Pigeonhole \#6}. Integers in a range; show two have a sum or a difference divisible by a given even number. Residues, folded.
\item \seminar{Pigeonhole \#7}. Among more than half of the integers up to an even bound, two are coprime. Adjacent pairs.
\item \seminar{Pigeonhole \#8}. Among more than half of the integers up to an even bound, one divides another. Largest odd divisors as boxes.
\item \putnam{2021 A1}. Fixed-length hops in the plane; the minimum number needed to reach a given point. The bound is extremal; the obstruction to doing better is modular.
\item \putnam{2023 B2}. Binary digits of the multiples of an integer; a minimum over the multiples. The extremal multiple and what it must look like.
\item \putnam{2022 A2}. The square of a polynomial; how many coefficients can be negative. Extremal on the coefficient pattern.
\item \putnam{2025 B5}. Modular inverses and a descent; a lower bound on how often something happens. Counting residues against boxes.
\item \putnam{2025 A5}. Permutations with a sign pattern; which pattern is most frequent. Extremal over patterns.
\item \putnam{2024 A3}. Order statistics of a random tableau; whether a balanced pair must exist. Averaging over positions.
\item \putnam{2016 B3}. A finite point set in the plane; a bounded triangle that covers it. The extremal triangle.
\end{ladder}

The seminar set \seminar{Pigeonhole} is the warm-up; its own notes say \seminar{\#9} is much harder and not really pigeonhole, and \seminar{\#17(a)} is Fermat's classic. Do \seminar{\#1} to \seminar{\#8} before the Putnam rungs.

\begin{hides}
Number theory, as a bound on a count of residues: a set of residues too large to avoid a coincidence is pigeonhole, whatever the problem is tagged. Analysis, as the averaging lemma: a continuous function whose integral is zero vanishes somewhere. Combinatorics, as the strong form: a two-colouring of enough objects has a large monochromatic class. Geometry, as pieces of a region: enough points in a square force two close together. Any ``find the smallest $k$'' problem, whose upper half is pigeonhole and lower half is a construction.
\end{hides}


\section{Invariants and monovariants}
\label{sec:invariants}

An invariant is a quantity a process leaves unchanged. A monovariant is a quantity the process moves in one direction only. The first proves that a state is unreachable; the second proves that a process stops.

\begin{recognise}
The problem describes moves: a token slides, two numbers are replaced by one, a square is recoloured, a stone is removed, a player acts. It then asks one of three things. Whether a given state can be reached from the start. Whether the process must stop, and if so after how many steps or in what state. What the possible final states are. Any of the three with a board, a multiset of numbers, or a configuration of objects is an invariant problem until shown otherwise.
\end{recognise}

\subsection*{The model}

Let a process have a set of states and a set of allowed moves, each taking a state to another state.

\begin{definition}
A function $I$ on the states is an \emph{invariant} if $I(s') = I(s)$ whenever a move takes $s$ to $s'$. It is a \emph{monovariant} if $I(s') < I(s)$ for every move (or $>$ for every move), with values in a set where this matters.
\end{definition}

\begin{theorem}[Invariant principle]
If $I$ is an invariant and $I(\text{start}) \ne I(\text{target})$, then no sequence of moves reaches the target from the start.
\end{theorem}

\begin{proof}
By induction on the number of moves, $I$ at every reachable state equals $I(\text{start})$. The target has a different value, so it is not reachable.
\end{proof}

\begin{theorem}[Monovariant principle]
If $I$ takes values in the non-negative integers and strictly decreases at every move, then every sequence of moves is finite, and has at most $I(\text{start})$ moves.
\end{theorem}

\begin{proof}
After $k$ moves the value is at most $I(\text{start}) - k$ and at least $0$, so $k \le I(\text{start})$.
\end{proof}

The same conclusion holds for any strictly decreasing quantity in a well-ordered set, for example a tuple of non-negative integers compared lexicographically, or a bounded set of multiples of a fixed positive real. It fails for real numbers in general: a quantity can strictly decrease forever and stay positive. When the monovariant is real, one must also show that each step decreases it by at least a fixed amount, or argue separately.

\begin{proposition}[Both directions]
To determine exactly which states are reachable, two arguments are needed: an invariant that rules out every unreachable state, and a construction that reaches every remaining state.
\end{proposition}

This is not a theorem so much as a reminder of the shape of a complete answer. An invariant on its own proves only the negative half.

The invariants that recur, in rough order of how often they are the answer:

\begin{itemize}[nosep,leftmargin=*]
\item Parity of a sum, of a count, or of the number of objects on squares of one colour.
\item A sum or a product modulo a small number: modulo 3 when a move changes two things by $1$ in opposite directions and a third by $2$, modulo $n$ when positions live on a cycle of length $n$.
\item A weighted sum, $\sum c_i x_i$ with weights chosen so that every move's net change is $0$. Finding the weights is a small linear system, and this is the systematic way to search for an invariant when parity fails.
\item The greatest common divisor of a multiset, under moves that add or subtract elements from each other.
\item A colouring count: colour the board so that each move changes the count on each colour in a controlled way.
\item For monovariants: the sum of the numbers on the board, the sum of squares, the number of inversions of a sequence, the number of adjacent pairs in some relation, the largest element.
\end{itemize}

\begin{firstmoves}
Write the move as a change in a vector of counts or positions, and compute what it does to the sum, the sum modulo 2, the sum modulo 3, and the sum of squares. Colour the board like a chessboard and count objects on black squares. If a move changes two quantities, look at their difference and their sum. If a move replaces $a$ and $b$ by something, look for a function $f$ with $f(\text{new}) = f(a) + f(b)$ or $f(a) f(b)$. For termination, write down the quantity the process is obviously eroding and check that every move erodes it, not just the typical move. Then do the smallest case by hand and list the states actually reached, which shows what the invariant has to separate.
\end{firstmoves}

\subsection*{Worked examples}

\begin{example}
Twelve lamps stand in a row at positions $1$ to $12$, all off. A move chooses two lamps whose positions differ by an even number and toggles both. Which configurations can be reached?
\end{example}

\begin{proof}[Solution]
Call a lamp odd or even by the parity of its position. Two positions differ by an even number exactly when they have the same parity, so every move toggles two odd lamps or two even lamps.

\emph{Invariant.} Let $E$ be the number of even lamps that are on and $O$ the number of odd lamps that are on. A move toggling two even lamps changes $E$ by $-2$, $0$ or $+2$ and leaves $O$ alone; likewise for odd. So the parities of $E$ and $O$ are each invariant. At the start both are $0$. Hence in every reachable configuration, $E$ and $O$ are both even.

\emph{Construction.} Conversely, let a configuration have $E$ and $O$ both even. Toggling positions $i$ and $i+2$ is a legal move, and a chain of such moves toggles any two lamps of the same parity: to toggle $i$ and $j$ with $i < j$ of the same parity, toggle $(i, i+2)$, then $(i+2, i+4)$, and so on to $(j-2, j)$; every intermediate lamp is toggled twice and ends where it began. The even lamps that are on in the target number $E$, an even count, so they can be paired off and each pair switched on by one chain; likewise the odd ones. So the configuration is reachable.

The reachable configurations are exactly those with an even number of lit lamps among the even positions and an even number among the odd positions. There are $2^5 \cdot 2^5 = 1024$ of them out of $4096$.
\end{proof}

The move allowed differences $4$, $6$ and so on, and the construction used only difference $2$. That is common: a move set larger than needed for the construction is still constrained by the same invariant, and the invariant is the half of the answer that uses every move.

\begin{example}
On a board are finitely many positive integers, at least two. A move chooses two of them, $a > b$, and replaces $a$ by $a - b$. Show that the process must stop, and determine what is on the board when it does.
\end{example}

\begin{proof}[Solution]
\emph{Termination.} The sum of the numbers on the board is a positive integer, and each move decreases it by $b \ge 1$. By the monovariant principle, there are at most as many moves as the initial sum. The process stops exactly when no two numbers differ, that is, when all numbers on the board are equal.

\emph{The final state.} Let $g$ be the greatest common divisor of the numbers on the board. A move replaces $a$ by $a - b$. Any common divisor of the old numbers divides $a - b$, and any common divisor of the new numbers divides $(a - b) + b = a$; so the set of common divisors is unchanged and $g$ is invariant. At the end all numbers equal some $d$, and the greatest common divisor of a board of equal numbers $d$ is $d$ itself. So $d = g$: the board ends with every entry equal to the greatest common divisor of the starting numbers, and the number of entries never changes.
\end{proof}

A monovariant for the stopping and an invariant for the final state, in that order. Most termination problems that also ask ``in what state'' have this two-part shape, and it is worth naming both quantities before proving anything about either.

\begin{trap}
A quantity preserved by most moves and not by one. The lamp problem above would break if a single move of odd difference were allowed; check every move in the list, especially the ones at the boundary of the board where a move may behave differently.

An invariant proved and the construction forgotten. ``The parity shows configuration $X$ is unreachable'' answers a yes or no question; a ``which configurations'' question needs the construction for everything the invariant does not exclude, and the construction is where the writing-up points are lost.

A monovariant in the reals with no lower bound on the step. Strictly decreasing and bounded below does not mean finite; the quantity might converge without arriving.

A monovariant that is only weakly monotone. If some move leaves the quantity unchanged, the process may cycle among states of equal value. Either strengthen the quantity, by adding a tie-breaking second coordinate compared lexicographically, or show directly that the equal-value moves cannot repeat.
\end{trap}

\begin{ladder}
\item \putnam{2023 B1}. Coins sliding diagonally on a grid; which configurations can be reached. A chessboard colouring is the invariant; a construction finishes.
\item \putnam{2024 B1}. A grid with prescribed values along a transversal; whether such a grid can be filled. An invariant of the arrangement, then a construction for the feasible cases.
\item \putnam{2017 A4}. Quiz scores in a class; split the class into two halves with equal averages. A monovariant on an adjustment process, or an invariant of a swap.
\item \putnam{2022 B3}. Recolouring the positive reals by a rule; whether everything can end up one colour. An invariant on the set of colours, with the reals as the surprise.
\item \putnam{2020 B2}. A peg-moving game on a line; when the first player wins. The invariant is the winning position's signature, see Section~\ref{sec:games}.
\item \putnam{2025 A3}. A game on digit strings with no repeats; who wins. Parity of the remaining moves.
\end{ladder}

Seminar: \seminar{PS 10 \#13}, a stone-removal game with a divisibility rule; \seminar{PS 10 \#14}, a signing game on the faces of a polyhedron; \seminar{PS 10 \#27}, a redistribution between boxes and which states it reaches. All three are invariant problems wearing other clothes.

\begin{hides}
Games: the set of losing positions is the level set of an invariant that the winning player restores at each turn. Number theory: the Euclidean algorithm is a monovariant with an invariant, as in the second example. Combinatorics on sequences: sorting arguments count inversions, which is a monovariant. Linear algebra: a weighted sum invariant is a vector in the kernel of the move matrix, and finding it is solving a linear system. Any problem that says ``operation'' or ``step'' and asks about the end.
\end{hides}


\section{Induction}
\label{sec:induction}

Induction proves a statement about every positive integer by proving it for the first one and showing that each case follows from the ones before it. Everyone knows the shape. The Putnam skill is in two choices: which smaller cases to use, and what statement to prove, since the statement given is often not the one that the step can be run on.

\begin{recognise}
A natural number $n$ in the statement, and an object for $n$ that is built from objects for smaller values: a sequence from its earlier terms, a board from a smaller board, a set from a set with one fewer element, a polynomial of degree $n$ from one of degree $n-1$. ``Show that for all $n$'' is the usual signal; so is ``show that the process always'' with a size that shrinks. A closed form guessed from the first terms of a recurrence is proved by induction; an inequality with a product or a sum of $n$ terms is often one too.
\end{recognise}

\subsection*{The model}

\begin{theorem}[Induction]
Let $P(n)$ be a statement about the integer $n$. If $P(1)$ holds, and $P(n-1)$ implies $P(n)$ for every $n \ge 2$, then $P(n)$ holds for every $n \ge 1$.
\end{theorem}

\begin{proof}
Suppose not, and let $n$ be the smallest integer for which $P(n)$ fails; it exists because the positive integers are well ordered. Then $n \ge 2$ since $P(1)$ holds, so $P(n-1)$ holds by minimality, so $P(n)$ holds by the step. Contradiction.
\end{proof}

\begin{theorem}[Strong induction]
If $P(1)$ holds, and for every $n \ge 2$ the statements $P(1), \dots, P(n-1)$ together imply $P(n)$, then $P(n)$ holds for every $n \ge 1$.
\end{theorem}

\begin{proof}
Apply ordinary induction to the statement $Q(n)$: ``$P(k)$ holds for all $k \le n$''. $Q(1)$ is $P(1)$, and $Q(n-1)$ implies $Q(n)$ by the hypothesis.
\end{proof}

The proof shows that strong induction is ordinary induction applied to a stronger statement. That is the general pattern: whenever the step from $n-1$ to $n$ cannot be made, the remedy is to prove more.

\begin{proposition}[Strengthening]
Suppose $P(n)$ is to be proved for all $n$, and let $P'(n)$ be a statement that implies $P(n)$. If $P'(1)$ holds and $P'(n-1)$ implies $P'(n)$, then $P(n)$ holds for all $n$.
\end{proposition}

This is immediate, and it is the most useful sentence in the section. A bound of the form ``$S_n < C$'' almost never passes from $n-1$ to $n$, because the hypothesis $S_{n-1} < C$ gives no room: adding a positive term to something below $C$ need not stay below $C$. The strengthened statement ``$S_n \le C - f(n)$'' with $f$ decreasing to $0$ does pass, provided $f(n-1) - f(n)$ is at least the $n$-th term. Choosing $f$ is the work, and the first example below does it.

Two further shapes:

\begin{itemize}[nosep,leftmargin=*]
\item \emph{Induction on a different variable.} A statement about pairs $(m, n)$ can be proved by induction on $m$ with $n$ fixed, on $n$ with $m$ fixed, on $m + n$, or on $\max(m, n)$. The right choice is the one for which the step has something to use.
\item \emph{Induction with several base cases.} If the step uses $P(n-2)$, or $P(\lfloor n/2 \rfloor)$, the base cases must cover every $n$ that the step cannot reach; the step from $P(n-2)$ needs both $P(1)$ and $P(2)$.
\end{itemize}

\begin{firstmoves}
Compute the cases $n = 1, 2, 3, 4$ by hand. Try the step on the statement as given. When it fails, write down exactly what fact about $n - 1$ you wished you had, and add that fact to the claim. For a bound, guess a correction term $c/n$ or $c/n^2$ and solve for the constant $c$ that makes the step work; then check the base case against it. For a structural statement, ask what object for $n - 1$ sits inside the object for $n$, and whether there might be several (then strong induction). For a recurrence, guess the closed form from the first six terms and let the induction verify it.
\end{firstmoves}

\subsection*{Worked examples}

\begin{example}
Show that for every positive integer $n$,
\[
\prod_{k=1}^{n} \left(1 + \frac{1}{k^3}\right) < 3 .
\]
\end{example}

\begin{proof}[Solution]
Let $P_n$ be the product. The values begin $P_1 = 2$, $P_2 = 9/4$, $P_3 = 7/3$, $P_4 \approx 2.37$, so the statement is plausible. The direct step fails: from $P_{n-1} < 3$ one gets only $P_n < 3(1 + 1/n^3)$, which is more than $3$.

Strengthen the claim to
\[
P_n \le 3 - \frac{1}{n} \qquad (n \ge 1).
\]
The base case is $P_1 = 2 \le 3 - 1$, with equality. For the step, let $n \ge 2$ and assume $P_{n-1} \le 3 - 1/(n-1)$. Then
\[
P_n = P_{n-1}\left(1 + \frac{1}{n^3}\right) \le \left(3 - \frac{1}{n-1}\right)\left(1 + \frac{1}{n^3}\right)
= 3 - \frac{1}{n-1} + \frac{3}{n^3} - \frac{1}{(n-1)n^3}.
\]
We want this to be at most $3 - 1/n$. Since $\frac{1}{n-1} - \frac{1}{n} = \frac{1}{n(n-1)}$, the requirement is
\[
\frac{3}{n^3} - \frac{1}{(n-1)n^3} \le \frac{1}{n(n-1)} .
\]
Multiplying by the positive number $(n-1) n^3$ turns this into $3(n-1) - 1 \le n^2$, that is, $n^2 - 3n + 4 \ge 0$, which holds for every real $n$ because $n^2 - 3n + 4 = (n - \tfrac32)^2 + \tfrac74 > 0$. So $P_n \le 3 - 1/n < 3$.
\end{proof}

The constant $1$ in the correction $1/n$ was found by running the step with an unknown $c$ and asking what $c$ makes the resulting inequality true for all $n \ge 2$; the inequality becomes $3(n-1) \le c(n^2 + 1)$, and $c = 1$ is the smallest integer that works. Then the base case was checked, and it happened to hold with equality, which is a sign the strengthening is tight. The actual limit of $P_n$ is about $2.43$.

\begin{example}
Let $t_k = 1 + \dfrac{k(k-1)}{2}$ for $k \ge 1$, so the sequence begins $1, 2, 4, 7, 11, 16, 22$. Show that every positive integer is a sum of distinct terms of this sequence.
\end{example}

\begin{proof}[Solution]
Strong induction on $n$. For $n = 1$, take $t_1$ alone.

Let $n \ge 2$ and suppose every positive integer less than $n$ is a sum of distinct terms. Let $k$ be the largest index with $t_k \le n$; it exists because the sequence is increasing and unbounded, and $k \ge 2$ since $t_2 = 2 \le n$. By the choice of $k$, $n < t_{k+1} = t_k + k$. Put $m = n - t_k$, so $0 \le m < k$.

If $m = 0$ then $n = t_k$ and we are done. Otherwise $1 \le m < k$, and we need the small fact
\[
t_k - k = 1 + \frac{k(k-1)}{2} - k = \frac{(k-1)(k-2)}{2} \ge 0 ,
\]
so $k \le t_k$ and hence $m < t_k$. By the induction hypothesis, $m$ is a sum of distinct terms of the sequence. Each of those terms is at most $m$, hence less than $t_k$, so none of them is $t_k$. Adding $t_k$ gives a representation of $n$ as a sum of distinct terms.
\end{proof}

The induction hypothesis was used at $m$, which is not $n - 1$; ordinary induction would have been useless here. The fact $k \le t_k$ is the kind of lemma that the step demands and the statement never mentions, and finding it is the whole problem. The same argument shows that any increasing sequence starting at $1$ in which each term is at most one more than the sum of all earlier terms has this property; the sequence above satisfies the stronger condition $t_{k+1} - t_k \le t_k$, which is what was actually used.

\begin{trap}
A base case that is not the one the step starts from. If the step is valid only for $n \ge 3$, then $P(1)$ and $P(2)$ both need checking by hand. If the step uses $P(n-2)$, two base cases. If the step uses $P(\lfloor n/2 \rfloor)$, the base case is $n = 1$ and the step must handle $n = 2$ and $n = 3$ correctly, since both reduce to $n = 1$.

A step that silently uses $n \ge 2$, or divides by $n - 1$, or picks ``an element other than the first'' from a set that might have one element. Read the step at $n = 2$ specifically.

A strengthening that is too strong. The claim $S_n \le C - 2/n$ might make the step easier and the base case false. Always check the base case against the strengthened claim, not the original one.

Induction that proves the closed form but not the recurrence. If a problem defines a sequence by a recurrence and asks for a property, proving the property for a guessed closed form only helps once the closed form is itself proved by induction.
\end{trap}

\begin{ladder}
\item \putnam{2017 A1}. A set of positive integers closed under two rules; find the elements not in it. Induction on the integer, with the small cases as base.
\item \seminar{PS 4 \#1}. A third-order recurrence; express it through known sequences. Guess and verify.
\item \seminar{PS 4 \#3}. The parity of the integer part of a power; a pattern in $n$. Induction carried by a conjugate.
\item \seminar{PS 4 \#5}. A ceiling of a power of a surd; divisibility by a power of two. Strengthen to a statement about the conjugate pair.
\item \putnam{2016 A5}. A finite group; every element is a short word in two generators. Induction on the size of the set reached.
\end{ladder}

\begin{hides}
Recurrences: every closed form guessed from the first terms is proved by induction, and the strengthening is usually a second sequence (a conjugate, a companion) carried alongside. Inequalities: products of the form $\prod (1 + a_k)$ and Bernoulli-type bounds, as in the first example. Combinatorics: a counting formula verified by removing one element; a tiling of a board built from tilings of smaller boards. Games: the characterisation of losing positions in Section~\ref{sec:games} is an induction on the distance to the end. Analysis: a sequence shown monotone and bounded by proving both at once, which is a strengthening.
\end{hides}


\section{Inequalities}

\begin{recognise}
A sum, product or integral has to be bounded. The largest or smallest value of an expression is asked for, usually under a constraint such as a fixed sum or a fixed product. One function has to be shown to stay on one side of another on an interval. A sequence has to be shown increasing, or bounded, or both. The words ``best constant'' or ``for all positive reals'' appear. Behind almost all of these sit three results and one habit: the arithmetic-geometric mean inequality, Cauchy-Schwarz, convexity, and the habit of finding the equality case first.
\end{recognise}

\begin{definition}
A function $f$ on an interval $I$ is convex if $f(tx+(1-t)y)\le tf(x)+(1-t)f(y)$ for all $x,y\in I$ and $t\in[0,1]$, and strictly convex if the inequality is strict whenever $x\ne y$ and $0<t<1$. It is concave if $-f$ is convex. A twice differentiable $f$ is convex exactly when $f''\ge 0$ on $I$, and strictly convex when $f''>0$ except at isolated points.
\end{definition}

\begin{lemma}[Tangent line]
Let $f$ be differentiable and convex on an interval $I$ and let $a\in I$. Then $f(x)\ge f(a)+f'(a)(x-a)$ for all $x\in I$. If $f$ is strictly convex, equality holds only at $x=a$.
\end{lemma}

\begin{proof}
Convexity makes $f'$ non-decreasing. For $x>a$ the mean value theorem gives $f(x)-f(a)=f'(\xi)(x-a)$ with $a<\xi<x$, and $f'(\xi)\ge f'(a)$. For $x<a$ the same with $f'(\xi)\le f'(a)$ and $x-a<0$. Strict convexity makes $f'$ strictly increasing and the inequalities strict.
\end{proof}

\begin{theorem}[Jensen]
Let $f$ be convex on $I$, let $x_1,\dots,x_n\in I$, and let $w_1,\dots,w_n\ge 0$ with $\sum w_i=1$. Then
\[
f\Bigl(\sum_i w_ix_i\Bigr)\le\sum_i w_if(x_i).
\]
For strictly convex $f$ and all $w_i>0$, equality holds only when all $x_i$ are equal. For concave $f$ the inequality reverses.
\end{theorem}

\begin{proof}
Let $m=\sum w_ix_i$, which lies in $I$. Convex functions have a supporting line at every interior point: a number $c$ with $f(x)\ge f(m)+c(x-m)$ for all $x\in I$ (for differentiable $f$, $c=f'(m)$ by the tangent line lemma; in general any $c$ between the one-sided derivatives works). Multiply by $w_i$ at $x=x_i$ and add: $\sum w_if(x_i)\ge f(m)+c\sum w_i(x_i-m)=f(m)$. In the strict case each inequality $f(x_i)\ge f(m)+c(x_i-m)$ is strict unless $x_i=m$.
\end{proof}

\begin{theorem}[Arithmetic-geometric mean]
For positive reals $a_1,\dots,a_n$,
\[
\frac{a_1+\dots+a_n}{n}\ \ge\ (a_1\cdots a_n)^{1/n},
\]
with equality if and only if all $a_i$ are equal. More generally, for weights $w_i>0$ with $\sum w_i=1$, $\sum w_ia_i\ge\prod a_i^{w_i}$, with the same equality case.
\end{theorem}

\begin{proof}
$\ln$ is strictly concave on $(0,\infty)$ since its second derivative is $-1/x^2<0$. Jensen gives $\ln(\sum w_ia_i)\ge\sum w_i\ln a_i=\ln\prod a_i^{w_i}$, with equality only when all $a_i$ coincide. Exponentiate.
\end{proof}

\begin{theorem}[Cauchy-Schwarz]
For real $a_1,\dots,a_n$ and $b_1,\dots,b_n$,
\[
\Bigl(\sum_i a_ib_i\Bigr)^2\le\Bigl(\sum_i a_i^2\Bigr)\Bigl(\sum_i b_i^2\Bigr),
\]
with equality if and only if one sequence is a scalar multiple of the other. The same holds for integrals: $(\int fg)^2\le\int f^2\int g^2$ for continuous $f,g$ on an interval, with equality if and only if $f$ and $g$ are proportional.
\end{theorem}

\begin{proof}
If all $b_i=0$ both sides vanish. Otherwise $q(t)=\sum_i(a_i-tb_i)^2=\sum a_i^2-2t\sum a_ib_i+t^2\sum b_i^2$ is a quadratic in $t$ that is never negative, so its discriminant $4(\sum a_ib_i)^2-4\sum a_i^2\sum b_i^2$ is at most $0$. Equality means $q$ has a real root $t_0$, so $a_i=t_0b_i$ for all $i$. The integral case is the same computation with $\int(f-tg)^2\ge 0$, and a non-negative continuous function with zero integral is identically zero.
\end{proof}

Two consequences are used constantly. With $b_i=1$: $(\sum a_i)^2\le n\sum a_i^2$, the mean of squares dominates the square of the mean. With $a_i=\sqrt{x_i}$ and $b_i=\sqrt{y_i}$ rearranged: for positive $y_i$, $\sum x_i^2/y_i\ge(\sum x_i)^2/\sum y_i$, which is how a sum of fractions is bounded below in one line.

\begin{theorem}[Bernoulli]
For real $x\ge -1$: $(1+x)^r\ge 1+rx$ when $r\ge 1$ or $r\le 0$, and $(1+x)^r\le 1+rx$ when $0\le r\le 1$. Equality holds only at $x=0$ unless $r\in\{0,1\}$.
\end{theorem}

\begin{proof}
$g(x)=(1+x)^r$ has $g''(x)=r(r-1)(1+x)^{r-2}$, which is non-negative on $(-1,\infty)$ exactly when $r\ge 1$ or $r\le 0$. The tangent line at $0$ is $1+rx$, and the tangent line lemma gives the result; for $0<r<1$ the function is concave and the inequality reverses. For integer $r\ge 1$ a direct induction also works: $(1+x)^{k+1}\ge(1+kx)(1+x)=1+(k+1)x+kx^2\ge 1+(k+1)x$, where the first step uses $1+x\ge 0$.
\end{proof}

\begin{lemma}[Comparison of derivatives]
Let $f,g$ be differentiable on $[a,b]$ with $f(a)=g(a)$ and $f'(x)\ge g'(x)$ for $a<x<b$. Then $f\ge g$ on $[a,b]$; if the derivative inequality is strict then $f>g$ on $(a,b]$.
\end{lemma}

\begin{proof}
$h=f-g$ has $h(a)=0$ and $h'\ge 0$, so $h$ is non-decreasing by the mean value theorem.
\end{proof}

The lemma is used repeatedly: to prove $f>g$ near a shared point, differentiate both, check the derivatives at the point, and if they also agree, differentiate again. The first derivative where they differ decides, which is the same first-surviving-term idea as in the limits section.

\begin{firstmoves}
Find where equality should hold and check the claim there; if the claim is false at the symmetric point, the symmetric point is not the extreme and the extreme is on the boundary. For $n$ variables under a symmetric constraint, try all equal and try one variable carrying everything. Name the inequality by the equality case: all equal points to AM-GM or Jensen; proportional sequences point to Cauchy-Schwarz. For a function on an interval, evaluate at the endpoints and at the stationary points. For a function that is not convex on the whole domain, try the tangent line at the equality point anyway and check the one-variable inequality it produces by factoring. For a product, take logarithms. For a sum of $n$ terms with a parameter $n$, compute $n=1$ and $n=2$ exactly.
\end{firstmoves}

\begin{example}
Let $a,b,c>0$ with $a+b+c=3$. Show that
\[
\frac{a^2}{a^2-a+1}+\frac{b^2}{b^2-b+1}+\frac{c^2}{c^2-c+1}\le 3,
\]
and determine when equality holds.
\end{example}

\begin{proof}[Solution]
Let $f(x)=x^2/(x^2-x+1)$ for $x>0$; the denominator is $(x-\tfrac12)^2+\tfrac34>0$. The natural equality case is $a=b=c=1$, where $f(1)=1$ and the sum is $3$. Jensen is not available: $f(0)=0$, $f(\tfrac12)=\tfrac13$, $f(1)=1$, $f(2)=\tfrac43$, $f(3)=\tfrac97$, so $f$ is convex near $0$ and concave past $1$, and neither direction of Jensen applies on $(0,3)$.

Use the tangent line at $x=1$ instead. $f'(x)=\dfrac{2x(x^2-x+1)-x^2(2x-1)}{(x^2-x+1)^2}=\dfrac{x(2-x)}{(x^2-x+1)^2}$, so $f'(1)=1$ and the tangent line is $y=x$. The claim $f(x)\le x$ for $x>0$ reads $x^2\le x(x^2-x+1)$, that is $0\le x^2-2x+1=(x-1)^2$, true for all $x$, with equality only at $x=1$. Summing, $f(a)+f(b)+f(c)\le a+b+c=3$, with equality exactly when $a=b=c=1$.

The move to notice: the tangent line was guessed from the equality case, and the one-variable inequality it produced was checked by factoring, not by convexity. The factor $(x-1)^2$ is the signature of a correct tangent at $x=1$; if it had not appeared, the tangent line would have been the wrong tool.
\end{proof}

\begin{example}
Let $f\colon[0,1]\to\R$ be continuous with $\int_0^1 f(x)\,dx=1$ and $\int_0^1 xf(x)\,dx=1$. Show that $\int_0^1 f(x)^2\,dx\ge 4$, and find all $f$ for which equality holds.
\end{example}

\begin{proof}[Solution]
First find the equality case. Equality in Cauchy-Schwarz needs $f$ proportional to the function we pair it with, and the only functions whose integrals against $f$ we know are combinations of $1$ and $x$. So try $f(x)=\alpha+\beta x$ and impose the constraints: $\alpha+\beta/2=1$ and $\alpha/2+\beta/3=1$, which give $\beta=6$ and $\alpha=-2$. The candidate is $g(x)=6x-2$, with $\int_0^1 g^2=\int_0^1(36x^2-24x+4)\,dx=12-12+4=4$.

Now the proof. For any $f$ satisfying the constraints, $\int_0^1 g f=6\int xf-2\int f=6-2=4$. By Cauchy-Schwarz,
\[
16=\Bigl(\int_0^1 gf\Bigr)^2\le\int_0^1 g^2\int_0^1 f^2=4\int_0^1 f^2,
\]
so $\int_0^1 f^2\ge 4$. Equality forces $f=\lambda g$ for a constant $\lambda$, and $\int f=1$ then forces $\lambda=1$. So equality holds only for $f(x)=6x-2$.

The order of work matters: the equality case was computed first and it told us which function to pair $f$ with. Pairing with $x$ alone gives $1\le\int x^2\int f^2=\tfrac13\int f^2$, the weaker bound $3$, because $x$ is not the function that equality wants.
\end{proof}

\begin{trap}
Applying an inequality in the direction that bounds the wrong way: AM-GM gives an upper bound on a product and a lower bound on a sum, never the reverse, and a chain of inequalities all pointing the same way is the only kind that proves anything. An extreme at the boundary rather than at the symmetric point: for $x+y+z=1$ with $x,y,z\ge 0$ the maximum of $xyz$ is at the centre but the maximum of $x^2+y^2+z^2$ is at a corner, and a Jensen argument that assumed the centre would be simply false. A convex function assumed where the second derivative changes sign, which is why the first example had to be checked by factoring. Forgetting the equality case in the write-up, when the problem asked for the best constant: the bound without the example achieving it is half the points. Dividing by an expression that may be zero, or multiplying through by one whose sign is not known, both of which silently flip or break the inequality.
\end{trap}

\begin{ladder}
\item \seminar{PS 5 \#3}. A ratio of a power to a factorial; which term of the sequence is largest.
\item \seminar{PS 5 \#4}. A finite product of factors slightly above one; a lower bound.
\item \seminar{PS 5 \#13}. Three positive reals satisfying one inequality; the same for their powers.
\item \putnam{2025 A2}. A sine bounded by a quadratic on an interval; the best constants.
\item \putnam{2018 B2}. A polynomial with decreasing coefficients; no roots in a disc.
\item \putnam{2021 B2}. A series of weighted geometric means under a sum constraint; its supremum.
\item \putnam{2018 A3}. A trigonometric sum under a constraint on another; its maximum.
\item \putnam{2019 A3}. Polynomials with increasing coefficients; the least possible mean root modulus.
\item \putnam{2017 A3}. A sequence of ratios of integrals; increasing and unbounded.
\item \putnam{2025 B4}. A zero-one matrix with monotone rows and columns; a sum bounded by a count.
\end{ladder}

\begin{hides}
Inside limits, where the squeeze needs an upper and a lower bound and each is an inequality (\putnam{2019 B2}, \putnam{2024 B3}). Inside polynomials, where a bound on the roots is a bound on an expression in the coefficients (\putnam{2018 B2}, \putnam{2019 A3}, \putnam{2016 A6}). Inside probability, where a concentration statement is a bound on an expectation or a tail (\putnam{2021 B6}, \putnam{2022 B5}). Inside functional equations, where growth conditions pin the exponent (\putnam{2025 B6}). Inside counting, where the extremal configuration is found by an inequality and then realised by a construction (\putnam{2025 B4}).
\end{hides}

\section{Polynomials and their roots}
\label{sec:polynomials}

A polynomial has three descriptions: its coefficients, its roots, and its values at enough points. Most polynomial problems are solved by moving from the description the problem gives to the one it hides.

\begin{recognise}
A polynomial is given by a property of its roots (``all roots real'', ``roots in arithmetic progression'') or of its values (``$p(k)$ is an integer for every integer $k$'', ``$p(1) = p(2) = 0$'') rather than by its coefficients. An expression in several variables is symmetric, so the variables are the roots of something. A sequence of values is asked about at the next index. A question about the real roots of $p$ mentions $p'$, or could. A sum of coefficients taken every $k$-th one is asked for. A recurrence or congruence is really a statement about a polynomial identity.
\end{recognise}

\subsection*{The model}

Throughout, $p(x) = a_n x^n + \dots + a_1 x + a_0$ with $a_n \ne 0$ has degree $n$, with coefficients in a field, which for the Putnam is $\Q$, $\R$, $\C$ or $\Z/p\Z$.

\begin{theorem}[Factor theorem]
$p(r) = 0$ if and only if $x - r$ divides $p(x)$.
\end{theorem}

\begin{proof}
Divide $p$ by $x - r$ with remainder: $p(x) = (x - r) q(x) + c$ with $c$ a constant. Setting $x = r$ gives $c = p(r)$.
\end{proof}

\begin{theorem}[Degree counting]
A nonzero polynomial of degree $n$ has at most $n$ roots, counted with multiplicity. Hence two polynomials of degree at most $n$ that agree at $n + 1$ points are identical, and a polynomial of degree at most $n$ with $n+1$ roots is the zero polynomial.
\end{theorem}

\begin{proof}
By the factor theorem, each root $r$ splits off a factor $x - r$ and lowers the degree by one, so after $n$ roots the quotient is a nonzero constant with no roots. For the second sentence, apply the first to the difference of the two polynomials.
\end{proof}

Over $\C$ every polynomial of degree $n$ has exactly $n$ roots with multiplicity; this is the fundamental theorem of algebra, used freely. So over $\C$, $p(x) = a_n \prod_{i=1}^{n} (x - r_i)$.

\begin{theorem}[Vieta]
If $r_1, \dots, r_n$ are the roots of $p$ with multiplicity, then for each $k$,
\[
\sum_{i_1 < \dots < i_k} r_{i_1} \cdots r_{i_k} = (-1)^k \frac{a_{n-k}}{a_n} .
\]
In particular $\sum r_i = -a_{n-1}/a_n$ and $\prod r_i = (-1)^n a_0 / a_n$.
\end{theorem}

\begin{proof}
Expand $a_n \prod (x - r_i)$ and compare the coefficient of $x^{n-k}$ with $a_{n-k}$.
\end{proof}

Vieta turns a symmetric expression in the roots into an expression in the coefficients. Power sums $\sum r_i^k$ are reached from the elementary symmetric functions either by Newton's identities or, faster on the day, by using $p(r_i) = 0$ to reduce $r_i^k$ to lower powers, as in the first example below.

\begin{theorem}[Rolle, for real roots]
Let $p$ have real coefficients. Between two distinct real roots of $p$ there is a real root of $p'$. Consequently, if $p$ has $r$ real roots counted with multiplicity, then $p'$ has at least $r - 1$; equivalently, $p$ has at most one more real root than $p'$.
\end{theorem}

\begin{proof}
The first sentence is Rolle's theorem from calculus. For the second, let the distinct real roots be $x_1 < \dots < x_s$ with multiplicities $m_1, \dots, m_s$ summing to $r$. A root of multiplicity $m_i \ge 2$ of $p$ is a root of multiplicity $m_i - 1$ of $p'$, and each of the $s - 1$ open gaps contains a further root of $p'$. So $p'$ has at least $\sum (m_i - 1) + (s - 1) = r - 1$ real roots.
\end{proof}

Used the other way round: if $p'$ has few real roots, so does $p$. If $p'$ has $k$ real roots then $p$ has at most $k + 1$, and the sign of $p$ at the roots of $p'$ decides how many of those are achieved.

\begin{theorem}[Lagrange interpolation]
Let $x_0, \dots, x_n$ be distinct and let $y_0, \dots, y_n$ be arbitrary. There is exactly one polynomial $p$ of degree at most $n$ with $p(x_i) = y_i$ for every $i$, namely
\[
p(x) = \sum_{i=0}^{n} y_i \prod_{j \ne i} \frac{x - x_j}{x_i - x_j} .
\]
\end{theorem}

\begin{proof}
Each product is a polynomial of degree $n$ equal to $1$ at $x_i$ and $0$ at the other nodes, so the sum has the right values. Uniqueness is degree counting: two such polynomials differ by a polynomial of degree at most $n$ with $n + 1$ roots.
\end{proof}

\begin{theorem}[Finite differences]
Define $\Delta p(x) = p(x + 1) - p(x)$. If $p$ has degree $n \ge 1$ with leading coefficient $a_n$, then $\Delta p$ has degree $n - 1$ with leading coefficient $n a_n$. Hence $\Delta^n p$ is the constant $n! \, a_n$, and $\Delta^{n+1} p = 0$, which written out at $x = 0$ reads
\[
\sum_{k=0}^{n+1} (-1)^{n+1-k} \binom{n+1}{k} p(k) = 0 .
\]
\end{theorem}

\begin{proof}
$(x+1)^n - x^n = n x^{n-1} + \dots$, and lower terms of $p$ contribute lower degrees. The displayed identity is $\Delta^{n+1} p(0)$ expanded by the binomial theorem applied to the operator $\Delta = E - 1$, where $E p(x) = p(x+1)$.
\end{proof}

The identity gives $p(n+1)$ from $p(0), \dots, p(n)$ without finding $p$. It is also the reason a polynomial that takes integer values at $0, 1, \dots, n$ takes integer values at every integer: the differences are integers, and $p(k)$ is a sum of binomial coefficients times differences.

\begin{theorem}[Roots-of-unity filter]
Let $\omega = e^{2\pi i/m}$. For any polynomial $p(x) = \sum a_k x^k$ and any residue $r$ modulo $m$,
\[
\sum_{k \equiv r \ (\mathrm{mod}\ m)} a_k = \frac{1}{m} \sum_{j=0}^{m-1} \omega^{-jr} \, p(\omega^j) .
\]
\end{theorem}

\begin{proof}
The right side is $\frac1m \sum_k a_k \sum_{j=0}^{m-1} \omega^{j(k-r)}$. The inner sum is $m$ when $m$ divides $k - r$ and otherwise is a geometric series with ratio $\omega^{k-r} \ne 1$ summing to $(\omega^{m(k-r)} - 1)/(\omega^{k-r} - 1) = 0$.
\end{proof}

With $m = 2$ this is the familiar $\sum a_{2k} = (p(1) + p(-1))/2$.

\begin{firstmoves}
Write the polynomial by its roots, with the leading coefficient in front, not by its coefficients. Evaluate at $0$, $1$, $-1$, and if anything is cyclic, at the roots of unity. Take a derivative and apply Rolle. Reduce the coefficients modulo $2$. For a polynomial given by values, write down Lagrange or the finite-difference table for the small case and see what it says. For a symmetric expression in several variables, name the variables as the roots of a polynomial and let Vieta translate. Count degrees on both sides of any identity before believing it.
\end{firstmoves}

\subsection*{Worked examples}

\begin{example}
Let $r$, $s$, $t$ be the three roots of $x^3 - 2x - 5$. Compute $r^5 + s^5 + t^5$.
\end{example}

\begin{proof}[Solution]
By Vieta, $r + s + t = 0$, $rs + st + tr = -2$, $rst = 5$. Then
\[
r^2 + s^2 + t^2 = (r + s + t)^2 - 2(rs + st + tr) = 0 + 4 = 4 .
\]
Each root satisfies its own equation, $r^3 = 2r + 5$, so
\[
r^5 = r^2 \cdot r^3 = r^2 (2r + 5) = 2 r^3 + 5 r^2 = 2(2r + 5) + 5 r^2 = 5 r^2 + 4 r + 10 ,
\]
and the same for $s$ and $t$. Summing,
\[
r^5 + s^5 + t^5 = 5 \cdot 4 + 4 \cdot 0 + 3 \cdot 10 = 50 .
\]
\end{proof}

Nothing was computed about the roots themselves; the polynomial has one real root and two complex ones, and it does not matter. The move ``reduce $r^k$ using $p(r) = 0$'' is the polynomial version of reducing modulo something, and it works for any power sum.

\begin{example}
Show that for every real $a$ the polynomial $x^5 - 5x + a$ has at most three real roots, and determine the values of $a$ for which it has exactly three distinct real roots.
\end{example}

\begin{proof}[Solution]
Let $p(x) = x^5 - 5x + a$. Then $p'(x) = 5x^4 - 5 = 5(x^2 - 1)(x^2 + 1)$, which has exactly two real roots, $x = -1$ and $x = 1$. By the Rolle count, $p$ has at most $2 + 1 = 3$ real roots with multiplicity.

For the second part, $p$ is increasing on $(-\infty, -1)$, decreasing on $(-1, 1)$, and increasing on $(1, \infty)$, with $p(x) \to -\infty$ as $x \to -\infty$ and $p(x) \to +\infty$ as $x \to +\infty$. Its local maximum is $p(-1) = -1 + 5 + a = a + 4$ and its local minimum is $p(1) = 1 - 5 + a = a - 4$. On each of the three monotone pieces $p$ has at most one root, and it has a root on all three exactly when the local maximum is positive and the local minimum is negative, that is, when $a + 4 > 0$ and $a - 4 < 0$. So $p$ has exactly three distinct real roots if and only if $-4 < a < 4$. At $a = \pm 4$ one root is double, and for $|a| > 4$ there is exactly one real root.
\end{proof}

The derivative had a simple factorisation and the problem fell to a sign table. Most ``how many real roots'' questions on the Putnam are this argument with a less convenient derivative, and the derivative's own real roots are found by the same method one level down.

\begin{example}
Let $p$ be a polynomial with real coefficients such that $p(x) = p(2 - x)$ for all real $x$. Show that $p(x)$ is a polynomial in $(x - 1)^2$.
\end{example}

\begin{proof}[Solution]
Put $q(y) = p(1 + y)$, a polynomial in $y$. The hypothesis at $x = 1 + y$ says $q(y) = p(1 + y) = p(1 - y) = q(-y)$ for all real $y$. Write $q(y) = \sum c_k y^k$. Then $q(y) - q(-y) = 2 \sum_{k \text{ odd}} c_k y^k$ is a polynomial vanishing at every real $y$, so by degree counting it is the zero polynomial and every odd coefficient $c_k$ is $0$. Hence $q(y) = \sum_j c_{2j} y^{2j}$ and $p(x) = q(x - 1) = \sum_j c_{2j} (x - 1)^{2j}$.
\end{proof}

A functional identity between polynomials becomes a statement about coefficients only through degree counting: ``equal at infinitely many points'' is what licenses ``equal as polynomials''. The shift to the symmetric point is the standard first move for any identity of the form $p(x) = p(c - x)$.

\begin{trap}
Roots that are complex or repeated. Vieta counts with multiplicity and over $\C$; ``the roots are real'' is a hypothesis to be used, not assumed. A polynomial identity checked at finitely many points without counting degrees proves nothing. The leading coefficient forgotten when writing $p$ by its roots: $p(x) = a_n \prod (x - r_i)$, and $a_n$ is often what the problem is about. Rolle gives an upper bound on real roots, never a lower one; the lower bound needs sign changes. Lagrange interpolation produces the unique polynomial of degree at most $n$; a polynomial of higher degree through the same points is not determined. The finite-difference identity needs equally spaced nodes.
\end{trap}

\begin{ladder}
\item \putnam{2019 A1}. A symmetric cubic expression in three integers; which values it takes. Factor the symmetric form.
\item \putnam{2023 A2}. An even-degree polynomial with prescribed values at prescribed points; find its other roots. Build $p$ from its values.
\item \seminar{Roots \#6}. A polynomial with coefficients from a two-element set; can all its roots be real. Vieta against an inequality.
\item \seminar{Roots \#7}. A polynomial derived from one with all roots real; show the derived one also has. Rolle and sign changes.
\item \putnam{2016 A1}. Derivatives of an integer polynomial; the smallest order always divisible by a given integer. Polynomials meet congruences.
\item \putnam{2022 B1}. A power series built from an integer polynomial; show its coefficients are all nonzero. Work modulo $2$.
\item \putnam{2024 A2}. Iterates of a polynomial and a divisibility relation; which polynomials. Degree counting on the iterate.
\item \putnam{2017 A2}. A rational recurrence whose terms are integer polynomials. The polynomial identity behind the recurrence.
\item \putnam{2019 B5}. A polynomial that interpolates a famous recurrence at many points; its value at the next one. Finite differences.
\item \putnam{2018 B2}. A polynomial with decreasing coefficients; no roots in a disc. A bound on the roots through a telescoping product.
\item \putnam{2019 A3}. Polynomials with increasing coefficients; the smallest mean of the root moduli. Vieta and the AM--GM inequality.
\item \putnam{2022 A2}. The square of a polynomial; how many coefficients can be negative. Extremal on the coefficient pattern.
\end{ladder}

Seminar: \seminar{Roots \#1, \#8, \#11a}; the set notes that \seminar{\#11(b)} is considerably harder, that \seminar{\#23} is the Chebyshev bound, and that \seminar{\#25} is Descartes's rule of signs.

\begin{hides}
Recurrences (\putnam{2017 A2}, \putnam{2019 B5}): a sequence of polynomials, or a polynomial through a sequence's values, is a polynomial problem with a recurrence in front. Congruences (\putnam{2016 A1}, \putnam{2022 B1}): an integer polynomial reduced modulo a prime is still a polynomial, with at most as many roots as its degree, which is the bound most such problems need. Inequalities: a symmetric inequality in $n$ variables is often Vieta applied to the polynomial with those roots. Linear algebra: the characteristic polynomial, and the space of polynomials of degree at most $n$ as a vector space of dimension $n + 1$. Generating functions: a generating function is a polynomial or power series, and extracting every $k$-th coefficient is the roots-of-unity filter.
\end{hides}


\section{Congruences and number theory}
\label{sec:numbertheory}

The largest tag in the easy slots. The tools are few and old, and the problems are mostly about choosing what to divide by. This section starts from division with remainder and defines every word before it is used.

\begin{recognise}
Integers, divisibility, remainders, digits, primes, factorials, greatest common divisors. An equation to be solved in integers. A polynomial or a recurrence asked about ``modulo'' something, which means only its remainders matter. A count of integers with a divisibility property. Anything with ``for all $n$'' where $n$ is an exponent. A statement that is false over the reals and true over the integers is a number theory statement.
\end{recognise}

\subsection*{The model, from the ground}

\paragraph{Dividing and remainders.} For integers $d \ne 0$ and $m$, the phrase ``$d$ divides $m$'', written $d \mid m$, means $m = dk$ for some integer $k$. Written $d \nmid m$, it does not. So $3 \mid 12$ and $3 \nmid 14$, and every nonzero integer divides $0$. One fact is used constantly: if $d$ divides $a$ and $b$, then $d$ divides $a + b$, $a - b$ and every multiple of $a$.

For any integer $a$ and any integer $n \ge 1$ there is exactly one way to write
\[
a = qn + r \qquad \text{with integers } q, r \text{ and } 0 \le r < n .
\]
The number $r$ is the \emph{remainder} of $a$ on division by $n$. For $n = 5$: $17 = 3 \cdot 5 + 2$ and $-3 = (-1) \cdot 5 + 2$, so $17$ and $-3$ both have remainder $2$.

\begin{definition}[Congruence]
Fix an integer $n \ge 1$, called the \emph{modulus}. Two integers $a$ and $b$ are \emph{congruent modulo $n$}, written $a \equiv b \pmod n$, when $n$ divides $a - b$.
\end{definition}

Congruent means the same remainder. Write $a = qn + r$ and $b = q'n + r'$. Then $a - b = (q - q')n + (r - r')$, and $r - r'$ lies strictly between $-n$ and $n$, so $n$ divides $a - b$ exactly when $r - r' = 0$. So $17 \equiv 2 \equiv -3 \pmod 5$. A clock is arithmetic modulo $12$: five hours after ten o'clock is three o'clock, because $15 \equiv 3 \pmod{12}$.

\paragraph{Classes.} Congruence modulo $n$ sorts the integers into $n$ piles, one for each possible remainder $0, 1, \dots, n - 1$. Each pile is called a \emph{class} modulo $n$, or a \emph{residue class}; \emph{residue} is another word for remainder. For $n = 5$ the five piles are

\begin{center}\small
\begin{tabular}{ll}
\toprule
Remainder & The pile \\
\midrule
$0$ & $\dots, -10, -5, 0, 5, 10, \dots$ \\
$1$ & $\dots, -9, -4, 1, 6, 11, \dots$ \\
$2$ & $\dots, -8, -3, 2, 7, 12, \dots$ \\
$3$ & $\dots, -7, -2, 3, 8, 13, \dots$ \\
$4$ & $\dots, -6, -1, 4, 9, 14, \dots$ \\
\bottomrule
\end{tabular}
\end{center}

Every integer is in exactly one pile, and two integers are congruent exactly when they are in the same pile. Working modulo $n$ means treating a whole pile as one object and naming it by any of its members, usually the remainder.

That is all the phrase ``an equivalence relation with $n$ classes'' says. An \emph{equivalence relation} is any rule for calling two things the same for present purposes that obeys three laws: each thing is the same as itself; if $a$ is the same as $b$ then $b$ is the same as $a$; and if $a$ is the same as $b$ and $b$ the same as $c$, then $a$ is the same as $c$. Any such rule sorts things into piles, called its \emph{equivalence classes}. Congruence obeys the three laws: $n$ divides $a - a = 0$; if $n$ divides $a - b$ it divides $b - a$; and if $n$ divides $a - b$ and $b - c$ it divides their sum $a - c$. Its piles are the $n$ classes in the table.

\begin{proposition}[Arithmetic respects classes]
If $a \equiv b$ and $c \equiv d \pmod n$, then $a + c \equiv b + d$, $a - c \equiv b - d$ and $ac \equiv bd \pmod n$. Consequently $a^k \equiv b^k \pmod n$ for every $k \ge 1$.
\end{proposition}

\begin{proof}
By assumption $a - b$ and $c - d$ are multiples of $n$. Then $(a + c) - (b + d) = (a - b) + (c - d)$ is a sum of two multiples of $n$, and the difference is handled the same way. For the product, $ac - bd = a(c - d) + d(a - b)$, again a sum of two multiples of $n$. The statement about powers is the product rule applied $k - 1$ times.
\end{proof}

So in any expression built by adding, subtracting and multiplying integers, any number may be replaced at any moment by another member of its pile, and the pile of the result does not change. Three uses.

\begin{itemize}[nosep,leftmargin=*]
\item The remainder of $2^{100}$ on division by $7$. Since $2^3 = 8 \equiv 1$, we get $2^{100} = (2^3)^{33} \cdot 2 \equiv 1^{33} \cdot 2 = 2 \pmod 7$. Nobody computed $2^{100}$.
\item The last digit of an integer is its class modulo $10$.
\item An integer is congruent modulo $9$ to the sum of its digits, because $10 \equiv 1 \pmod 9$ gives $10^k \equiv 1$. So $4725 \equiv 4 + 7 + 2 + 5 = 18 \equiv 0 \pmod 9$, and $9$ divides $4725$.
\end{itemize}

\paragraph{Division, inverses and units.} Division is the exception. Modulo $6$, $2 \cdot 3 \equiv 2 \cdot 0$, since both sides are multiples of $6$, yet $3 \not\equiv 0$. The factor $2$ could not be cancelled. To see when cancelling is allowed, ask what dividing means. Among real numbers, dividing by $c$ is multiplying by $1/c$, the number whose product with $c$ is $1$. The same idea is used for classes.

\begin{definition}[Inverse, unit]
An \emph{inverse} of $c$ modulo $n$ is an integer $x$ with $cx \equiv 1 \pmod n$. If one exists, $c$ is called \emph{invertible} modulo $n$, or a \emph{unit} modulo $n$.
\end{definition}

Modulo $7$: $3 \cdot 5 = 15 \equiv 1$, so $5$ is an inverse of $3$. Also $1 \cdot 1$, $2 \cdot 4 = 8$ and $6 \cdot 6 = 36$ are all congruent to $1$, so each of $1, 2, 3, 4, 5, 6$ is a unit modulo $7$. Modulo $6$: $1 \cdot 1 \equiv 1$ and $5 \cdot 5 = 25 \equiv 1$, so $1$ and $5$ are units. But $2x$ and $4x$ are always even and $3x$ is always a multiple of $3$, while a number congruent to $1$ modulo $6$ is neither, so $2$, $3$ and $4$ have no inverse. Being a unit is a property of the whole pile, because replacing $c$ by another member of its pile does not change the pile of $cx$.

A unit can be cancelled. If $ac \equiv bc \pmod n$ and $x$ is an inverse of $c$, multiply both sides by $x$: $a \equiv acx \equiv bcx \equiv b$.

Which numbers are units is decided by common factors. The \emph{greatest common divisor} $\gcd(a, b)$ is the largest positive integer dividing both $a$ and $b$. The two are \emph{coprime} when $\gcd(a, b) = 1$, that is, when no prime divides both.

\begin{lemma}[B\'ezout]
For integers $c$ and $n \ge 1$ there are integers $x$ and $y$ with $cx + ny = \gcd(c, n)$.
\end{lemma}

\begin{proof}
Look at all the positive integers of the form $cx + ny$ with $x, y$ integers. There is at least one, $n$ itself. Let $d = c x_0 + n y_0$ be the smallest. Divide $c$ by $d$: $c = qd + r$ with $0 \le r < d$. Then $r = c - qd = c(1 - q x_0) + n(-q y_0)$ is of the same form and is smaller than $d$, so it cannot be positive. Hence $r = 0$ and $d$ divides $c$. The same argument shows $d$ divides $n$. So $d$ is a common divisor. And every common divisor of $c$ and $n$ divides $c x_0 + n y_0 = d$, so none is larger than $d$. Therefore $d = \gcd(c, n)$.
\end{proof}

\begin{proposition}[Which numbers are units]
$c$ is invertible modulo $n$ exactly when $\gcd(c, n) = 1$.
\end{proposition}

\begin{proof}
If $\gcd(c, n) = 1$, B\'ezout gives integers with $cx + ny = 1$, so $cx - 1$ is a multiple of $n$, which is $cx \equiv 1 \pmod n$. Conversely, if $cx \equiv 1$ then $cx - 1 = nk$ for some integer $k$, so $cx - nk = 1$. A common divisor of $c$ and $n$ divides the left side, hence divides $1$, so it is $1$.
\end{proof}

For $c = 3$ and $n = 7$ the B\'ezout identity is $3 \cdot 5 + 7 \cdot (-2) = 1$, which is the inverse $5$ found above.

\begin{definition}[Euler's function]
$\varphi(n)$ is the number of integers among $1, 2, \dots, n$ that are coprime to $n$. By the proposition it is the number of invertible classes modulo $n$.
\end{definition}

\begin{center}\small
\begin{tabular}{lll}
\toprule
$n$ & Units modulo $n$ & $\varphi(n)$ \\
\midrule
$5$ & $1, 2, 3, 4$ & $4$ \\
$6$ & $1, 5$ & $2$ \\
$7$ & $1, 2, 3, 4, 5, 6$ & $6$ \\
$8$ & $1, 3, 5, 7$ & $4$ \\
$9$ & $1, 2, 4, 5, 7, 8$ & $6$ \\
$10$ & $1, 3, 7, 9$ & $4$ \\
$12$ & $1, 5, 7, 11$ & $4$ \\
\bottomrule
\end{tabular}
\end{center}

For a prime $p$, every one of $1, 2, \dots, p - 1$ is coprime to $p$, because a prime has no divisors other than $1$ and itself. So $\varphi(p) = p - 1$. This is what ``for a prime $p$ every nonzero class is invertible'' says. The class of $0$ is the pile of multiples of $p$. A nonzero class is any other pile. Each of those has an inverse. Modulo a prime you may divide by anything that is not congruent to $0$, exactly as with real numbers.

\paragraph{The word group.} Under multiplication, the units modulo $n$ have four properties.

\begin{enumerate}[nosep,leftmargin=*]
\item The product of two units is a unit: if $cx \equiv 1$ and $dy \equiv 1$ then $(cd)(xy) \equiv 1$.
\item Multiplication does not care about bracketing: $(ab)c = a(bc)$.
\item The number $1$ is a unit, and multiplying by it changes nothing.
\item Every unit has an inverse, and that inverse is again a unit.
\end{enumerate}

A set with an operation that has these four properties is called a \emph{group}, and the number of its members is its \emph{size}. So the sentence ``the classes coprime to $n$ form a group under multiplication, of size $\varphi(n)$'' says exactly this: there are $\varphi(n)$ invertible classes, the product of two of them is another, and each can be divided by. Nothing more is meant, and nothing about groups beyond this list is used in this section. The multiplication table of the units modulo $8$ shows the four properties at work:

\begin{center}\small
\begin{tabular}{c|cccc}
$\times$ & $1$ & $3$ & $5$ & $7$ \\
\midrule
$1$ & $1$ & $3$ & $5$ & $7$ \\
$3$ & $3$ & $1$ & $7$ & $5$ \\
$5$ & $5$ & $7$ & $1$ & $3$ \\
$7$ & $7$ & $5$ & $3$ & $1$ \\
\end{tabular}
\end{center}

Every row contains each unit exactly once. Multiplying all the units by one fixed unit only shuffles them. That observation is the whole proof of the next theorem.

\begin{theorem}[Fermat and Euler]
If $p$ is prime and $p \nmid a$, then $a^{p-1} \equiv 1 \pmod p$. More generally, if $\gcd(a, n) = 1$ then $a^{\varphi(n)} \equiv 1 \pmod n$.
\end{theorem}

\begin{proof}
One instance first, $n = 7$ and $a = 3$. The units are $1, 2, 3, 4, 5, 6$. Multiplied by $3$ they become $3, 6, 9, 12, 15, 18$, which are congruent to $3, 6, 2, 5, 1, 4$: the same six classes in another order. So the two products agree, $(3 \cdot 1)(3 \cdot 2) \cdots (3 \cdot 6) \equiv 1 \cdot 2 \cdots 6$, that is, $3^6 \cdot 720 \equiv 720 \pmod 7$. The number $720$ is a product of units, so it is a unit and can be cancelled, leaving $3^6 \equiv 1$. Indeed $3^6 = 729 = 104 \cdot 7 + 1$.

In general, let $u_1, \dots, u_k$ be the units modulo $n$, one from each invertible class, so $k = \varphi(n)$. Each $a u_i$ is a unit, by the first property above. If $a u_i \equiv a u_j$ then $u_i \equiv u_j$, by cancelling the unit $a$. So $a u_1, \dots, a u_k$ are $k$ units lying in $k$ different classes, which makes them all of the units in some order. Multiplying them together, $a^k U \equiv U \pmod n$ where $U = u_1 u_2 \cdots u_k$. The product $U$ is a unit, and cancelling it gives $a^k \equiv 1$. For a prime $p$, $k = p - 1$, and $\gcd(a, p) = 1$ is the same as $p \nmid a$.
\end{proof}

\begin{definition}[Order]
For $\gcd(a, n) = 1$, the \emph{order} of $a$ modulo $n$, written $\ord_n(a)$, is the smallest positive integer $d$ with $a^d \equiv 1 \pmod n$. Euler's theorem guarantees that such a $d$ exists and is at most $\varphi(n)$.
\end{definition}

Modulo $7$ the powers of $2$ are $2, 4, 1$, so $\ord_7(2) = 3$. The powers of $3$ are $3, 2, 6, 4, 5, 1$, so $\ord_7(3) = 6$. The powers of $6$ are $6, 1$, so $\ord_7(6) = 2$. After reaching $1$ the powers repeat, since $a^{m + d} = a^m a^d \equiv a^m$. The three orders $3$, $6$ and $2$ all divide $6 = \varphi(7)$, and that is no accident.

\begin{lemma}[Orders]
$a^m \equiv 1 \pmod n$ if and only if $\ord_n(a)$ divides $m$. In particular $\ord_n(a)$ divides $\varphi(n)$, and for a prime $p$ it divides $p - 1$.
\end{lemma}

\begin{proof}
Write $d = \ord_n(a)$ and divide $m$ by $d$: $m = qd + r$ with $0 \le r < d$. Then $a^m = (a^d)^q a^r \equiv 1^q a^r = a^r$. So $a^m \equiv 1$ exactly when $a^r \equiv 1$. Since $r < d$ and $d$ is the smallest positive exponent that gives $1$, this happens only for $r = 0$, which is the statement that $d$ divides $m$. The particular case takes $m = \varphi(n)$, where Euler's theorem says $a^m \equiv 1$.
\end{proof}

A unit whose order is the full $\varphi(n)$, so that its powers run through every unit, is called a \emph{primitive root} modulo $n$. Modulo $7$, $3$ is a primitive root and $2$ is not. Every prime has a primitive root. That fact is not proved here, and it is needed only on the harder rungs of the ladder.

\begin{lemma}[Smallest prime factor]
Let $n > 1$ and let $p$ be the smallest prime factor of $n$. Then $\gcd(n, p - 1) = 1$. Consequently, if $a^n \equiv 1 \pmod p$ and $p \nmid a$, then $a \equiv 1 \pmod p$.
\end{lemma}

\begin{proof}
Every prime factor of $p - 1$ is at most $p - 1$, so it is smaller than $p$. If it divided $n$, it would be a prime factor of $n$ smaller than the smallest one. So $n$ and $p - 1$ have no common prime factor, and $\gcd(n, p - 1) = 1$.

For the consequence, let $d = \ord_p(a)$. By the lemma on orders, $d$ divides $n$, because $a^n \equiv 1$. By the same lemma $d$ divides $p - 1$. B\'ezout gives integers with $nx + (p - 1)y = 1$, and $d$ divides the left side, so $d$ divides $1$. Hence $d = 1$, which says $a^1 \equiv 1 \pmod p$.
\end{proof}

This lemma is the standard tool against statements of the form ``$n$ never divides $a^n - 1$'' or ``$n$ divides $a^n - b^n$ only when''. Taking the smallest prime factor turns a statement about $n$ into one about a single prime where the order is pinned down.

\begin{theorem}[Chinese remainder theorem]
If $\gcd(m, n) = 1$, then for any integers $a$ and $b$ there is exactly one class $x$ modulo $mn$ with $x \equiv a \pmod m$ and $x \equiv b \pmod n$.
\end{theorem}

In words: when $m$ and $n$ are coprime, a remainder modulo $m$ and a remainder modulo $n$ can be chosen freely, and together they fix the remainder modulo $mn$. For $m = 3$ and $n = 5$, here is each $x$ from $0$ to $14$ with its pair of remainders (modulo $3$, modulo $5$):

\begin{center}\small
\begin{tabular}{lllll}
$0\colon (0,0)$ & $1\colon (1,1)$ & $2\colon (2,2)$ & $3\colon (0,3)$ & $4\colon (1,4)$ \\
$5\colon (2,0)$ & $6\colon (0,1)$ & $7\colon (1,2)$ & $8\colon (2,3)$ & $9\colon (0,4)$ \\
$10\colon (1,0)$ & $11\colon (2,1)$ & $12\colon (0,2)$ & $13\colon (1,3)$ & $14\colon (2,4)$ \\
\end{tabular}
\end{center}

All fifteen possible pairs appear, each exactly once. The number with remainder $2$ modulo $3$ and remainder $3$ modulo $5$ is $8$, and every other such number is congruent to $8$ modulo $15$.

\begin{proof}
Send each class $x$ modulo $mn$ to the pair of its remainders modulo $m$ and modulo $n$. There are $mn$ classes and $mn$ possible pairs. Two different classes never give the same pair. Suppose $x$ and $x'$ agree in both remainders. Then $m$ and $n$ both divide $x - x'$. Write $x - x' = mk$. B\'ezout gives $mu + nv = 1$, and multiplying by $k$ gives $k = mku + nkv$. Here $n$ divides $mk$ and so divides the first term, and it visibly divides the second, so $n$ divides $k$. Hence $mn$ divides $mk = x - x'$, and $x$ and $x'$ are the same class. So $mn$ classes go to $mn$ pairs with no pair used twice, which means every pair is used exactly once.
\end{proof}

So a congruence modulo a composite number is a system of congruences modulo prime powers, solved one prime at a time, and a count of classes modulo $mn$ with a property that can be checked separately modulo $m$ and modulo $n$ is a product of two counts.

\begin{theorem}[Lifting a simple root]
Let $f$ be a polynomial with integer coefficients and $f'$ its derivative, computed term by term by the usual rule. Let $p$ be a prime and $a$ an integer with $f(a) \equiv 0 \pmod p$ and $f'(a) \not\equiv 0 \pmod p$. Then there is exactly one $t$ modulo $p$ such that $f(a + tp) \equiv 0 \pmod{p^2}$.
\end{theorem}

An integer $a$ with $f(a) \equiv 0 \pmod p$ is called a \emph{root of $f$ modulo $p$}. The theorem says that a root modulo $p$ can be corrected, in exactly one way, into a root modulo $p^2$.

\begin{proof}
First, for integers $a$ and $h$ and any $k \ge 1$, the binomial theorem gives $(a + h)^k = a^k + k a^{k-1} h + (\text{terms containing } h^2)$, so $(a + h)^k \equiv a^k + k a^{k-1} h \pmod{h^2}$. Multiplying by the coefficients of $f$ and adding over its terms, $f(a + h) \equiv f(a) + h f'(a) \pmod{h^2}$.

Now put $h = tp$, so that $h^2$ is a multiple of $p^2$: $f(a + tp) \equiv f(a) + tp f'(a) \pmod{p^2}$. Since $p$ divides $f(a)$, write $f(a) = pm$. The right side is $p(m + t f'(a))$, and it is a multiple of $p^2$ exactly when $p$ divides $m + t f'(a)$, that is, when $t f'(a) \equiv -m \pmod p$. Because $f'(a)$ is not congruent to $0$ modulo the prime $p$, it is a unit, with an inverse $x$. Multiplying by $x$ gives $t \equiv -mx \pmod p$, exactly one class.
\end{proof}

An instance: $f(x) = x^2 - 2$, $p = 7$, $a = 3$. Here $f(3) = 7$, so $m = 1$, and $f'(3) = 6$. The condition is $6t \equiv -1 \equiv 6 \pmod 7$, so $t = 1$ and the corrected root is $3 + 7 = 10$. Indeed $10^2 - 2 = 98 = 2 \cdot 49$.

The same step lifts from $p^k$ to $p^{k+1}$, so a simple root modulo $p$ gives a root modulo every power of $p$. A root with $f'(a) \equiv 0$ may lift to several classes or to none, which is the case Putnam problems like to hide.

\paragraph{Primes inside a factorial.} Three pieces of notation. For a prime $p$ and a positive integer $m$, $v_p(m)$ is the exponent of $p$ in $m$: the largest $e$ such that $p^e$ divides $m$. So $v_2(40) = 3$, since $40 = 2^3 \cdot 5$. The \emph{floor} $\lfloor x \rfloor$ is the largest integer not exceeding $x$, so $\lfloor 10/4 \rfloor = 2$. And $s_p(n)$ is the sum of the digits of $n$ written in base $p$.

\begin{theorem}[The exponent of a prime in a factorial (Legendre's formula)]
For a prime $p$ and $n \ge 1$,
\[
v_p(n!) = \sum_{k \ge 1} \left\lfloor \frac{n}{p^k} \right\rfloor = \frac{n - s_p(n)}{p - 1}.
\]
\end{theorem}

For $n = 10$ and $p = 2$ the sum is $\lfloor 10/2 \rfloor + \lfloor 10/4 \rfloor + \lfloor 10/8 \rfloor = 5 + 2 + 1 = 8$. In base $2$, ten is $1010$, with digit sum $2$, and $(10 - 2)/(2 - 1) = 8$ as well. Indeed $10! = 3628800 = 2^8 \cdot 14175$ with $14175$ odd.

\begin{proof}
Among $1, 2, \dots, n$ there are $\lfloor n/p \rfloor$ multiples of $p$, $\lfloor n/p^2 \rfloor$ multiples of $p^2$, and so on. A number whose exponent of $p$ is exactly $j$ is counted $j$ times, once as a multiple of each of $p, p^2, \dots, p^j$, and $j$ is what it contributes to $v_p(n!)$. That is the sum. For the closed form, write $n = \sum_i d_i p^i$ in base $p$. Then $\lfloor n / p^k \rfloor = \sum_{i \ge k} d_i p^{i-k}$, and summing over $k \ge 1$, the digit $d_i$ is multiplied by $p^{i-1} + \dots + p + 1 = (p^i - 1)/(p - 1)$. So the total is $\sum_i d_i (p^i - 1)/(p - 1) = (n - s_p(n))/(p - 1)$.
\end{proof}

\paragraph{Equations in integers.} A \emph{Diophantine equation} is an equation whose solutions are required to be integers. For these the model is a procedure rather than a theorem: factor what factors, bound what can be bounded, and reduce modulo $3$, $4$, $5$, $7$, $8$, $9$ until the cases die. Reducing works because a square can only fall in a few classes. The list for a modulus $m$ is found by squaring $0, 1, \dots, m - 1$ and taking remainders. Modulo $5$: $0, 1, 4, 9, 16$ give $0, 1, 4, 4, 1$, so a square is $0$, $1$ or $4$ modulo $5$ and never $2$ or $3$. The lists worth knowing cold: squares are $0, 1 \pmod 3$; $0, 1 \pmod 4$; $0, 1, 4 \pmod 5$; $0, 1, 4 \pmod 8$; $0, 1, 4, 7 \pmod 9$. Cubes are $0, \pm 1 \pmod 7$ and $0, \pm 1 \pmod 9$. Fourth powers are $0, 1 \pmod{16}$ and $0, 1 \pmod 5$.

\begin{firstmoves}
The smallest cases by hand, and more of them than feels necessary: a solution at $n = 2$ that does not persist is the shape of the whole answer. Reduce modulo $4$ and modulo $3$ first, then modulo $8$ and $9$. Factor anything that factors, including differences of powers and forms like $x^2 + x + 1$. For an exponent in the statement, compute the order of the base modulo the small primes involved. For a factorial, apply Legendre's formula. For ``find all $n$'', separate $n$ prime from $n$ composite, and for composite $n$ take the smallest prime factor.
\end{firstmoves}

\subsection*{Worked examples}

\begin{example}
Let $n$ be a positive integer. Show that every prime factor of $n^4 + n^2 + 1$ is either $3$ or congruent to $1$ modulo $6$.
\end{example}

\begin{proof}[Solution]
Let $p$ be a prime dividing $N = n^4 + n^2 + 1$. First, $p \nmid n$, since $p \mid n$ would give $p \mid N - n^4 - n^2 = 1$. Second, $p$ is odd, because $N$ is odd: if $n$ is even then $N$ is even plus even plus $1$, and if $n$ is odd then $N$ is a sum of three odd numbers.

Now $(n^2 - 1) N = n^6 - 1$, so $n^6 \equiv 1 \pmod p$, and by the lemma on orders $d = \ord_p(n)$ divides $6$. If $d$ is $1$ or $2$, then $n^2 \equiv 1 \pmod p$, so $N \equiv 1 + 1 + 1 = 3 \pmod p$. Since $p$ divides $N$, this gives $p \mid 3$, so $p = 3$. If $d$ is $3$ or $6$, then $3 \mid d$, and $d \mid p - 1$ by the lemma on orders, so $3 \mid p - 1$. Since $p$ is odd, also $2 \mid p - 1$, and therefore $6 \mid p - 1$, which says $p \equiv 1 \pmod 6$.
\end{proof}

Both branches occur: $n = 1$ gives $3$, and $n = 2$ gives $21 = 3 \times 7$. The argument is the template for every ``prime factors of $f(n)$'' statement: find an identity that makes some power of $n$ congruent to $1$ modulo the prime, then let the order do the sorting.

\begin{example}
Find all pairs of positive integers $(x, y)$ with $x^2 = y! + 7$.
\end{example}

\begin{proof}[Solution]
Small cases first. $y = 1$: $1 + 7 = 8$, not a square. $y = 2$: $2 + 7 = 9 = 3^2$, a solution. $y = 3$: $6 + 7 = 13$, no. $y = 4$: $24 + 7 = 31$, no.

For $y \ge 5$, the factorial $y!$ is divisible by $5$, so $x^2 \equiv 7 \equiv 2 \pmod 5$. But the squares modulo $5$ are $0$, $1$ and $4$. So there is no solution with $y \ge 5$, and the only pair is $(x, y) = (3, 2)$.
\end{proof}

The modulus was chosen to make the factorial vanish and to leave the constant in a class that is not a square. Modulo $3$ would not have worked ($7 \equiv 1$ is a square), nor modulo $4$ ($7 \equiv 3$, not a square, but $y!$ vanishes modulo $4$ only from $y = 4$, which still works here and is worth noticing as the alternative). The habit is to try the small moduli in turn and keep the first one that kills the equation, then check that every $y$ below the threshold was handled by hand.

\begin{trap}
Dividing a congruence by something not coprime to the modulus: $2 \cdot 3 \equiv 2 \cdot 0 \pmod 6$ does not give $3 \equiv 0$. A modulus that kills the information needed: reducing $x^2 + y^2 = z^2$ modulo $2$ says nothing, modulo $4$ says a great deal. Forgetting that $2$ is prime, so that an argument ``for every odd prime'' has a case left over, and forgetting that $p = 3$ behaves differently from other primes in anything involving $n^2 + n + 1$ or sixth powers. Euler's theorem used with $\gcd(a, n) \ne 1$. Applying Fermat with the modulus and the exponent swapped. A lift from $p$ to $p^2$ asserted for a root where the derivative vanishes.
\end{trap}

\begin{ladder}
\item \putnam{2018 A1}. An equation in positive integers built from unit fractions; find all solutions. Clear denominators, factor, bound.
\item \putnam{2017 B2}. Sums of consecutive integers; the smallest first term with a property. Parity and divisibility of a sum of an arithmetic progression.
\item \seminar{PS 9 \#14}. A classic statement that an integer never divides a certain exponential expression. Smallest prime factor and orders.
\item \seminar{PS 9 \#15}. The companion statement for odd integers and a shifted expression. The same lemma, one twist.
\item \putnam{2020 A1}. Integers whose decimal digits follow a block pattern; how many are divisible by a given integer. Factor the pattern as a product, then orders.
\item \putnam{2024 A1}. An exponential equation in integers; for which exponents it has solutions. Reduce modulo small numbers until the cases die.
\item \putnam{2021 A1}. Fixed-length hops in the plane; the minimum number to reach a given point. Tagged pigeonhole; the modular obstruction on the hop lengths is the number theory.
\item \putnam{2023 B2}. The binary digits of the multiples of an integer; a minimum over the multiples. Structure of multiples with few ones.
\item \putnam{2018 B3}. Integers satisfying linked divisibility conditions; find all. Orders and the size of exponents.
\item \putnam{2024 A4}. Powers rearranged into an arithmetic progression modulo a prime; which primes. Primitive roots and residues.
\item \putnam{2021 A5}. Power sums over the residues coprime to a modulus; which exponents give divisibility. The group structure of the units.
\item \putnam{2025 B3}. A set of integers closed under divisors, built from exponential differences; whether it is everything. Orders modulo each candidate.
\end{ladder}

Seminar: \seminar{PS 9}, with \seminar{\#1, \#4, \#12, \#31} as warm-ups and \seminar{\#16, \#17, \#19} for power towers, where Euler's theorem is applied to the exponent. \seminar{Floor \#2} is the exponent of a prime in a factorial; \seminar{Floor \#12} to \seminar{\#15} are Putnam floor-function problems with their score statistics printed, which is a useful calibration.

\begin{hides}
Polynomials (week 2 of the ladders): an integer polynomial modulo a prime has at most its degree in roots, and a power series with integer coefficients can be read modulo $2$. Counting: the binary digits of a count (\putnam{2020 B1}) are a number theory statement about a sum. Games with integers (\putnam{2023 A6}): the winning condition is a congruence. Recurrences: a sequence modulo $m$ is eventually periodic, and the period is an order. Abstract algebra: the units modulo $n$ are the first example of a group, and Euler's theorem is Lagrange's theorem for that group.
\end{hides}

\section{Recurrences}

\begin{recognise}
A sequence is defined from its earlier terms and something is asked about a far term, about divisibility of the terms, about their growth or their limit. A count for $n$ splits according to what happens first or last, which builds it from counts for smaller $n$. A sequence of polynomials, matrices or functions is defined by iteration. A problem about ``the $2026$th term'' is a problem about the period. A rational recurrence whose terms stay integers, or a nonlinear one whose terms stay bounded, is almost always a linear one in disguise.
\end{recognise}

\begin{definition}
A linear recurrence of order $k$ with constant coefficients is $a_{n+k}=c_{k-1}a_{n+k-1}+\dots+c_0a_n$ for all $n\ge 0$, with $c_0\ne 0$. Its characteristic polynomial is $p(x)=x^k-c_{k-1}x^{k-1}-\dots-c_0$. Writing $E$ for the shift $(Ea)_n=a_{n+1}$, the recurrence says $p(E)a=0$.
\end{definition}

\begin{theorem}[Solving a linear recurrence]
The sequences satisfying a linear recurrence of order $k$ form a vector space $V$ of dimension $k$. If $p(x)=\prod_{i=1}^s(x-r_i)^{m_i}$ with distinct $r_i\ne 0$, then the $k$ sequences $n^jr_i^{\,n}$, $0\le j<m_i$, form a basis of $V$. So every solution is
\[
a_n=\sum_{i=1}^s P_i(n)\,r_i^{\,n},\qquad \deg P_i<m_i,
\]
and when all roots are simple, $a_n=\sum_i C_ir_i^{\,n}$ with constants $C_i$ fixed by the first $k$ terms.
\end{theorem}

\begin{proof}
A solution is determined by $(a_0,\dots,a_{k-1})$ and every choice of those values gives a solution, so the map $V\to\R^k$ (or $\C^k$) sending a solution to its first $k$ terms is a linear bijection and $\dim V=k$.

For a single factor: $(E-r)$ applied to $n^jr^n$ gives $(n+1)^jr^{n+1}-r\cdot n^jr^n=r^{n+1}\bigl((n+1)^j-n^j\bigr)$, and $(n+1)^j-n^j$ is a polynomial in $n$ of degree $j-1$ (zero when $j=0$). So each application of $(E-r)$ lowers the polynomial degree by one, and $(E-r)^{m}$ kills $n^jr^n$ for every $j<m$. The factors $(E-r_i)^{m_i}$ commute, so $p(E)$ kills all $k$ listed sequences and they lie in $V$.

Linear independence: suppose $\sum_iP_i(n)r_i^n=0$ for all $n$ with some $P_i\ne 0$. Order the roots with $|r_1|\ge|r_2|\ge\dots$; among the roots of largest modulus and largest $\deg P_i$, divide by $n^{d}r_1^{\,n}$ where $d$ is that degree. Terms with smaller modulus or smaller degree tend to $0$, and what remains is a finite sum of the form $\sum_j c_j\,e^{in\theta_j}$ with distinct $\theta_j$ and not all $c_j$ zero, which does not tend to $0$ (its average over $n$ of $|\cdot|^2$ is $\sum|c_j|^2>0$). Contradiction. Hence the $k$ sequences are independent and, by dimension, a basis.
\end{proof}

A check worth doing once by hand: $a_{n+2}=4a_{n+1}-4a_n$ has $p(x)=(x-2)^2$, so $a_n=(\alpha+\beta n)2^n$; with $a_0=1$, $a_1=4$ this is $(1+n)2^n$, and indeed $a_2=4\cdot4-4\cdot1=12=3\cdot 4$. Forgetting the factor $n$ for the repeated root is the commonest error in this section.

\begin{proposition}[Eventual periodicity modulo $m$]
Let $m\ge 2$ and let a sequence satisfy
\[
a_{n+k}=f(a_n,\dots,a_{n+k-1}),
\]
where $f$ maps $(\Z/m)^k$ to $\Z/m$. Then the sequence $(a_n \bmod m)$ is eventually periodic with period at most $m^k$. If moreover $f$ is a bijection in its first argument for every fixed value of the others (for a linear recurrence, when $c_0$ is a unit modulo $m$), the sequence is periodic from the start.
\end{proposition}

\begin{proof}
The state $s_n=(a_n,\dots,a_{n+k-1})\bmod m$ determines all later states, and there are only $m^k$ states, so two of $s_0,\dots,s_{m^k}$ coincide, say $s_i=s_j$ with $i<j$; then $s_{n+(j-i)}=s_n$ for all $n\ge i$. When $f$ is invertible in its first argument, $s_{n}$ is recovered from $s_{n+1}$ as well, so the map on states is a bijection of a finite set and every orbit is a cycle: $s_i=s_j$ implies $s_0=s_{j-i}$.
\end{proof}

\begin{proposition}[Linearisation by substitution]
Each of the following nonlinear recurrences is a linear one after a change of variable.
\begin{enumerate}[nosep]
\item $x_{n+1}=\dfrac{\alpha x_n+\beta}{\gamma x_n+\delta}$ with $M=\begin{pmatrix}\alpha&\beta\\ \gamma&\delta\end{pmatrix}$ invertible. If $(p_{n+1},q_{n+1})^{T}=M(p_n,q_n)^{T}$ and $x_0=p_0/q_0$, then $x_n=p_n/q_n$ whenever $q_n\ne 0$. So $x_n$ is a ratio of two sequences satisfying the linear recurrence with characteristic polynomial $\det(xI-M)$, and $M^k$ scalar forces period $k$.
\item $x_{n+1}=x_n^2-2$: with $x_n=t_n+t_n^{-1}$ one gets $t_{n+1}=t_n^2$, so $x_n=t_0^{2^n}+t_0^{-2^n}$; this needs $|x_0|\ge 2$ for real $t_0$, else $x_0=2\cos\theta$ and $x_n=2\cos(2^n\theta)$.
\item $x_{n+1}=\dfrac{x_n+c}{1-cx_n}$: with $x_n=\tan\theta_n$ and $c=\tan\varphi$ one gets $\theta_{n+1}=\theta_n+\varphi$.
\item $x_{n+1}=x_n^{\,r}\cdot K$: take logarithms, $\ln x_{n+1}=r\ln x_n+\ln K$, a first-order linear recurrence.
\end{enumerate}
\end{proposition}

\begin{proof}
(1) $\dfrac{p_{n+1}}{q_{n+1}}=\dfrac{\alpha p_n+\beta q_n}{\gamma p_n+\delta q_n}=\dfrac{\alpha x_n+\beta}{\gamma x_n+\delta}$ after dividing by $q_n$. By Cayley-Hamilton both $p_n$ and $q_n$ satisfy $u_{n+2}=\tr(M)u_{n+1}-\det(M)u_n$. (2) $(t+t^{-1})^2-2=t^2+t^{-2}$. (3) is the tangent addition formula. (4) is immediate.
\end{proof}

The other two pieces of the model come from neighbouring sections. A closed form guessed from the first terms is proved by induction, and the induction step is usually just the recurrence. A count $c_n$ that splits by the first object satisfies $c_n=\sum_j(\text{ways the first part is }j)\cdot c_{n-j}$, which is the convolution shape that generating functions handle.

\begin{firstmoves}
Compute six terms and look at them: Fibonacci, Lucas, Catalan, powers of $2$, factorials, triangular numbers and their small relatives are the usual disguises. Compute the first differences and the ratios of consecutive terms. Compute modulo $2$, $3$, $4$ and $5$ and look for the period. For a linear recurrence write the characteristic polynomial and factor it; a repeated root or a root of unity is a flag. For a rational or square-root recurrence try the substitutions above, in the order ratio, tangent, $t+1/t$, logarithm. For a recurrence with polynomial coefficients, write the generating function.
\end{firstmoves}

\begin{example}
Define $a_0=0$, $a_1=1$ and $a_{n+2}=6a_{n+1}-a_n$ for $n\ge 0$. Find a closed form for $a_n$, and show that $7$ divides $a_n$ if and only if $3$ divides $n$.
\end{example}

\begin{proof}[Solution]
The characteristic polynomial is $x^2-6x+1$ with roots $r=3+2\sqrt2$ and $\bar r=3-2\sqrt2$, distinct, so $a_n=Cr^n+D\bar r^{\,n}$. From $a_0=0$, $D=-C$; from $a_1=1$, $C(r-\bar r)=4\sqrt2\,C=1$. Hence
\[
a_n=\frac{(3+2\sqrt2)^n-(3-2\sqrt2)^n}{4\sqrt2}.
\]
The first terms are $0,1,6,35,204,1189,6930$, and $35=5\cdot7$, $6930=7\cdot 990$, while the others are not multiples of $7$. This suggests the claim.

Modulo $7$ the coefficient of $a_n$ in the recurrence is $-1$, a unit, so the residues are purely periodic by the proposition. Compute: $a_0\equiv0$, $a_1\equiv1$, $a_2\equiv6$, $a_3\equiv 36-1=35\equiv0$, $a_4\equiv 0-6\equiv1$, $a_5\equiv 6-0=6$, and the state $(a_3,a_4)\equiv(0,1)=(a_0,a_1)$ has returned. So the residues repeat with period $3$: $a_n\equiv 0,1,6$ according as $n\equiv 0,1,2\pmod 3$, and $7\mid a_n$ exactly when $3\mid n$.

Two remarks. The closed form was not needed for the divisibility, and could not have given it without more work; the finite state space did. And the proof of periodicity is the two lines of computation plus the sentence saying why a returning state closes the cycle, which is the sentence graders look for.
\end{proof}

\begin{example}
Let $x_0=2$ and $x_{n+1}=\dfrac{x_n+1}{1-x_n}$ for $n\ge0$. Show that $x_n$ is defined for every $n$, and find $x_{2026}$.
\end{example}

\begin{proof}[Solution]
The recurrence is the Möbius map with matrix $M=\begin{pmatrix}1&1\\-1&1\end{pmatrix}$. Compute $M^2=\begin{pmatrix}0&2\\-2&0\end{pmatrix}$ and $M^4=(M^2)^2=\begin{pmatrix}-4&0\\0&-4\end{pmatrix}=-4I$, a scalar. So if $(p_n,q_n)^T=M^n(2,1)^T$, then $(p_{n+4},q_{n+4})=-4(p_n,q_n)$ and $x_{n+4}=x_n$ whenever the quotients are defined.

Directly: $x_0=2$, $x_1=3/(-1)=-3$, $x_2=(-2)/4=-\tfrac12$, $x_3=(\tfrac12)/(\tfrac32)=\tfrac13$, $x_4=(\tfrac43)/(\tfrac23)=2=x_0$. None of these four values is $1$, so the denominator never vanishes, and the sequence is periodic with period $4$ from the start. Since $2026=4\cdot506+2$, $x_{2026}=x_2=-\tfrac12$.

The substitution view says the same thing: with $x_n=\tan\theta_n$ the map is $\theta_{n+1}=\theta_n+\pi/4$, and four steps add $\pi$, which does not change the tangent. Either view is a complete proof once the four values are listed; the matrix view also explains why the period is $4$ before any term is computed, since $M$ is $\sqrt2$ times a rotation by $45^\circ$.
\end{proof}

\begin{trap}
A repeated characteristic root, which needs the factor $n$; check by computing one more term than the number of constants. Initial terms that break the pattern: a recurrence valid for $n\ge 2$ with $a_0$, $a_1$ not fitting the closed form, so the closed form must be stated for $n\ge 2$ and the early terms listed separately. Assuming periodicity starts at the first term when the map on states is not a bijection; $a_{n+1}=a_n^2 \bmod 10$ from $a_0=2$ runs $2,4,6,6,6,\dots$ and is eventually constant, not periodic from the start. A substitution that is only valid on part of the domain, like $t+1/t$ for $|x|\ge 2$. Reporting a period found numerically without the sentence that a repeated state forces it. A limit of a recurrence computed by solving $L=f(L)$ before the sequence has been shown to converge.
\end{trap}

\begin{ladder}
\item \putnam{2016 B1}. A sequence defined through logarithms; sum a series built from it.
\item \seminar{PS 4 \#1}. A third-order linear recurrence; the terms as a sum of familiar sequences.
\item \seminar{PS 4 \#3}. The parity of the integer part of a power; a recurrence hidden in a surd.
\item \seminar{PS 4 \#5}. The ceiling of a power of a surd; divisibility by a power of two.
\item \seminar{PS 4 \#9}. A second-order recurrence; one specific far term.
\item \putnam{2017 A2}. A rational recurrence whose terms turn out to be integer polynomials.
\item \putnam{2025 A1}. An iteration on fractions; the terms eventually coprime.
\item \putnam{2020 A3}. An iterated sine; convergence of a series of the terms.
\item \putnam{2018 B4}. A cosine-like recurrence; one zero term forces periodicity.
\item \putnam{2021 B4}. A product of powers reduced modulo a Fibonacci number; the remainder.
\item \putnam{2019 B5}. A polynomial interpolating a famous recurrence; extrapolate one step.
\end{ladder}

\begin{hides}
Inside counting, where the count for $n$ splits by the first part (\putnam{2020 A5}, \putnam{2022 A3}, \putnam{2024 B5}). Inside probability, where conditioning on the first step turns a probability or an expectation into a recurrence (\putnam{2020 A4}, \putnam{2020 B3}, \putnam{2024 B4}). Inside number theory, where the exact power of a prime dividing a term follows the recurrence (\putnam{2025 A6}, \putnam{2021 B4}). Inside polynomials, where a sequence of polynomials is defined by iteration (\putnam{2017 A2}, \putnam{2019 B5}). Inside linear algebra, where the recurrence is a power of a matrix (\seminar{PS 7 \#24}, \seminar{PS 7 \#25}).
\end{hides}

\section{Sums, integrals and limits}

\begin{recognise}
A sum or an integral has to be evaluated in closed form, or bounded above and below. A limit of a sequence or a function is asked for, often of the form something to the power of something, or of a sum with $n$ terms each depending on $n$. A sequence has to be shown to converge. A product with many factors appears, which is a sum after a logarithm. A double sum or a double integral appears, which is usually easier in the other order. The seminar's second set is the reference; the notes before its problems state every theorem below.
\end{recognise}

\begin{lemma}[Telescoping]
For any sequence $f$, $\sum_{k=m}^{n}\bigl(f(k+1)-f(k)\bigr)=f(n+1)-f(m)$. The series $\sum_{k\ge m}(f(k+1)-f(k))$ converges if and only if $\lim f(n)$ exists, and then equals $\lim f(n)-f(m)$.
\end{lemma}

\begin{proof}
Write the terms out; every $f(k)$ with $m<k\le n$ appears once with each sign. The series statement is the finite one with $n\to\infty$.
\end{proof}

The work in a telescoping problem is finding $f$. Partial fractions give it for rational terms; the difference of a product with one more factor gives it for products; $\arctan$ addition gives it for $\arctan$ of a rational function; $\ln$ gives it for products.

\begin{theorem}[Tonelli and Fubini for sums]
Let $(a_{jk})_{j,k\ge 1}$ be a double sequence.
\begin{enumerate}[nosep]
\item (Tonelli) If $a_{jk}\ge 0$ for all $j,k$, then $\sum_j\sum_k a_{jk}=\sum_k\sum_j a_{jk}$, both sides being possibly $+\infty$.
\item (Fubini) If $\sum_j\sum_k|a_{jk}|<\infty$, then both iterated sums $\sum_j\sum_k a_{jk}$ and $\sum_k\sum_j a_{jk}$ converge and are equal, and so is the sum in any order of enumeration of the pairs.
\end{enumerate}
The same statements hold for integrals and for mixed sums and integrals, with ``$\ge 0$'' and ``absolutely integrable'' in place of the two hypotheses.
\end{theorem}

\begin{proof}[Proof sketch]
For Tonelli, let $S$ be the supremum of the finite sums $\sum_{j\le J}\sum_{k\le K}a_{jk}$ over all rectangles. Every partial sum of either iterated sum is at most $S$, and every rectangle sum is at most either iterated sum, so both iterated sums equal $S$. For Fubini, write $a_{jk}=a^+_{jk}-a^-_{jk}$ with $a^\pm\ge 0$ and $a^\pm\le|a|$; Tonelli applies to each part separately with finite value, and subtraction is allowed because both parts are finite.
\end{proof}

Why the hypothesis matters: take $a_{jj}=1$, $a_{j,j+1}=-1$, all other entries $0$. Every row sums to $0$, so $\sum_j\sum_k a_{jk}=0$; the first column sums to $1$ and every other column to $0$, so $\sum_k\sum_j a_{jk}=1$. The absolute sum is infinite, and the two orders disagree.

\begin{theorem}[Riemann sums]
Let $f$ be continuous on $[a,b]$. Then
\[
\lim_{n\to\infty}\frac{b-a}{n}\sum_{k=1}^{n}f\Bigl(a+k\,\frac{b-a}{n}\Bigr)=\int_a^b f(x)\,dx,
\]
and the same holds with $k$ running from $0$ to $n-1$, or with any choice of sample point in each subinterval.
\end{theorem}

\begin{proof}
Continuous on a closed interval means uniformly continuous: for $\varepsilon>0$ there is $\delta$ with $|f(x)-f(y)|<\varepsilon$ when $|x-y|<\delta$. For $(b-a)/n<\delta$, on each subinterval $I_k$ the sample value differs from $f(x)$ by less than $\varepsilon$ for all $x\in I_k$, so the sum differs from $\sum_k\int_{I_k}f=\int_a^b f$ by less than $\varepsilon(b-a)$.
\end{proof}

The recognition pattern: a sum of $n$ terms, a factor $1/n$ (sometimes hidden as $n/n^2$ or as $1/\sqrt{n^2}$), and each term a function of $k/n$ alone. A product $\prod_{k=1}^n g(k/n)^{1/n}$ is $\exp$ of a Riemann sum of $\ln g$.

\begin{lemma}[Integral comparison]
Let $f$ be non-negative and non-increasing on $[m,\infty)$, with $m$ an integer. Then for $n\ge m$,
\[
\int_m^{n+1}f(x)\,dx\ \le\ \sum_{k=m}^{n}f(k)\ \le\ f(m)+\int_m^{n}f(x)\,dx.
\]
In particular $\sum_{k\ge m}f(k)$ converges if and only if $\int_m^\infty f$ does.
\end{lemma}

\begin{proof}
On $[k,k+1]$, $f(k+1)\le f(x)\le f(k)$. Integrating gives $f(k+1)\le\int_k^{k+1}f\le f(k)$; sum the right inequality over $k=m,\dots,n$ for the lower bound and the left one over $k=m,\dots,n-1$ for the upper bound.
\end{proof}

With $f(x)=1/x$: $\ln(n+1)\le H_n\le 1+\ln n$. With $f(x)=x^{-1/2}$: $2\sqrt{n+1}-2\le\sum_{k=1}^n k^{-1/2}\le 2\sqrt n-1$. These two are worth knowing by heart, and the method gives the first-order growth of any monotone sum.

\begin{proposition}[First surviving term]
Let $u(x)=c\,x^{p}+o(x^{p})$ and $v(x)=d\,x^{q}+o(x^{q})$ as $x\to 0$, with $c,d\ne 0$. Then $u(x)/v(x)\to c/d$ if $p=q$, to $0$ if $p>q$, and is unbounded if $p<q$. If $u(x)\to 0$ and $w(x)\to\infty$ with $u(x)w(x)\to L$, then $(1+u(x))^{w(x)}\to e^{L}$.
\end{proposition}

\begin{proof}
Divide numerator and denominator by $x^q$. For the second statement, $\ln(1+u)=u+O(u^2)$ as $u\to 0$, so $w\ln(1+u)=wu+O(wu^2)=wu+O(u\cdot wu)\to L$, and $\exp$ is continuous.
\end{proof}

In practice: take the logarithm of anything of the form $A^B$, expand each factor with Taylor's theorem to one more order than seems necessary, cancel, and the first term that does not cancel gives the limit. If it is of order $x^0$ the limit is finite; the orders tell you how many terms were needed only after the cancellation is seen.

\begin{theorem}[Monotone convergence of sequences]
A non-decreasing sequence of reals that is bounded above converges to its supremum; a non-increasing one bounded below converges to its infimum. If $|x_{n+1}-L|\le q|x_n-L|$ for all $n$ with a fixed $q<1$, then $x_n\to L$.
\end{theorem}

\begin{proof}
Let $s=\sup x_n$. For $\varepsilon>0$ some $x_N>s-\varepsilon$, and then $s-\varepsilon<x_n\le s$ for all $n\ge N$. The contraction statement is $|x_n-L|\le q^n|x_0-L|\to 0$.
\end{proof}

For a recurrence $x_{n+1}=g(x_n)$ with continuous $g$, the route is: show the sequence is monotone and bounded (often by induction, comparing $x_n$ with the fixed point), conclude it converges to some $L$, and only then pass to the limit in the recurrence to get $L=g(L)$.

\begin{firstmoves}
Write the first three terms of a sum out in full. Look for a difference $f(k+1)-f(k)$; try partial fractions. Put $1/n$ in front and ask whether what remains is a function of $k/n$. Take a logarithm of a product or a power. Compare with an integral for growth. For a limit of a sequence given by a recurrence, find the fixed point, then prove monotone and bounded. For a double sum, write the region of summation and swap, after checking non-negativity or absolute convergence. For a limit of an integral depending on a parameter, substitute to make the interval fixed and then look for a dominating function.
\end{firstmoves}

\begin{example}
Evaluate $\displaystyle\lim_{n\to\infty}\sum_{k=1}^{n}\frac{n}{n^2+nk+k^2}$.
\end{example}

\begin{proof}[Solution]
Divide numerator and denominator by $n^2$:
\[
\sum_{k=1}^n\frac{n}{n^2+nk+k^2}=\frac1n\sum_{k=1}^n\frac{1}{1+(k/n)+(k/n)^2},
\]
a Riemann sum for $f(x)=1/(1+x+x^2)$ on $[0,1]$ with right endpoints. $f$ is continuous, so the limit is
\[
\int_0^1\frac{dx}{1+x+x^2}=\int_0^1\frac{dx}{(x+\frac12)^2+\frac34}=\frac{2}{\sqrt3}\Bigl[\arctan\frac{2x+1}{\sqrt3}\Bigr]_0^1=\frac{2}{\sqrt3}\Bigl(\frac\pi3-\frac\pi6\Bigr)=\frac{\pi}{3\sqrt3}.
\]
Sanity check: $f$ lies between $1/3$ and $1$ on $[0,1]$, so the answer must lie between $1/3$ and $1$; $\pi/(3\sqrt3)\approx 0.605$. At $n=1$ the sum is $1/3$, at $n=2$ it is $2/7+2/12\approx 0.452$, rising towards the limit as the terms near $k/n\approx 0$ gain weight.
\end{proof}

\begin{example}
Evaluate $\displaystyle\lim_{n\to\infty}n^2\Bigl(\bigl(1+\tfrac1n\bigr)^{n+\frac12}-e\Bigr)$.
\end{example}

\begin{proof}[Solution]
Let $L_n=(n+\tfrac12)\ln(1+\tfrac1n)$, so the expression is $n^2(e^{L_n}-e)=n^2e\,(e^{L_n-1}-1)$. Expand with $h=1/n$:
\begin{align*}
\ln(1+h)&=h-\frac{h^2}{2}+\frac{h^3}{3}-O(h^4),\\
L_n=\Bigl(\frac1h+\frac12\Bigr)\ln(1+h)&=1-\frac h2+\frac{h^2}{3}+\frac h2-\frac{h^2}{4}+O(h^3)=1+\frac{h^2}{12}+O(h^3).
\end{align*}
The $h$ terms cancel, which is why the exponent $n+\tfrac12$ was chosen, and the first surviving term is $h^2/12$. Then $e^{L_n-1}-1=\frac{h^2}{12}+O(h^3)$, and
\[
n^2\Bigl(\bigl(1+\tfrac1n\bigr)^{n+\frac12}-e\Bigr)=e\cdot\frac{1}{h^2}\Bigl(\frac{h^2}{12}+O(h^3)\Bigr)\longrightarrow\frac{e}{12}.
\]
Had the exponent been $n$ alone, $L_n=1-h/2+O(h^2)$, the $h$ term would survive, and $n\bigl((1+\frac1n)^n-e\bigr)\to-e/2$ while the $n^2$ version diverges. The number of Taylor terms needed is not visible until the cancellation is done, so expand one order further than the question seems to want.
\end{proof}

\begin{trap}
Swapping the order of summation when the double series is only conditionally convergent; the diagonal example above is the picture to keep. A Riemann sum whose endpoints are off by one term that matters: $\frac1n\sum_{k=1}^n f(k/n)$ and $\frac1n\sum_{k=0}^{n-1}f(k/n)$ have the same limit for continuous $f$, but if $f$ blows up at an endpoint the single term can carry the whole answer. A limit taken inside an integral or an infinite sum without a reason; the reasons are monotone convergence (the integrands increase to the limit) and dominated convergence (a fixed integrable function bounds them all), and one of the two must be named. Keeping too few terms of an expansion, so the first surviving term is actually an error term. Solving $L=g(L)$ for the limit of a recurrence that does not converge; $x_{n+1}=2x_n$ has fixed point $0$. Comparing with an integral when the function is not monotone.
\end{trap}

\begin{ladder}
\item \seminar{Limits \#15}. A sum of $n$ rational terms; recognise an integral.
\item \seminar{Limits \#12}. A product of $n$ factors with a normalising power; logarithms.
\item \putnam{2016 B1}. A sequence defined through logarithms; sum a series from it.
\item \putnam{2022 A1}. A line and a logarithmic curve meeting exactly once; the parameters.
\item \putnam{2019 B2}. A trigonometric sum; its asymptotic limit.
\item \putnam{2021 A2}. A function defined by a limit; an asymptotic ratio.
\item \putnam{2020 A2}. A sum of binomial coefficients; a closed form.
\item \putnam{2024 B3}. Gaps between the fixed points of a tangent-type function; a two-sided bound.
\item \putnam{2023 A1}. A product of cosines; its second derivative at a point exceeds a threshold.
\item \putnam{2017 B4}. An alternating series of logarithms; evaluate.
\item \putnam{2016 B6}. A double alternating series; evaluate.
\end{ladder}

\begin{hides}
Inside inequalities, as the bound that squeezes a sequence or a function (\putnam{2017 A3}, \putnam{2023 A3}, \putnam{2025 B2}). Inside probability, where an expectation is a sum or an integral and a limiting probability is a limit of a recurrence (\putnam{2022 A4}, \putnam{2020 A4}). Inside recurrences, where the asymptotic growth of a sequence is found by comparing with an integral (\seminar{PS 4 \#18}, \putnam{2020 A3}). Inside functional equations, where the answer is an integral of the unknown function (\putnam{2016 A3}). Inside polynomials, where a product of many factors is bounded through its logarithm (\putnam{2023 A1}).
\end{hides}

\section{Probability and expectation}

\begin{recognise}
The words random, expected, probability, on average, or almost surely appear. A counting problem asks for a proportion, which is a probability after dividing by the total. A process runs in steps and a stopping time or a final state is asked about. Points are chosen in a region and a geometric condition is imposed. Behind most Putnam probability there is no measure theory at all: a finite sample space, linearity of expectation, a recurrence from the first step, or an area.
\end{recognise}

\begin{definition}
On a finite or countable sample space $\Omega$ with probabilities $p(\omega)\ge 0$ summing to $1$, a random variable is a function $X\colon\Omega\to\R$, its expectation is $\E X=\sum_\omega X(\omega)p(\omega)$ when the sum converges absolutely, and the indicator of an event $A\subseteq\Omega$ is $\mathbf 1_A$, with $\E\,\mathbf 1_A=\Pr(A)$. The conditional probability of $A$ given $B$ with $\Pr(B)>0$ is $\Pr(A\mid B)=\Pr(A\cap B)/\Pr(B)$, and $\E(X\mid B)$ is the expectation under the probabilities $p(\omega)/\Pr(B)$ on $B$.
\end{definition}

\begin{theorem}[Linearity and indicators]
$\E(aX+bY)=a\E X+b\E Y$ for any random variables with finite expectation, independent or not. If $N=\sum_{i=1}^n\mathbf 1_{A_i}$ counts how many of the events $A_i$ occur, then
\[
\E N=\sum_i\Pr(A_i),\qquad \E N^2=\sum_i\Pr(A_i)+2\sum_{i<j}\Pr(A_i\cap A_j).
\]
\end{theorem}

\begin{proof}
The first is the distributive law inside the defining sum. For the second, $N^2=\sum_i\mathbf 1_{A_i}^2+2\sum_{i<j}\mathbf 1_{A_i}\mathbf 1_{A_j}$, and $\mathbf 1_{A}^2=\mathbf 1_A$, $\mathbf 1_A\mathbf 1_B=\mathbf 1_{A\cap B}$.
\end{proof}

The second formula is how a variance is computed in a problem that asks whether a count is concentrated, and it is the only place where the pairwise interactions enter.

\begin{theorem}[Total probability and conditioning on the first step]
Let $B_1,B_2,\dots$ be a partition of $\Omega$ into events of positive probability. Then for any event $A$ and any random variable $X$ with finite expectation,
\[
\Pr(A)=\sum_i\Pr(A\mid B_i)\Pr(B_i),\qquad \E X=\sum_i\E(X\mid B_i)\Pr(B_i).
\]
\end{theorem}

\begin{proof}
$A$ is the disjoint union of the $A\cap B_i$, and $\Pr(A\cap B_i)=\Pr(A\mid B_i)\Pr(B_i)$. The expectation version is the same decomposition of the defining sum.
\end{proof}

For a process that moves between finitely many states, condition on the first step: if $h(s)$ is the probability of eventually reaching a target from state $s$, or the expected number of steps to reach it, then $h(s)=\sum_{s'}P(s\to s')h(s')$, plus $1$ in the expected-time case, with $h$ known on the target. This is a linear system, and a Putnam problem is built so that it has a short solution: a symmetry that makes most unknowns equal, or a telescoping along the states.

\begin{proposition}[Tail sum]
For a random variable $N$ taking values in $\{0,1,2,\dots\}$, $\E N=\sum_{k\ge 1}\Pr(N\ge k)$. For $X\ge 0$ with a density, $\E X=\int_0^\infty\Pr(X>t)\,dt$.
\end{proposition}

\begin{proof}
$N=\sum_{k\ge1}\mathbf 1[N\ge k]$ pointwise; take expectations and swap the sum and the expectation, which Tonelli allows since everything is non-negative. The continuous case is the same with an integral.
\end{proof}

\begin{proposition}[Geometric probability]
If a point is chosen uniformly in a region $R\subset\R^d$ of finite positive volume, the probability that it lies in a measurable $S\subseteq R$ is $\operatorname{vol}(S)/\operatorname{vol}(R)$. If $X_1,\dots,X_d$ are independent and each uniform on $[0,1]$, then $(X_1,\dots,X_d)$ is uniform on the cube $[0,1]^d$, so a probability about them is a volume.
\end{proposition}

\begin{proof}
The first is the definition of the uniform distribution. For the second, independence means the joint density is the product of the marginal densities, which is $1$ on the cube.
\end{proof}

\begin{proposition}[The probabilistic method]
If $\E X=\mu$ then $\Pr(X\ge\mu)>0$ and $\Pr(X\le\mu)>0$. In particular, if a random object has expected score $\mu$, some object has score at least $\mu$. If $X\ge 0$ then $\Pr(X\ge a)\le\E X/a$ for $a>0$ (Markov), so if $\E X<1$ for a non-negative integer $X$ then $\Pr(X=0)>0$.
\end{proposition}

\begin{proof}
If $X<\mu$ everywhere then $\E X<\mu$. For Markov, $a\,\mathbf 1[X\ge a]\le X$ pointwise; take expectations.
\end{proof}

\begin{firstmoves}
Write the quantity as a sum of indicators and compute one probability. Condition on the first step and write the recurrence. List every outcome of the smallest case and count. For points in a region, draw the region where the condition holds and compute its area; look for a line or a curve as the boundary. For an expected stopping time, ask whether every step has the same chance of stopping, which makes it geometric. For a question of the form ``show that some arrangement has property $P$'', choose the arrangement at random and compute the expected number of failures.
\end{firstmoves}

\begin{example}
A fair six-sided die is rolled repeatedly and the running total of the rolls is kept. Find the expected number of rolls until the running total is first divisible by $4$.
\end{example}

\begin{proof}[Solution]
Only the residue of the total modulo $4$ matters. A roll adds a residue $d\in\{0,1,2,3\}$ with probabilities $p_0=\tfrac16$, $p_1=\tfrac26$, $p_2=\tfrac26$, $p_3=\tfrac16$, since the faces $1,\dots,6$ reduce to $1,2,3,0,1,2$. Let $T_r$ be the expected number of further rolls needed to reach residue $0$ from residue $r$, with $T_0=0$, and let $P(r,s)=p_{s-r}$ (indices modulo $4$). A step of $+1$ has positive probability, so every residue is reachable from every other and all $T_r$ are finite. Conditioning on the next roll, for $r\ne 0$,
\[
T_r=1+\sum_{s}P(r,s)\,T_s .
\]
The answer asked for is the expected return time to $0$ starting from a fresh total of $0$: $T^{*}=1+\sum_s P(0,s)T_s$.

Rather than solving the system, add the three equations for $r=1,2,3$:
\[
\sum_{r\ne 0}T_r=3+\sum_s T_s\sum_{r\ne 0}P(r,s)=3+\sum_s T_s\bigl(1-P(0,s)\bigr),
\]
because each column of $P$ sums to $1$: $\sum_rP(r,s)=\sum_r p_{s-r}=1$. Since $T_0=0$, the left side is $\sum_sT_s$, and cancelling it gives $\sum_sP(0,s)T_s=3$. Hence $T^{*}=1+3=4$.

As a check, the system solves to $T_1=\tfrac{98}{25}$, $T_2=\tfrac{84}{25}$, $T_3=\tfrac{86}{25}$, and $1+\tfrac16\bigl(2\cdot\tfrac{98}{25}+2\cdot\tfrac{84}{25}+\tfrac{86}{25}\bigr)=1+\tfrac{450}{150}=4$. The argument used nothing about the die except that each column of $P$ sums to $1$, which holds for any step distribution on any modulus $m$ with the walk irreducible, and gives $m$ every time. That is the uniformity of the previous section hiding in a probability problem: the stationary distribution of a random walk on a group is uniform, so the mean return time to any state is the number of states.
\end{proof}

\begin{example}
Two numbers $x$ and $y$ are chosen independently and uniformly from $[0,1]$. Find the probability that $x+y>1$ and $xy<\tfrac14$.
\end{example}

\begin{proof}[Solution]
The pair $(x,y)$ is uniform on the unit square, so the probability is the area of the region $\{x+y>1,\ xy<\tfrac14\}$. The region $\{x+y>1\}$ is the triangle above the diagonal, of area $\tfrac12$. Inside it, remove the part where $xy\ge\tfrac14$.

By the arithmetic-geometric mean inequality, $xy\ge\tfrac14$ forces $x+y\ge 2\sqrt{xy}\ge 1$, with equality only at the single point $(\tfrac12,\tfrac12)$. So $\{xy\ge\tfrac14\}$ lies inside $\{x+y\ge 1\}$ up to a set of area zero, and the area to remove is the whole of $\{xy\ge\tfrac14\}$ in the square. That region is $\tfrac14\le x\le 1$, $\tfrac1{4x}\le y\le 1$, with area
\[
\int_{1/4}^{1}\Bigl(1-\frac{1}{4x}\Bigr)dx=\frac34-\frac14\ln 4=\frac34-\frac{\ln 2}{2}.
\]
Therefore the probability is
\[
\frac12-\Bigl(\frac34-\frac{\ln2}{2}\Bigr)=\frac{\ln 2}{2}-\frac14=\frac{2\ln2-1}{4}\approx 0.097.
\]
It is positive, as it should be since $(0.9,0.2)$ is in the region, and small, since the hyperbola $xy=\tfrac14$ hugs the line $x+y=1$ near the centre. The inequality step is what made the picture simple: without it, the region to remove would have had to be intersected with the triangle by hand.
\end{proof}

\begin{trap}
Multiplying probabilities of events that share randomness; two positions of a random permutation being fixed are not independent, and the indicator method is exactly the way to avoid needing them to be. A sample space that is not uniform: two dice summing to $7$ is six outcomes out of $36$, not one of eleven. Conditioning on the first step and forgetting the $+1$ for the step itself. A stopping time whose expectation is infinite, so the recurrence has no finite solution; check irreducibility or a geometric bound first. In geometric probability, a region drawn with a straight boundary where the condition gives a curve. Computing an expectation and reporting it as the most likely value; the expected number of fixed points of a permutation is $1$, which says nothing about what happens most often.
\end{trap}

\begin{ladder}
\item \seminar{PS 6 \#3}. A biased coin tossed many times; the parity of the number of heads.
\item \seminar{PS 6 \#5}. Uniform draws added until a threshold is crossed; the expected count.
\item \seminar{PS 6 \#13}. A random point in a disc; nearer the centre or the edge.
\item \putnam{2021 B1}. A random square placed on a grid; the chance it avoids certain points.
\item \putnam{2020 A4}. Random blackening of a strip; the limiting white fraction.
\item \putnam{2020 B3}. A randomly shrinking interval; the expected number of steps below a threshold.
\item \putnam{2023 B3}. A random sequence and its longest alternating subsequence; the expectation.
\item \putnam{2024 B4}. A walk pulled toward uniform draws; the limiting mean.
\item \putnam{2022 A4}. A weighted sum of random terms until an ascent; the expectation.
\item \putnam{2016 B4}. A random antisymmetric matrix; the expected determinant.
\end{ladder}

\begin{hides}
Inside counting, by the probabilistic method: a random object has the property with positive probability, so one exists (\putnam{2024 A3}, \putnam{2018 B6}). Inside recurrences, where the first-step equation is the recurrence (\putnam{2020 A4}, \putnam{2020 B3}). Inside geometry, where a chance is a ratio of areas (\putnam{2021 B1}, \putnam{2024 A5}). Inside inequalities, where a bound on an expectation or a tail is the claim (\putnam{2021 B6}, \putnam{2022 B5}). Inside linear algebra, where the expected determinant of a random matrix is a signed sum of expected products (\putnam{2016 B4}, \seminar{PS 7 \#12}). Inside games, where the two players' chances are compared (\putnam{2017 A5}, \seminar{PS 6 \#24} to \seminar{\#28}).
\end{hides}

\section{Linear algebra}

\begin{recognise}
The problem names matrices, determinants, ranks, eigenvalues or vector spaces. Or it does not, and the problem is about polynomials of bounded degree, functions on a finite set, sequences satisfying a linear recurrence, lamps and switches, or a set of vectors with a condition on their inner products: these are vector spaces and linear maps in other clothes. A question that asks how many of something can exist is often a dimension bound. A question that says ``for all $n$'' about an $n\times n$ matrix is settled by the $2\times2$ and $3\times3$ cases and the right general fact.
\end{recognise}

\begin{theorem}[Rank-nullity]
Let $T\colon V\to W$ be linear with $V$ finite-dimensional. Then $\dim\ker T+\dim\operatorname{im}T=\dim V$. Consequently: more than $\dim V$ vectors in $V$ are linearly dependent; a linear map between spaces of the same finite dimension is injective if and only if it is surjective; and a square matrix is invertible if and only if its kernel is zero.
\end{theorem}

\begin{proof}
Choose a basis $u_1,\dots,u_k$ of $\ker T$ and extend it to a basis $u_1,\dots,u_k,v_1,\dots,v_r$ of $V$. The vectors $Tv_1,\dots,Tv_r$ span $\operatorname{im}T$ because $T$ kills the $u_i$, and they are independent because $\sum c_iTv_i=0$ puts $\sum c_iv_i$ in the kernel, where it is a combination of the $u_i$, so all $c_i=0$ by independence of the full basis. Hence $\dim\operatorname{im}T=r=\dim V-k$. The consequences are the equation read with one term zero or with $\dim V=\dim W$.
\end{proof}

The hidden vector spaces to know: polynomials of degree at most $n$ have dimension $n+1$, so $n+2$ of them are dependent and a polynomial of degree at most $n$ vanishing at $n+1$ points is zero; real functions on a set of size $m$ have dimension $m$; solutions of a linear recurrence of order $k$ have dimension $k$; $n\times n$ matrices have dimension $n^2$, so the powers $I,A,\dots,A^{n^2}$ are dependent and every matrix satisfies a polynomial.

\begin{theorem}[Eigenvalues, trace, determinant]
Let $A$ be an $n\times n$ matrix over $\C$ with characteristic polynomial $\chi_A(x)=\det(xI-A)=\prod_{i=1}^n(x-\lambda_i)$, the $\lambda_i$ listed with multiplicity. Then $\tr A=\sum\lambda_i$ and $\det A=\prod\lambda_i$. Moreover $\chi_A(A)=0$ (Cayley-Hamilton), and for any polynomial $f$ the eigenvalues of $f(A)$ are $f(\lambda_i)$ with multiplicity.
\end{theorem}

\begin{proof}
Expanding $\det(xI-A)$, the only term containing $x^{n}$ or $x^{n-1}$ is the product of the diagonal entries $\prod(x-a_{ii})$, since any other permutation misses at least two diagonal positions. So the coefficient of $x^{n-1}$ is $-\tr A$, which in the factored form is $-\sum\lambda_i$. Setting $x=0$ gives $\det(-A)=(-1)^n\det A$ on one side and $(-1)^n\prod\lambda_i$ on the other. For the eigenvalues of $f(A)$: over $\C$, $A$ is similar to an upper triangular matrix $U$ with the $\lambda_i$ on the diagonal, and $f(U)$ is upper triangular with $f(\lambda_i)$ on the diagonal. Cayley-Hamilton follows from the same triangularisation: $\prod(U-\lambda_iI)$ kills the standard basis vectors one at a time, in order.
\end{proof}

\begin{theorem}[Real symmetric matrices]
Let $A$ be a real symmetric $n\times n$ matrix. Then every eigenvalue of $A$ is real, eigenvectors for distinct eigenvalues are orthogonal, and $\R^n$ has an orthonormal basis of eigenvectors of $A$. In particular $A=Q\,\mathrm{diag}(\lambda_1,\dots,\lambda_n)\,Q^{T}$ with $Q$ orthogonal, and $x^TAx=\sum\lambda_iy_i^2$ in the coordinates $y=Q^Tx$, so $A$ is positive definite if and only if all $\lambda_i>0$.
\end{theorem}

\begin{proof}
Let $Av=\lambda v$ with $v\ne 0$ complex. Then $\bar v^{\,T}Av=\lambda|v|^2$, and the left side is real because its conjugate transpose is $\bar v^{\,T}A^{T}v=\bar v^{\,T}Av$ itself. So $\lambda\in\R$, and a real eigenvector exists since $A-\lambda I$ is a real singular matrix. If $Av=\lambda v$ and $Aw=\mu w$ with $\lambda\ne\mu$, then $\lambda\langle v,w\rangle=\langle Av,w\rangle=\langle v,Aw\rangle=\mu\langle v,w\rangle$, so $\langle v,w\rangle=0$. For the basis, take a unit eigenvector $v_1$; $A$ maps the orthogonal complement $v_1^\perp$ to itself (if $\langle w,v_1\rangle=0$ then $\langle Aw,v_1\rangle=\langle w,Av_1\rangle=0$) and is symmetric there, so induct on the dimension.
\end{proof}

\begin{theorem}[Determinants as alternating multilinear forms]
There is exactly one function $\det$ from $n\times n$ matrices over a field to the field that is linear in each row, changes sign when two rows are swapped, and takes the value $1$ at $I$. It satisfies: $\det$ is zero when two rows are equal or one row is zero; adding a multiple of one row to another leaves it unchanged; $\det(AB)=\det A\det B$; $\det A^{T}=\det A$; and $\det A\ne 0$ if and only if the rows are linearly independent.
\end{theorem}

\begin{proof}[Proof sketch]
Expanding each row in the standard basis by multilinearity writes $\det A$ as
\[
\sum_{\sigma}\ \prod_ia_{i\sigma(i)}\,\det(e_{\sigma(1)},\dots,e_{\sigma(n)}),
\]
the sum over all maps $\sigma$ of $\{1,\dots,n\}$ to itself, and the alternating property forces $\det(e_{\sigma(1)},\dots,e_{\sigma(n)})=\operatorname{sgn}(\sigma)$ for permutations and $0$ otherwise. That gives uniqueness and the formula; checking the formula has the three properties gives existence. Two equal rows: swapping them negates $\det$ and changes nothing, so $\det=-\det$ (in characteristic $2$ use the expansion directly). Row operations: linearity in the changed row plus the equal-rows case. For $\det(AB)$: the map $X\mapsto\det(XB)$ is alternating multilinear in the rows of $X$, so it equals $c\det X$ with $c=\det B$. Transpose: the permutation formula is symmetric under $\sigma\mapsto\sigma^{-1}$. Independence: dependent rows reduce by row operations to a zero row; independent rows reduce to $I$ by operations that only scale $\det$ by nonzero factors.
\end{proof}

Row operations are the working tool. Subtract each row from the next before expanding; add all columns to the first when every row has the same sum; factor a common factor out of a row. A determinant with a parameter $x$ is a polynomial in $x$, and its roots and degree often identify it without expansion.

\begin{proposition}[Linear algebra over the field with two elements]
Let $\mathbb F_2=\{0,1\}$ with addition modulo $2$. A configuration of $n$ lamps is a vector in $\mathbb F_2^n$, a switch that toggles a fixed set of lamps is a vector $s$, and pressing switches is adding vectors, so the reachable configurations from the all-off state are the span of the switch vectors. For a matrix $A$ over $\mathbb F_2$ (or any field), $Ax=b$ has a solution if and only if $b\cdot y=0$ for every $y$ with $A^{T}y=0$; and if a solution exists, the number of solutions is $2^{\dim\ker A}$.
\end{proposition}

\begin{proof}
Pressing a switch twice does nothing and the order does not matter, so the reachable set is the set of sums of subsets of switch vectors, which is the span. For the solvability criterion: if $Ax=b$ and $A^Ty=0$ then $b\cdot y=x\cdot A^Ty=0$, so $\operatorname{im}A\subseteq(\ker A^T)^\perp$; both have dimension $\rank A=n-\dim\ker A^T$ (rank-nullity for $A^T$, and $\rank A=\rank A^T$), so they are equal. Solutions form a coset of $\ker A$.
\end{proof}

\begin{firstmoves}
Compute the $2\times2$ and $3\times3$ cases by hand. Look at the trace and the determinant, and at what a polynomial identity in the matrix says about its eigenvalues. Ask what the kernel is, and what the dimension count gives. Name the hidden vector space and write the hidden linear map. For a determinant, do row operations before anything else, and treat a parameter as a polynomial variable. For symmetric matrices, diagonalise. For anything with parity, work over $\mathbb F_2$. For a set of vectors with a condition on their inner products, form the Gram matrix.
\end{firstmoves}

\begin{example}
Let $A$ be an $n\times n$ real matrix with $A^3=A$. Prove that $\rank A+\rank(I-A^2)=n$.
\end{example}

\begin{proof}[Solution]
Two inclusions give two inequalities. First, $A(I-A^2)=A-A^3=0$, so every vector in $\operatorname{im}(I-A^2)$ lies in $\ker A$, and $\rank(I-A^2)\le\dim\ker A=n-\rank A$ by rank-nullity. Second, every $v$ satisfies $v=A^2v+(I-A^2)v$, so $\R^n=\operatorname{im}A^2+\operatorname{im}(I-A^2)$ and
\[
n\le\rank A^2+\rank(I-A^2)\le\rank A+\rank(I-A^2),
\]
using $\operatorname{im}A^2\subseteq\operatorname{im}A$. Together, $\rank A+\rank(I-A^2)=n$.

The eigenvalue view explains the answer. $A$ satisfies $x^3-x=x(x-1)(x+1)$, a polynomial with distinct roots, so $A$ is diagonalisable with eigenvalues in $\{0,1,-1\}$; $\rank A$ is the number of nonzero eigenvalues and $\rank(I-A^2)$ the number of zero ones. But the rank-nullity argument above needs no diagonalisation and works over any field, including $\mathbb F_2$, where $x(x-1)(x+1)=x(x-1)^2$ has a repeated root and $A$ need not be diagonalisable. Writing the inclusion argument is also faster under exam conditions than justifying diagonalisability.
\end{proof}

\begin{example}
A row of $n\ge 1$ lamps is controlled by $n$ switches; switch $i$ toggles lamp $i$ and each of its immediate neighbours (lamps $i-1$ and $i+1$ when they exist). Starting with all lamps off, for which $n$ can every one of the $2^n$ configurations be reached?
\end{example}

\begin{proof}[Solution]
Over $\mathbb F_2$, switch $i$ is the vector with ones in positions $i-1,i,i+1$, and the reachable configurations are the span of the $n$ switch vectors. Every configuration is reachable if and only if the span is all of $\mathbb F_2^n$, that is, if and only if the $n\times n$ matrix $M_n$ with ones on the diagonal and both adjacent diagonals is invertible over $\mathbb F_2$, that is, $\det M_n\ne 0$ in $\mathbb F_2$.

Let $D_n=\det M_n\in\mathbb F_2$. Expanding along the first row, the entry $1$ at $(1,1)$ gives $D_{n-1}$, and the entry $1$ at $(1,2)$ gives the minor whose first column is $(1,0,\dots,0)^T$, which expands to $D_{n-2}$; signs are irrelevant modulo $2$. So $D_n=D_{n-1}+D_{n-2}$, with $D_1=1$ and $D_2=1\cdot1-1\cdot1=0$. The sequence is the Fibonacci numbers modulo $2$: $1,0,1,1,0,1,1,0,\dots$, periodic with period $3$ from $n=1$ (period from the start since the recurrence runs backwards too). Hence $D_n=0$ exactly when $n\equiv 2\pmod 3$, and every configuration is reachable if and only if $n\not\equiv 2\pmod 3$.

Check: $n=2$, both switches toggle both lamps, so only the all-off and all-on states are reachable; $n=3$, switch $2$ toggles everything, switch $1$ toggles lamps $1,2$, so lamp $3$ alone is reachable, and by symmetry lamp $1$, and then everything. When $n\equiv 2\pmod 3$ the kernel is one-dimensional (the vector with ones in positions $\equiv 1,2\pmod 3$ and zeros in positions $\equiv 0$), so exactly half the configurations are reachable, each in two ways.
\end{proof}

\begin{trap}
A matrix that is not diagonalisable: $\begin{pmatrix}0&1\\0&0\end{pmatrix}$ has all eigenvalues zero and is not zero, so an argument that works ``on each eigenvector'' needs a basis of eigenvectors, which a polynomial identity with distinct roots or symmetry provides and nothing else does for free. A determinant computed by cofactor expansion when row operations would finish in two lines. A dimension argument that asserts independence without proving it; the count of vectors is not the dimension of their span. Eigenvalues of a real matrix assumed real; a rotation has none. Over $\mathbb F_2$, $1=-1$, so symmetric and skew-symmetric coincide and a determinant that is a square over $\R$ is just a number. The trace of $AB$ equals the trace of $BA$, but $AB$ and $BA$ need not have the same rank or be similar when the matrices are not square.
\end{trap}

\begin{ladder}
\item \seminar{PS 7 \#1}. A matrix of signs with a circulant pattern; its rank.
\item \seminar{PS 7 \#7}. A matrix with a fixed off-diagonal entry and a few diagonal values; its determinant.
\item \seminar{PS 7 \#9}. A skew-symmetric matrix of even order; its determinant is a perfect square.
\item \putnam{2018 A2}. A matrix indexed by subsets, entries by intersection; its determinant.
\item \putnam{2019 B3}. An orthogonal matrix times a reflection; an eigenvalue exists.
\item \putnam{2022 B2}. A finite set of vectors closed under the cross product; its possible sizes.
\item \putnam{2017 B1}. Two distinct lines in the plane; a point of one scaled onto the other.
\item \putnam{2021 B5}. Integer matrices with odd principal minors; closure under powers.
\item \putnam{2025 A4}. A family of matrices commuting around a cycle; the least possible size.
\end{ladder}

\begin{hides}
Inside counting, where a determinant of an incidence or intersection matrix counts something and its nonvanishing bounds a family (\putnam{2018 A2}, \putnam{2019 B6}, \putnam{2025 A4}). Inside geometry, where lines and points are vectors (\putnam{2017 B1}, \putnam{2018 A6}). Inside probability, where an expected determinant is computed by linearity over the permutation expansion (\putnam{2016 B4}). Inside polynomials, where a space of polynomials has a dimension and a linear condition cuts it down (\putnam{2022 A6}). Inside generating functions, where a Hankel matrix of coefficients has a determinant (\putnam{2024 A6}). Inside number theory, where integer matrices are reduced modulo a prime and the rank there controls divisibility (\putnam{2021 B5}, \putnam{2023 B6}, \putnam{2023 A4}).
\end{hides}

\section{Counting and configurations}

\begin{recognise}
Arrangements, colourings, tilings, subsets with a property, lattice points, sequences with a forbidden pattern. People and acquaintances, teams and matches, committees and members, which are graphs and incidence structures in plain clothes. The question is how many, or whether a configuration exists, or whether it is impossible, or the largest size a configuration can have. The last two usually need two arguments: a proof that nothing bigger works and a construction that achieves the bound.
\end{recognise}

\begin{theorem}[Double counting]
Let $A$ and $B$ be finite sets and $R\subseteq A\times B$. For $a\in A$ let $d(a)$ be the number of $b$ with $(a,b)\in R$, and for $b\in B$ let $d(b)$ be the number of $a$ with $(a,b)\in R$. Then $\sum_{a\in A}d(a)=|R|=\sum_{b\in B}d(b)$.
\end{theorem}

\begin{proof}
Both sides count the pairs in $R$, once grouped by the first coordinate and once by the second.
\end{proof}

Consequences used constantly: in a finite graph the degrees sum to twice the number of edges, so the number of odd-degree vertices is even; in a family of $k$-element subsets of an $n$-set in which every element lies in exactly $r$ sets, $kF=rn$ where $F$ is the number of sets; and $\sum_k k\binom nk=n2^{n-1}$, by counting pairs (subset, marked element) both ways.

\begin{theorem}[Bijections and inclusion-exclusion]
If $f\colon X\to Y$ is a bijection of finite sets then $|X|=|Y|$. For finite sets $A_1,\dots,A_n$ inside a finite $U$,
\[
\Bigl|\bigcup_iA_i\Bigr|=\sum_{\emptyset\ne S\subseteq\{1,\dots,n\}}(-1)^{|S|+1}\Bigl|\bigcap_{i\in S}A_i\Bigr|,
\]
and the number of elements of $U$ in none of the $A_i$ is $|U|$ minus this, which for $n=1$ is complementary counting.
\end{theorem}

\begin{proof}
For a bijection, pair each element with its image. For inclusion-exclusion, fix $u\in U$ and let $m$ be the number of $A_i$ containing it. Pointwise, $1-\prod_i(1-\mathbf 1_{A_i}(u))$ equals $1$ if $m\ge 1$ and $0$ if $m=0$; expanding the product gives $\sum_{S\ne\emptyset}(-1)^{|S|+1}\prod_{i\in S}\mathbf 1_{A_i}(u)$. Sum over $u$.
\end{proof}

\begin{theorem}[Averaging and the probabilistic method]
If a finite set of objects has total score $T$ and there are $N$ objects, some object has score at least $T/N$ and some object has score at most $T/N$. If an object is chosen at random and the expected number of ``bad events'' is less than $1$, some object has no bad event.
\end{theorem}

\begin{proof}
If every score were less than $T/N$ the total would be less than $T$. The second statement is the first applied to the number of bad events, which is a non-negative integer, so an object with fewer than $1$ of them has none.
\end{proof}

\begin{proposition}[Colouring arguments for tilings]
Suppose the cells of a board are coloured, and every tile of a given shape, however placed, covers cells whose colour counts form the same vector $t=(t_1,\dots,t_c)$, or more generally whose colour counts satisfy a fixed linear relation. If the board's colour counts $(b_1,\dots,b_c)$ are not a non-negative integer multiple of $t$ (or fail the relation), the board cannot be tiled by that shape. The argument proves only impossibility; a possible case still needs an explicit tiling.
\end{proposition}

\begin{proof}
A tiling with $N$ tiles covers each cell once, so the colour counts of the board equal the sum over tiles of their colour counts, which is $Nt$.
\end{proof}

The standard colourings: the chessboard colouring, under which a domino covers one cell of each colour; the colouring of column $j$ by $j\bmod k$, under which a $1\times k$ bar covers every colour once but a $k\times 1$ bar covers $k$ cells of one colour; and a colouring by $(i+j)\bmod k$ or $(i+2j)\bmod k$, chosen so that each placement of an awkward tile covers a prescribed pattern. Weighted colourings, with a number in each cell, are the same argument with the sum in place of the count.

\begin{lemma}[Cycles]
A finite graph in which every vertex has degree at least $2$ contains a cycle. A connected graph on $n$ vertices has at least $n-1$ edges, with equality exactly for trees, and any graph on $n$ vertices with at least $n$ edges contains a cycle.
\end{lemma}

\begin{proof}
Walk from any vertex, never immediately reversing; degree at least $2$ makes this always possible, and the finite graph forces a repeated vertex, which closes a cycle. For the second part, remove edges from cycles until none remain; connectivity is preserved and a connected acyclic graph has exactly $n-1$ edges (induction on $n$, removing a vertex of degree $1$, which exists by the first statement since there is no cycle and the graph is connected).
\end{proof}

\begin{firstmoves}
List every object for $n=2$ and $n=3$ and count by hand. Colour the board. Count by rows and by columns, or by objects and by properties, and compare. Translate people and relations into a graph and ask about degrees. Build a bijection to something already counted: binary strings, lattice paths, subsets. For a maximum, guess the extremal configuration from the small cases, prove the bound by double counting or averaging, and present the construction separately. For a parity claim, look for an involution without fixed points.
\end{firstmoves}

\begin{example}
For a nonempty subset $S$ of $\{1,\dots,n\}$, let $w(S)=\max S-\min S$. Find $\sum_S w(S)$ over all nonempty subsets $S$.
\end{example}

\begin{proof}[Solution]
Count pairs $(S,k)$ with $S$ nonempty and $\min S\le k<\max S$, where $k$ is an integer. Grouped by $S$, there are $w(S)$ such $k$ for each $S$, so the number of pairs is the sum we want. Grouped by $k$, with $1\le k\le n-1$, the subsets $S$ that qualify are those containing at least one element in $\{1,\dots,k\}$ and at least one in $\{k+1,\dots,n\}$; there are $(2^k-1)(2^{n-k}-1)$ of them. Hence
\[
\sum_S w(S)=\sum_{k=1}^{n-1}(2^k-1)(2^{n-k}-1)=\sum_{k=1}^{n-1}\bigl(2^n-2^k-2^{n-k}+1\bigr)=(n-1)2^n-2(2^n-2)+(n-1),
\]
which simplifies to $(n-3)2^n+n+3$.

Check at $n=2$: only $\{1,2\}$ contributes, with $w=1$; the formula gives $-4+5=1$. At $n=3$: the three pairs contribute $1,2,1$ and the full set contributes $2$, total $6$; the formula gives $0+6=6$. The sum could also be done as $\sum_S\max S-\sum_S\min S$ with $\sum_S\max S=\sum_kk2^{k-1}$, but the pair count did both at once and needed no identity for $\sum k2^{k-1}$.
\end{proof}

\begin{example}
Let $n\ge 1$. Find the number of pairs $(A,B)$ of subsets of $\{1,\dots,n\}$ with $A\subseteq B$ and $|B\setminus A|$ odd.
\end{example}

\begin{proof}[Solution]
Each element is in one of three states: outside $B$, in $B\setminus A$, or in $A$. So there are $3^n$ pairs in all, and the number with exactly $k$ elements in $B\setminus A$ is $\binom nk2^{n-k}$. The answer is $\sum_{k\text{ odd}}\binom nk2^{n-k}$.

Rather than sum this, pair the configurations off. Given $(A,B)$ with $A\ne\{1,\dots,n\}$, let $j$ be the smallest element not in $A$ and move it between the states ``outside $B$'' and ``in $B\setminus A$''. This changes $|B\setminus A|$ by one, so it swaps odd and even, and doing it twice returns to the start; it is an involution on the pairs with $A\ne\{1,\dots,n\}$ with no fixed points. The only unpaired configuration is $A=B=\{1,\dots,n\}$, with $|B\setminus A|=0$, even. So the number of even pairs exceeds the number of odd pairs by exactly $1$, and with $3^n$ pairs in total the number of odd pairs is $(3^n-1)/2$.

Check at $n=1$: the pairs are $(\emptyset,\emptyset)$, $(\emptyset,\{1\})$, $(\{1\},\{1\})$, and exactly one has an odd difference; $(3-1)/2=1$. At $n=2$: nine pairs, and the odd ones have $|B\setminus A|=1$, of which there are $\binom21\cdot2=4=(9-1)/2$. The binomial sum confirms it: $\sum_{k\text{ odd}}\binom nk2^{n-k}=\tfrac12\bigl((2+1)^n-(2-1)^n\bigr)$, which is the generating-function view of the same involution.
\end{proof}

\begin{trap}
Overcounting symmetric arrangements: a necklace counted as a sequence, a handshake counted from both ends, a pair counted as an ordered pair. A bijection that is not onto, or not well defined on the whole domain, which is why the inverse map should be written down. A colouring that proves impossibility and forgets the construction for the possible cases, or a construction given for one $n$ when the claim is for all $n$. Inclusion-exclusion with the signs reversed or a term missed; write the $n=2$ case to fix the pattern. Reading a graph problem as being about the specific people in it rather than the degrees. A maximum claimed from the small cases without the proof that nothing larger exists, or proved without the example that attains it.
\end{trap}

\begin{ladder}
\item \seminar{PS 10 \#1}. A square grid and an L-shaped tile; can it be tiled.
\item \seminar{PS 10 \#4}. Digit strings with a divisibility condition; how many.
\item \seminar{PS 10 \#7}. The first $n$ integers into boxes of equal sum; when possible.
\item \putnam{2018 B1}. A set of lattice vectors; remove one so the rest split into two equal-sum halves.
\item \putnam{2019 B1}. Lattice points on a family of circles; count the squares among them.
\item \putnam{2023 B1}. Coins sliding on a grid; which configurations are reachable.
\item \putnam{2024 B1}. A grid with prescribed values along a transversal; feasibility.
\item \putnam{2017 A4}. Quiz scores; split the class into two halves with equal averages.
\item \putnam{2022 B4}. Cyclic sequences in which consecutive triples form progressions; which lengths.
\item \putnam{2016 A4}. Tiling a rectangle with a given tile; the minimum number of tiles.
\item \putnam{2025 A5}. Permutations with a sign pattern of differences; the pattern maximising the count.
\item \putnam{2024 A3}. Order statistics of a random arrangement in a rectangle; a balanced pair exists.
\end{ladder}

\begin{hides}
Everywhere. A number theory problem that asks how many residues or how many integers of a form (\putnam{2020 A1}, \putnam{2022 A3}). A probability that is a count divided by a total (\putnam{2021 B1}, \putnam{2024 A3}). A linear algebra problem whose matrix is an incidence matrix (\putnam{2018 A2}, \putnam{2022 B2}). A geometry problem about lattice points (\putnam{2019 B1}, \putnam{2025 B1}). An invariant problem where the reachable set has to be counted once the invariant is found (\putnam{2023 B1}, \putnam{2024 B1}). A recurrence whose terms count something (\putnam{2020 A5}, \putnam{2024 B5}).
\end{hides}

\section{Generating functions}

\begin{recognise}
The $n$th term of a sequence depends on sums of products of earlier terms, which is a convolution. A count splits an object into an ordered sequence of parts, each part counted by a known sequence. A recurrence has coefficients that are polynomials in $n$. A sum of binomial coefficients runs over every third index, or over an arithmetic progression of indices. A problem asks about the coefficients of a power of a polynomial. The seminar's eleventh set accepts only this method, and its problems are the ladder.
\end{recognise}

\begin{definition}
The ordinary generating function of a sequence $(a_n)_{n\ge0}$ is the formal power series $A(x)=\sum_{n\ge0}a_nx^n$; the exponential generating function is $\hat A(x)=\sum_{n\ge0}a_nx^n/n!$. Formal series are added and multiplied coefficientwise and by the Cauchy product; a series with nonzero constant term has a reciprocal; no convergence is involved until a number is substituted for $x$. We write $[x^n]A$ for the coefficient of $x^n$.
\end{definition}

\begin{theorem}[Products are convolutions]
For ordinary and for exponential generating functions respectively,
\[
[x^n]\,A(x)B(x)=\sum_{k=0}^na_kb_{n-k},\qquad [x^n/n!]\,\hat A(x)\hat B(x)=\sum_{k=0}^n\binom nka_kb_{n-k}.
\]
Hence: if an object of size $n$ is an ordered pair of a first part of size $k$ counted by $a_k$ and a second of size $n-k$ counted by $b_{n-k}$, the objects are counted by $AB$; if the parts are built on labelled elements and the labels are distributed between the parts, the objects are counted by $\hat A\hat B$. An ordered sequence of any number of parts, each counted by $C$ with $c_0=0$, is counted by $1/(1-C(x))$.
\end{theorem}

\begin{proof}
Multiply out: $x^k\cdot x^{n-k}=x^n$. For the exponential version, $\frac{x^k}{k!}\cdot\frac{x^{n-k}}{(n-k)!}=\binom nk\frac{x^n}{n!}$, and $\binom nk$ is the number of ways to choose which labels go to the first part. The sequence statement is $\sum_{j\ge0}C(x)^j$, a geometric series that makes sense formally because $c_0=0$ means $C^j$ starts at $x^j$ and each coefficient is a finite sum.
\end{proof}

\begin{theorem}[Linear recurrences and rational functions]
A sequence satisfies $a_{n+k}=c_{k-1}a_{n+k-1}+\dots+c_0a_n$ for all $n\ge0$ if and only if $Q(x)A(x)$ is a polynomial of degree less than $k$, where $Q(x)=1-c_{k-1}x-\dots-c_0x^k$ is the reversed characteristic polynomial. Partial fractions then give the closed form: if $Q(x)=\prod_i(1-r_ix)^{m_i}$ then $A(x)=\sum_i\sum_{j=1}^{m_i}\dfrac{\alpha_{ij}}{(1-r_ix)^j}$ and
\[
[x^n]\frac{1}{(1-rx)^j}=\binom{n+j-1}{j-1}r^n,
\]
a polynomial in $n$ of degree $j-1$ times $r^n$.
\end{theorem}

\begin{proof}
The coefficient of $x^{n+k}$ in $Q(x)A(x)$ is $a_{n+k}-c_{k-1}a_{n+k-1}-\dots-c_0a_n$, which vanishes for all $n\ge0$ exactly when the recurrence holds; so $QA$ has no terms beyond $x^{k-1}$. The coefficient formula is the binomial series $(1-y)^{-j}=\sum_n\binom{n+j-1}{j-1}y^n$, proved by differentiating the geometric series $j-1$ times, with $y=rx$.
\end{proof}

A recurrence with polynomial coefficients, like $a_{n+1}=(n+1)a_n+\dots$, becomes a differential equation for $A$, since $xA'(x)=\sum na_nx^n$; that is when the exponential generating function usually gives the simpler equation.

\begin{theorem}[Roots-of-unity filter]
Let $\omega=e^{2\pi i/k}$ and let $A(x)$ be a polynomial, or a series converging at the points $\omega^j$. Then for any residue $r$,
\[
\sum_{n\equiv r\ (\mathrm{mod}\ k)}a_n=\frac1k\sum_{j=0}^{k-1}\omega^{-jr}A(\omega^j).
\]
\end{theorem}

\begin{proof}
$\sum_{j=0}^{k-1}\omega^{j(n-r)}$ is $k$ when $k\mid n-r$ and otherwise a geometric sum $\frac{\omega^{k(n-r)}-1}{\omega^{n-r}-1}=0$. Multiply $a_n$ by this, sum over $n$, and swap the finite sum over $j$ with the sum over $n$.
\end{proof}

For $k=2$ this is $\sum_{n\text{ even}}a_n=\frac{A(1)+A(-1)}2$. For $k=3$ with $\omega^2+\omega+1=0$ the values $A(\omega)$ and $A(\omega^2)$ are conjugate when $A$ has real coefficients, so the filter is $\frac13\bigl(A(1)+2\operatorname{Re}\,\omega^{-r}A(\omega)\bigr)$.

\begin{lemma}[Standard series]
$\dfrac1{1-x}=\sum x^n$; $\dfrac{1}{(1-x)^{m}}=\sum_n\binom{n+m-1}{m-1}x^n$; $(1+x)^n=\sum_k\binom nkx^k$; $e^x=\sum x^n/n!$ is the exponential generating function of the all-ones sequence, and $e^x-1$ of the sequence that is $0$ at $n=0$ and $1$ after; $\dfrac{1-\sqrt{1-4x}}{2x}$ generates the Catalan numbers; $\dfrac{x}{1-x-x^2}$ generates the Fibonacci numbers with $F_0=0$, $F_1=1$.
\end{lemma}

\begin{proof}
The first three are the geometric and binomial series. The exponential ones are the definition. The Fibonacci one is the previous theorem with $Q=1-x-x^2$ and $QA=x$. For Catalan, the recurrence $C_{n+1}=\sum_kC_kC_{n-k}$ says $C(x)-1=xC(x)^2$, and the quadratic formula gives the root with $C(0)=1$.
\end{proof}

\begin{firstmoves}
Write the series. Multiply the recurrence by $x^{n+k}$ and sum over $n$; collect $A(x)$ and the boundary terms. Decompose the object into parts and multiply the parts' series. Find the closed form and do partial fractions. For every $k$th coefficient, evaluate at the $k$th roots of unity. For labelled objects, use exponential generating functions. When a closed form is not available, extract what is asked (a sign, a residue, a growth rate) directly from the functional equation for $A$.
\end{firstmoves}

\begin{example}
Call a string over the alphabet $\{A,B,C\}$ tidy if every maximal run of consecutive $A$'s has even length. Let $t_n$ be the number of tidy strings of length $n$. Find a recurrence for $t_n$ and a closed form.
\end{example}

\begin{proof}[Solution]
A tidy string is an alternation: a run of $A$'s of even length (possibly empty), then a letter $B$ or $C$, then an even run of $A$'s, then a letter $B$ or $C$, and so on, ending with an even run. Even runs are counted by $E(x)=1+x^2+x^4+\dots=1/(1-x^2)$, and a letter $B$ or $C$ by $2x$. A tidy string is an even run followed by an ordered sequence of blocks, each block a letter followed by an even run, so by the convolution theorem
\[
T(x)=E(x)\cdot\frac{1}{1-2xE(x)}=\frac{1}{1-x^2}\cdot\frac{1}{1-\frac{2x}{1-x^2}}=\frac{1}{1-2x-x^2}.
\]
Reading the denominator as a reversed characteristic polynomial, $t_n=2t_{n-1}+t_{n-2}$ for $n\ge2$, with $t_0=1$ and $t_1=2$. Check: $t_2=5$, and the tidy strings of length $2$ are $AA$ and the four strings with no $A$.

The characteristic roots are $1\pm\sqrt2$, so $t_n=\alpha(1+\sqrt2)^n+\beta(1-\sqrt2)^n$; from $t_0=1$ and $t_1=2$,
\[
t_n=\frac{(1+\sqrt2)^{n+1}-(1-\sqrt2)^{n+1}}{2\sqrt2}.
\]
At $n=0$ this is $2\sqrt2/(2\sqrt2)=1$ and at $n=1$ it is $4\sqrt2/(2\sqrt2)=2$. The recurrence has a direct explanation too: a tidy string ends in $B$ or $C$ preceded by a tidy string of length $n-1$, or in $AA$ preceded by a tidy string of length $n-2$; the generating function found it without that insight, and it also proves the decomposition is exhaustive and unambiguous.
\end{proof}

\begin{example}
Let $n\ge1$. Find the number of pairs $(A,B)$ of subsets of $\{1,\dots,n\}$ with $A\subseteq B$ and $|B\setminus A|$ divisible by $3$.
\end{example}

\begin{proof}[Solution]
As in the counting section, each element is outside $B$, in $B\setminus A$, or in $A$, and marking the middle state with $x$ gives the polynomial $(2+x)^n=\sum_k\binom nk2^{n-k}x^k$, whose coefficient of $x^k$ counts pairs with $|B\setminus A|=k$. The filter with $k=3$, $r=0$ and $\omega=e^{2\pi i/3}$ gives
\[
N=\frac13\Bigl(3^n+(2+\omega)^n+(2+\omega^2)^n\Bigr).
\]
Now $2+\omega=\tfrac32+\tfrac{\sqrt3}2i=\sqrt3\,e^{i\pi/6}$ and $2+\omega^2$ is its conjugate, so
\[
N=\frac13\Bigl(3^n+2\cdot3^{n/2}\cos\frac{n\pi}6\Bigr).
\]
Check at $n=1$: $\frac13(3+2\sqrt3\cdot\frac{\sqrt3}2)=2$, the pairs $(\emptyset,\emptyset)$ and $(\{1\},\{1\})$. At $n=2$: $\frac13(9+2\cdot3\cdot\frac12)=4$, the four pairs with $A=B$, since $|B\setminus A|=3$ is impossible and $|B\setminus A|\in\{1,2\}$ is excluded. At $n=3$: $\frac13(27+0)=9$, which is the $8$ pairs with $A=B$ plus $(\emptyset,\{1,2,3\})$.

The parity version, $|B\setminus A|$ odd, is the same filter with $k=2$ and gives $\frac12(3^n-1)$; the involution in the counting section proved that without complex numbers, and no such involution is in sight for divisibility by $3$, which is what the roots of unity are for.
\end{proof}

\begin{trap}
Convergence is irrelevant for a formal power series and relevant the moment a number is substituted: the roots-of-unity filter on an infinite series needs the series to converge at those points, and a closed form like $1/(1-2x)$ represents the sequence $2^n$ for $|x|<\tfrac12$ and nothing at $x=1$. A recurrence that starts late, so the first few coefficients need their own correction terms in $QA$; always compute $QA$ to the order $x^{k-1}$ explicitly rather than assuming it is $a_0$. Ordinary and exponential series mixed in one product. Partial fractions with a repeated root done as if the root were simple. A decomposition into parts that is not unique, so the product overcounts; the tidy-string decomposition was unambiguous because every block ends with the first non-$A$ letter.
\end{trap}

\begin{ladder}
\item \seminar{PS 11 \#2}. A binomial-type identity for falling factorials; prove it.
\item \seminar{PS 11 \#3}. Subsets counted by size modulo three; a closed form.
\item \seminar{PS 11 \#7}. Coefficients of a power of a short polynomial; a property of them.
\item \putnam{2017 B3}. A power series with zero-one coefficients built from a recurrence; irrationality of its value.
\item \putnam{2023 A5}. A signed sum indexed by digit counts; its complex roots.
\item \putnam{2020 B4}. A weighted average over lattice paths; evaluate.
\item \putnam{2024 A6}. A matrix built from the coefficients of a series; its determinant.
\end{ladder}

\begin{hides}
Inside recurrences and counting, as the method that does the bookkeeping (\putnam{2020 A5}, \putnam{2024 B5}, \seminar{PS 4 \#12}). Inside number theory, where coefficients are read modulo a prime or a power of a polynomial is expanded modulo $2$ (\putnam{2022 B1}, \seminar{PS 9 \#2}). Inside probability, where the generating function of a distribution gives its moments and the probability of a sum (\putnam{2020 B4}, \seminar{PS 6 \#12}). Inside sums, where a sum of binomial coefficients with a pattern is a coefficient or a value of a known series (\putnam{2020 A2}). Inside linear algebra, where Hankel determinants of a series encode a continued fraction (\putnam{2024 A6}, \seminar{PS 7 \#24}).
\end{hides}

\section{Abstract algebra}

\begin{recognise}
A set with an operation and a few rules is given, and something must be deduced from the rules alone. Functions that respect a structure are asked about: homomorphisms, automorphisms, permutations with a property. A number theory problem is really about the multiplicative group modulo $n$, or about the order of an element in it. A problem about a finite set and a composition law asks whether an identity or inverses exist. The eighth seminar set is the only one that prints its Putnam sources, so its first fifteen problems are also a decade-by-decade ladder of what the exam has asked.
\end{recognise}

\begin{definition}
A group is a set $G$ with an associative operation, an identity $e$, and an inverse for every element. The order of $G$ is $|G|$; the order $\ord(g)$ of an element is the least $n\ge1$ with $g^n=e$. A subgroup $H$ is a subset closed under the operation and inverses; its left cosets are the sets $gH$. A homomorphism is a map $\varphi$ with $\varphi(ab)=\varphi(a)\varphi(b)$. The group is cyclic if some element generates it, that is, $G=\{g^k\}$.
\end{definition}

\begin{theorem}[Lagrange]
If $H$ is a subgroup of a finite group $G$ then $|H|$ divides $|G|$. Consequently $\ord(g)$ divides $|G|$ for every $g$, and $g^{|G|}=e$.
\end{theorem}

\begin{proof}
Two cosets $aH$, $bH$ are either equal or disjoint: if $ah_1=bh_2$ then $a=bh_2h_1^{-1}\in bH$, so $aH\subseteq bH$, and symmetrically. Every $g$ lies in $gH$. The map $h\mapsto gh$ is a bijection $H\to gH$, so all cosets have $|H|$ elements, and $|G|$ is $|H|$ times the number of cosets. For the consequences, the powers of $g$ form a subgroup of size $\ord(g)$: they are closed, and $g^n=e$ makes $g^{n-1}$ the inverse of $g$, with $e,g,\dots,g^{n-1}$ distinct by minimality of $n$.
\end{proof}

\begin{proposition}
A group of prime order $p$ is cyclic, and every non-identity element generates it. In the group of units modulo $n$, Lagrange's theorem is Euler's theorem $a^{\varphi(n)}\equiv1$ for $\gcd(a,n)=1$, and for prime $n$ it is Fermat's little theorem.
\end{proposition}

\begin{proof}
The order of $g\ne e$ divides $p$ and is not $1$, so it is $p$, and the $p$ distinct powers of $g$ exhaust the group. The units modulo $n$ form a group of order $\varphi(n)$ under multiplication, since a product of units is a unit and every unit has an inverse by Bézout.
\end{proof}

\begin{lemma}
For every positive integer $n$, $\sum_{d\mid n}\varphi(d)=n$.
\end{lemma}

\begin{proof}
Sort the fractions $1/n,2/n,\dots,n/n$ by their denominator in lowest terms. The fractions with denominator $d$ are the $k/d$ with $\gcd(k,d)=1$ and $1\le k\le d$, of which there are $\varphi(d)$, and each $d$ dividing $n$ occurs.
\end{proof}

\begin{theorem}[The multiplicative group modulo a prime is cyclic]
For a prime $p$, the group $(\Z/p)^{\times}$ of nonzero residues under multiplication is cyclic of order $p-1$. A generator is called a primitive root. More generally every finite subgroup of the multiplicative group of a field is cyclic.
\end{theorem}

\begin{proof}[Proof sketch]
Let $G$ be the group, of order $n=p-1$. For each $d\mid n$ let $N(d)$ be the number of elements of order exactly $d$. If $N(d)>0$, take $g$ of order $d$; the $d$ powers of $g$ are distinct roots of $x^d-1$, and a polynomial of degree $d$ over a field has at most $d$ roots, so every element of order dividing $d$ is a power of $g$. Among the powers $g^k$, $0\le k<d$, the ones of order exactly $d$ are those with $\gcd(k,d)=1$, so $N(d)=\varphi(d)$. Hence $N(d)\in\{0,\varphi(d)\}$ for every $d\mid n$. Since every element has some order dividing $n$, $\sum_{d\mid n}N(d)=n=\sum_{d\mid n}\varphi(d)$ by the lemma, and this forces $N(d)=\varphi(d)$ for all $d$, in particular $N(n)=\varphi(n)>0$.
\end{proof}

What the theorem buys: with a primitive root $g$, every nonzero residue is $g^k$ for a unique $k$ modulo $p-1$, so multiplicative questions modulo $p$ become additive questions modulo $p-1$. The squares are the even powers, exactly half the nonzero residues; $x^m\equiv1$ has exactly $\gcd(m,p-1)$ solutions; $x\mapsto x^m$ is a bijection of the nonzero residues if and only if $\gcd(m,p-1)=1$; and $-1$ is $g^{(p-1)/2}$, the unique element of order $2$.

\begin{proposition}[Cancellation in finite structures]
Let $S$ be a finite nonempty set with an associative operation in which both cancellation laws hold: $ax=ay$ implies $x=y$, and $xa=ya$ implies $x=y$. Then $S$ is a group. In any finite group, left multiplication by a fixed element is a permutation of the group.
\end{proposition}

\begin{proof}
For fixed $a$, left multiplication $L_a\colon x\mapsto ax$ is injective by cancellation, hence a bijection of the finite set $S$; likewise right multiplication $R_a$. So there is $e$ with $ae=a$. For any $b$, surjectivity of $R_a$ gives $c$ with $ca=b$, and then $be=cae=ca=b$: $e$ is a right identity for all of $S$. Symmetrically there is a left identity $e'$, and $e'=e'e=e$. Surjectivity of $L_a$ and $R_a$ gives $a'$, $a''$ with $aa'=e=a''a$, and $a'=ea'=a''aa'=a''e=a''$, a two-sided inverse.
\end{proof}

\begin{proposition}[Homomorphisms and orders]
A homomorphism is determined by its values on a generating set. A homomorphism sends an element of order $n$ to an element whose order divides $n$. In a permutation, the order is the least common multiple of the cycle lengths, and the sign is $(-1)^{\text{(number of even-length cycles)}}$.
\end{proposition}

\begin{proof}
Every element is a product of generators and their inverses, and $\varphi$ respects products and inverses. $g^n=e$ gives $\varphi(g)^n=\varphi(e)=e$. Disjoint cycles commute, a $k$-cycle has order $k$, and a $k$-cycle is a product of $k-1$ transpositions.
\end{proof}

\begin{firstmoves}
Identify the structure as one you know in disguise: the integers modulo $n$, the units modulo $n$, the permutations of a small set, the invertible $2\times2$ matrices. Look for the identity and for inverses, and check whether the rules given actually provide them. Count elements by their orders, using Lagrange to restrict the possible orders. Write the operation table for the smallest case. Translate the question modulo $p$ into exponents of a primitive root. For a bare binary operation, test associativity and commutativity on three symbols before assuming either.
\end{firstmoves}

\begin{example}
Let $G$ be a finite group whose order is not divisible by $3$. Prove that every element of $G$ has exactly one cube root in $G$: for each $g$ there is a unique $x\in G$ with $x^3=g$.
\end{example}

\begin{proof}[Solution]
Let $n=|G|$. Since $\gcd(3,n)=1$, there is an integer $k\ge1$ with $3k\equiv1\pmod n$. By Lagrange, $y^n=e$ for every $y\in G$, so $y^{3k}=y^{3k-1}\cdot y=y$, because $n$ divides $3k-1$. Hence the maps $c(x)=x^3$ and $r(y)=y^k$ satisfy $r(c(x))=x^{3k}=x$ and $c(r(y))=y^{3k}=y$ for all $x,y$: they are mutually inverse bijections of $G$. So $x=g^k$ is a cube root of $g$, and it is the only one, since $c$ is injective.

Nothing was assumed about commutativity. The point to carry away is the shape of the argument: a power map on a finite group is a bijection whenever the exponent is coprime to the order, and its inverse is another power map. Modulo a prime $p$ this says that $x\mapsto x^3$ permutes the nonzero residues exactly when $p\not\equiv1\pmod3$, which is the way ``every residue is a cube modulo $p$'' is proved. When $3$ does divide $n$, Cauchy's theorem gives an element of order $3$, so $x^3=e$ has at least three solutions and uniqueness fails.
\end{proof}

\begin{example}
Let $p$ be an odd prime. Show that the number of pairs $(a,b)$ of residues modulo $p$ with $a^2\equiv b^3\pmod p$ is exactly $p$.
\end{example}

\begin{proof}[Solution]
Count the pairs by $b$. If $b\equiv0$, then $a^2\equiv0$ forces $a\equiv0$, one pair. If $b\not\equiv0$, write $b=g^k$ for a primitive root $g$; then $b^3=g^{3k}$, and a nonzero residue $g^m$ is a square if and only if $m$ is even modulo $p-1$ (the squares are exactly the even powers of $g$, because $p-1$ is even). Since $3k$ is even if and only if $k$ is, $b^3$ is a square exactly when $b$ is. When $b^3$ is a nonzero square it has exactly two square roots $\pm a_0$, which are distinct because $p$ is odd; when it is not a square it has none.

There are $(p-1)/2$ nonzero squares among the residues, so the number of pairs is
\[
1+2\cdot\frac{p-1}{2}+0\cdot\frac{p-1}{2}=p .
\]
Check at $p=3$: the pairs with $a^2\equiv b^3$ are $(0,0)$, $(1,1)$, $(2,1)$, three of them. At $p=5$: $b=1$ gives $a=\pm1$, $b=4$ gives $b^3=64\equiv4$ and $a=\pm2$, while $b=2,3$ give $b^3\equiv3,2$, non-squares; with $(0,0)$ that is five pairs.

The primitive root turned a question about cubes and squares into a question about the parity of an exponent, and every multiplicative fact about residues modulo $p$ that this document uses goes through the same translation.
\end{proof}

\begin{trap}
Assuming commutativity: $(ab)^2=a^2b^2$ fails in every non-abelian group, and an argument that rearranges a product has used it. Assuming associativity when the problem gives a bare binary operation; the first problem of the seminar's set is a commutative operation that is not associative. Dividing by an element that may have no inverse, which modulo a composite $n$ is any element sharing a factor with $n$. Using a primitive root modulo a composite modulus; the units modulo $8$ form a group in which every element squares to $1$, so no element generates it. Lagrange gives divisibility only; a divisor of the order need not be the order of any element, and need not even be the order of a subgroup (the alternating group on four letters has order $12$ and no subgroup of order $6$). Concluding $x=e$ from $x^m=e$ without knowing $\gcd(m,\ord x)$.
\end{trap}

\begin{ladder}
\item \seminar{PS 8 \#1}. A binary operation on a set with a given rule; commutative but not associative.
\item \seminar{PS 8 \#2}. A group with two relations between two elements; the elements are trivial.
\item \seminar{PS 8 \#3}. A finite group; the number of elements of a given prime order.
\item \putnam{2016 A5}. A finite group; every element is a short word in two generators.
\item \putnam{2018 A4}. A group identity with floor functions in the exponents; two elements commute.
\item \putnam{2023 B5}. A permutation of the residues whose square is multiplication by a constant; which moduli.
\end{ladder}

\begin{hides}
Inside congruences, where the order of an element and the primitive root do the work (\putnam{2021 A5}, \putnam{2024 A4}, \putnam{2023 B5}, \seminar{PS 9 \#7}, \seminar{PS 9 \#32}). Inside permutations and counting, where the cycle structure or the sign is the invariant (\putnam{2025 A5}, \seminar{PS 10 \#16}). Inside linear algebra, where a set of matrices closed under multiplication is a finite group and Lagrange applies to it (\putnam{2022 B2}, \seminar{PS 7 \#14}). Inside games and invariants, where the moves generate a group acting on the positions (\putnam{2023 B1}).
\end{hides}

\section{Functional equations}

\begin{recognise}
A function is defined by an equation it satisfies for all inputs, and the problem asks to find all such functions, or to prove one property of them, or to compute one value. A function is characterised by a relation between its values at related points: $x$ and $2x$, $x$ and $1/x$, $x$ and $f(x)$. Sometimes the function is a sequence, a polynomial or a function on the integers, and the same methods apply with induction in place of continuity.
\end{recognise}

\begin{definition}
$f\colon X\to Y$ is injective if $f(a)=f(b)$ implies $a=b$, surjective if every $y\in Y$ is $f(x)$ for some $x$, and bijective if both. A function $f\colon\R\to\R$ (or $\Q\to\Q$) is additive if $f(x+y)=f(x)+f(y)$ for all $x,y$.
\end{definition}

\begin{theorem}[Cauchy's equation]
Let $f$ be additive. Then $f(qx)=qf(x)$ for every rational $q$ and real $x$; in particular $f(q)=cq$ on $\Q$ with $c=f(1)$. If $f\colon\R\to\R$ is additive and also continuous at one point, or monotone on an interval, or bounded on an interval, then $f(x)=cx$ for all real $x$.
\end{theorem}

\begin{proof}
$f(0)=f(0)+f(0)$ gives $f(0)=0$, then $f(-x)=-f(x)$. Induction gives $f(nx)=nf(x)$ for positive integers $n$, so $f(x)=f(n\cdot x/n)=nf(x/n)$ and $f(x/n)=f(x)/n$; combining, $f(mx/n)=(m/n)f(x)$, and negatives follow from oddness.

Continuity at one point $x_0$ gives continuity everywhere: $f(x+h)-f(x)=f(x_0+h)-f(x_0)$. Then for real $x$ take rationals $q_n\to x$; $f(x)=\lim f(q_n)=\lim cq_n=cx$. If $f$ is monotone, say non-decreasing, then for rationals $q<x<q'$, $cq\le f(x)\le cq'$, and letting $q,q'\to x$ gives $f(x)=cx$ (if $c<0$ the inequalities reverse, and $c=0$ makes $f$ constant on $\Q$ hence everywhere by monotonicity). If $f$ is bounded on an interval $J$, then $g(x)=f(x)-cx$ is additive, vanishes on $\Q$, and is bounded on $J$ by some $M'$; for any real $x$ there is a rational $q$ with $x-q\in J$, and $g(x)=g(x-q)$, so $|g|\le M'$ on all of $\R$; but $g(nx)=ng(x)$ is unbounded unless $g(x)=0$.
\end{proof}

\begin{proposition}[Relatives of Cauchy's equation]
With continuity or monotonicity assumed:
\begin{enumerate}[nosep]
\item $f(x+y)=f(x)f(y)$ on $\R$ gives $f\equiv0$ or $f(x)=e^{cx}$;
\item $f(xy)=f(x)+f(y)$ on $(0,\infty)$ gives $f(x)=c\ln x$;
\item $f(xy)=f(x)f(y)$ on $(0,\infty)$ gives $f\equiv0$ or $f(x)=x^{c}$;
\item $f(x+y)+f(x-y)=2f(x)+2f(y)$ on $\R$ gives $f(x)=cx^2$.
\end{enumerate}
\end{proposition}

\begin{proof}
(1) If $f(a)=0$ for some $a$ then $f(x)=f(a)f(x-a)=0$ for all $x$. Otherwise $f(x)=f(x/2)^2>0$ and $\ln f$ is additive. (2) $f(e^t)$ is additive in $t$. (3) $\ln f(e^t)$ is additive in $t$ after ruling out zeros as in (1) and noting $f(x)=f(\sqrt x)^2\ge0$. (4) $x=y=0$ gives $f(0)=0$; $x=0$ gives $f$ even; induction on $n$ gives $f(nx)=n^2f(x)$, so $f(q)=cq^2$ on $\Q$ with $c=f(1)$, and continuity extends it.
\end{proof}

The pattern in every case is the same: a substitution that produces additivity, the rationals handled by induction, and the reals by continuity. Without the regularity hypothesis there are other, wild solutions, so a problem that gives no continuity, monotonicity or boundedness either lives on $\Q$ or $\Z$, or contains a hidden regularity (a positivity, a bound, an inequality) that must be extracted from the equation itself.

\begin{proposition}[Injectivity and surjectivity from the equation]
Let $f$ satisfy a functional equation.
\begin{enumerate}[nosep]
\item If some substitution yields $f(\text{expression in }y)=g(y)$ with $g$ injective, then $f$ is injective on the range of that expression; if $g$ is surjective onto a set $T$, then the range of $f$ contains $T$.
\item If $f$ is injective, then any two expressions whose images under $f$ are equal are themselves equal, which turns one equation into a simpler one.
\item If $f$ is surjective, a variable $f(y)$ may be replaced by a free variable $a$ ranging over the whole codomain.
\end{enumerate}
\end{proposition}

\begin{proof}
(1) If $f(E(y_1))=f(E(y_2))$ then $g(y_1)=g(y_2)$ and $y_1=y_2$; for the range, each $t\in T$ is $g(y)=f(E(y))$. (2) and (3) are the definitions.
\end{proof}

\begin{proposition}[Polynomial and integer cases]
If $f$ is a polynomial, compare degrees and leading coefficients on both sides of the equation first; a polynomial identity holding at infinitely many points holds everywhere. If $f\colon\Z\to\Z$ or $\N\to\N$, compute $f(0)$, $f(1)$, $f(2)$ from the equation, then prove the guessed formula by strong induction, using the equation to express $f(n)$ through smaller values.
\end{proposition}

\begin{proof}
A nonzero polynomial has finitely many roots, so two polynomials agreeing at infinitely many points are equal; the degree comparison is the comparison of the leading terms of equal polynomials.
\end{proof}

\begin{firstmoves}
Substitute $0$, $1$, $x=y$, $y=-x$, $y=1/x$, $y=f(x)$, and whatever makes an inner argument vanish or makes two terms cancel. Determine $f(0)$ and $f(1)$. Guess the form from the special values: constant, linear, a power, an exponential, a quadratic. Ask whether the equation forces injectivity or surjectivity, and use them to simplify. Subtract the guessed solution from $f$ and see what equation the difference satisfies, which is often Cauchy's. On $\Q$ or $\Z$, induct. At the end, verify that the candidate satisfies the original equation, and say so.
\end{firstmoves}

\begin{example}
Find all functions $f\colon\Q\to\Q$ such that $f(1)=1$ and
\[
f(x+y)=f(x)+f(y)+2xy+1\qquad\text{for all }x,y\in\Q.
\]
\end{example}

\begin{proof}[Solution]
Guess the form. With $x=y=0$, $f(0)=2f(0)+1$, so $f(0)=-1$. With $x=y=1$, $f(2)=2+2+1=5$. With $x=2,y=1$, $f(3)=5+1+4+1=11$. The values $-1,1,5,11$ at $0,1,2,3$ have second differences $2$, so a quadratic with leading coefficient $1$ fits: $x^2+x-1$ takes exactly these values.

Subtract it. Let $g(x)=f(x)-(x^2+x-1)$. Then
\[
g(x+y)=f(x)+f(y)+2xy+1-(x+y)^2-(x+y)+1=g(x)+g(y),
\]
after expanding $(x+y)^2=x^2+2xy+y^2$ and cancelling. So $g$ is additive on $\Q$, and by Cauchy's theorem $g(q)=cq$ with $c=g(1)=f(1)-1=0$. Hence $g\equiv0$ and $f(x)=x^2+x-1$ for all rational $x$.

Verification: $f(x)+f(y)+2xy+1=x^2+x+y^2+y-2+2xy+1=(x+y)^2+(x+y)-1=f(x+y)$. On $\R$ without a continuity hypothesis the same argument gives $f(x)=x^2+x-1+g(x)$ with $g$ any additive function vanishing at $1$, and there are uncountably many of those; the problem lives on $\Q$ precisely so that additivity alone finishes it.
\end{proof}

\begin{example}
Show that there is no function $f\colon\R\to\R$ such that
\[
f(x+f(y))=f(x)+y^3\qquad\text{for all }x,y\in\R.
\]
\end{example}

\begin{proof}[Solution]
Suppose $f$ exists. Put $x=0$: $f(f(y))=f(0)+y^3$. The right side is an injective function of $y$ that takes every real value, so $f$ is injective (if $f(y_1)=f(y_2)$ then $f(f(y_1))=f(f(y_2))$, so $y_1^3=y_2^3$) and surjective (every real is $f(f(y))$ for some $y$).

Use surjectivity to free the inner variable. For any real $a$ choose $y$ with $f(y)=a$; then $y^3=f(f(y))-f(0)=f(a)-f(0)$, and the equation reads $f(x+a)=f(x)+f(a)-f(0)$ for all $x,a$. So $g(x)=f(x)-f(0)$ is additive: $g(x+a)=g(x)+g(a)$. In particular $g(0)=0$ and $g(2t)=2g(t)$.

Rewrite $f(f(y))=f(0)+y^3$ in terms of $g$: $f(f(y))=g(f(y))+f(0)=g\bigl(g(y)+f(0)\bigr)+f(0)=g(g(y))+g(f(0))+f(0)$, so $g(g(y))+g(f(0))=y^3$. At $y=0$ this gives $g(f(0))=0$, hence $g(g(y))=y^3$ for all $y$. But then
\[
8y^3=g(g(2y))=g(2g(y))=2g(g(y))=2y^3,
\]
so $y^3=0$ for every $y$, which is false at $y=1$. No such $f$ exists.

Where the cube mattered: with $y$ in place of $y^3$ the same steps give $g(g(y))=y$, which is consistent with additivity, and $f(x)=x$ is a solution. The contradiction came from comparing how the two sides scale under $y\mapsto2y$, a homogeneity check that is worth running on any equation whose two sides have different degrees.
\end{proof}

\begin{trap}
Proving that the guess satisfies the equation and forgetting to prove it is the only solution, or the reverse; both halves are needed and the write-up should name them. Dividing by $f(\text{something})$ that might be zero; treat the zero case separately. Using continuity you were not given, and conversely, not noticing a regularity the equation does give (a square, a positivity, a bound on an interval), which is what makes Cauchy's theorem applicable. Substituting $y=1/x$ at $x=0$. Concluding injectivity from $f(a)=f(b)\Rightarrow a=b$ checked only for one special pair. Claiming a function on $\R$ is determined by its values on $\Q$ without a reason. Reporting ``$f(x)=cx$'' when the equation forces a specific $c$, or the reverse.
\end{trap}

\begin{ladder}
\item \seminar{PS 3 \#8}. Functions satisfying the addition laws of sine and cosine; identify them.
\item \putnam{2022 B6}. A continuous positive function with an equation linking $f$ at related points; find all.
\item \putnam{2016 A3}. A function with a self-referential relation involving $1/x$; compute an integral.
\item \putnam{2016 B5}. Functions preserving a constraint between powers and intervals; find all.
\item \putnam{2025 B6}. A function with an iterated growth condition; the largest admissible exponent.
\end{ladder}

\begin{hides}
Inside analysis, where the unknown function is pinned down by its equation and then integrated (\putnam{2016 A3}, \seminar{PS 3 \#28}, \seminar{PS 3 \#38}). Inside inequalities, where a functional inequality rather than an equation is given and the growth rate is the question (\putnam{2025 B6}, \seminar{PS 5 \#15}). Inside polynomials, where a polynomial identity in $x$ is a functional equation solved by degree comparison (\putnam{2024 A2}, \seminar{PS 4 \#13}). Inside recurrences, where a sequence satisfying a relation at $n$ and $2n$ is a function on the integers (\seminar{PS 3 \#26}). Inside number theory, where a function on the positive integers is multiplicative or satisfies a divisibility relation (\putnam{2018 B3}).
\end{hides}

\section{Geometry}

\begin{recognise}
Points, lines, circles, triangles, polygons or lattice points appear. A figure is built in steps (erect a square, take a midpoint, rotate) and a property of the result is asked for. A set of points in the plane or in space has a metric condition and a count or an existence is asked. A quantity is claimed not to depend on a point that was chosen freely. The seminar has no geometry set, but six of the forty easy Putnam slots in the last decade were geometry, so the tools below have to be at hand without a reference set to practise on.
\end{recognise}

\begin{proposition}[Coordinates]
Place the origin at the figure's natural centre (a midpoint, a circumcentre, a vertex with a right angle) and an axis along a natural line. Then the signed area of the triangle with vertices $A,B,C$ is $\tfrac12\det(B-A,\ C-A)=\tfrac12\bigl((b_1-a_1)(c_2-a_2)-(b_2-a_2)(c_1-a_1)\bigr)$, positive when $ABC$ is counterclockwise; three points are collinear exactly when this is $0$; and the area of a polygon with vertices $P_1,\dots,P_n$ in order is $\tfrac12\sum_i\det(P_i,P_{i+1})$ with $P_{n+1}=P_1$.
\end{proposition}

\begin{proof}
The determinant is the area of the parallelogram spanned by two vectors, which is base times height, and the triangle is half of it. The polygon formula is the triangle formula applied to the triangles $OP_iP_{i+1}$ and summed; the parts outside the polygon cancel in pairs by sign.
\end{proof}

\begin{proposition}[Complex numbers for rotations]
Identify the plane with $\C$. The rotation by angle $\theta$ about the point $c$ is $z\mapsto c+e^{i\theta}(z-c)$. The composition of a rotation by $\alpha$ and a rotation by $\beta$ is a rotation by $\alpha+\beta$ when $\alpha+\beta$ is not a multiple of $2\pi$, and a translation otherwise. The angle from the ray $AB$ to the ray $AC$ is the argument of $(c-a)/(b-a)$; this ratio is real exactly when $A,B,C$ are collinear. The vertices of a regular $n$-gon centred at $0$ are $r e^{i\varphi}\omega^k$ with $\omega=e^{2\pi i/n}$, and $1+\omega+\dots+\omega^{n-1}=0$. The triangle $abc$ is equilateral, with $a,b,c$ counterclockwise, if and only if $a+\omega b+\omega^2c=0$ for $\omega=e^{2\pi i/3}$.
\end{proposition}

\begin{proof}
A rotation about $0$ is multiplication by $e^{i\theta}$, and conjugating by the translation $z\mapsto z-c$ moves the centre. Composing $z\mapsto c_1+e^{i\alpha}(z-c_1)$ with $z\mapsto c_2+e^{i\beta}(z-c_2)$ gives $z\mapsto e^{i(\alpha+\beta)}z+\text{constant}$, which is a rotation about the fixed point $z_0=\text{constant}/(1-e^{i(\alpha+\beta)})$ when that denominator is nonzero. The angle statement is the polar form of a quotient. The roots of unity sum to $0$ because they are the roots of $x^n-1$, whose $x^{n-1}$ coefficient is $0$. For the last claim, $1+\omega+\omega^2=0$ gives $a+\omega b+\omega^2c=(a-c)+\omega(b-c)$, which vanishes exactly when $a-c=-\omega(b-c)=e^{-i\pi/3}(b-c)$, that is, when $A$ is obtained from $B$ by rotating about $C$ through $-60^\circ$.
\end{proof}

\begin{proposition}[Vectors and the classical centres]
Let $A,B,C$ be the position vectors of the vertices of a triangle.
\begin{enumerate}[nosep]
\item The centroid is $G=\tfrac13(A+B+C)$; it lies on each median and divides it in the ratio $2:1$ from the vertex.
\item If the origin is the circumcentre $O$, the orthocentre is $H=A+B+C$, so $O$, $G$, $H$ are collinear with $OH=3\,OG$.
\item The incentre is $I=(aA+bB+cC)/(a+b+c)$, where $a,b,c$ are the side lengths opposite $A,B,C$; it is where the angle bisectors meet and is equidistant from the three sides.
\item (Inscribed angle) An angle inscribed in a circle is half the central angle on the same arc; in particular, angles on the same arc are equal, and an angle in a semicircle is right.
\end{enumerate}
\end{proposition}

\begin{proof}
(1) The midpoint of $BC$ is $M=\tfrac12(B+C)$, and $A+\tfrac23(M-A)=\tfrac13(A+B+C)$, the same for all three medians. (2) $(H-A)\cdot(B-C)=(B+C)\cdot(B-C)=|B|^2-|C|^2=0$ since $|A|=|B|=|C|$; so the line $AH$ is the altitude from $A$, and likewise for $B$ and $C$; then $H=3G$. (3) The bisector from $A$ meets $BC$ at the point $D$ dividing $BC$ in the ratio $c:b$ (angle bisector theorem, from the sine rule in the two sub-triangles), so $D=(bB+cC)/(b+c)$; the point $I$ above lies on $AD$ since $I=\frac{a}{a+b+c}A+\frac{b+c}{a+b+c}D$, and by symmetry on the other two bisectors. (4) Join the centre to the three points; the two isosceles triangles give the inscribed angle as the sum of two base angles, each half of its exterior angle at the centre.
\end{proof}

\begin{proposition}[Lattice points]
A square is $0$ or $1$ modulo $4$, so $x^2+y^2$ is never $3$ modulo $4$, and $x^2+y^2\equiv 2\pmod4$ forces $x,y$ both odd. The area of a polygon with lattice-point vertices is a multiple of $\tfrac12$ (Pick: area $=I+B/2-1$ with $I$ interior and $B$ boundary lattice points). There is no equilateral triangle with all three vertices at lattice points. The number of lattice points on the segment from $(0,0)$ to $(p,q)$, excluding the origin, is $\gcd(p,q)$.
\end{proposition}

\begin{proof}
Squares modulo $4$: $(2k)^2=4k^2$ and $(2k+1)^2=4k(k+1)+1$. The half-integer area follows from the determinant formula, whose terms are integers. For an equilateral lattice triangle with side $s$, $s^2$ is an integer, and the area $\tfrac{\sqrt3}4s^2$ is irrational, contradicting the half-integer area. The lattice points $(tp,tq)$ with $0<t\le1$ require $t=k/\gcd(p,q)$.
\end{proof}

\begin{proposition}[Pigeonhole in the plane and extremal points]
To show two of $n$ points in a region are within distance $d$, partition the region into fewer than $n$ pieces each of diameter at most $d$. Among finitely many points, the pair at maximum distance, the point farthest from a line, and the triangle of largest area are the natural extremal objects; the largest-area triangle $ABC$ in a finite set has every other point inside the triangle with sides parallel to those of $ABC$ through the opposite vertices (four times the area).
\end{proposition}

\begin{proof}
Two points in one piece are at most its diameter apart. For the largest triangle, a point beyond the line through $A$ parallel to $BC$ would form a triangle with $B$ and $C$ of larger height, hence larger area.
\end{proof}

\begin{firstmoves}
Draw it to scale, twice, with different shapes for the free parts of the figure, and see what stays the same. Put coordinates with the origin at the symmetric point; for a right angle, put it at the origin with legs along the axes. For rotations and regular polygons, use complex numbers. For midpoints, centroids and parallel lines, use vectors. For lattice points, reduce modulo $4$ and compute an area. Compute the smallest case or a degenerate case (a point at the centre, a triangle collapsed to a segment) to find the answer before proving it.
\end{firstmoves}

\begin{example}
Let $A$ and $B$ be distinct points in the plane. A point $P$ is rotated by $90^\circ$ counterclockwise about $A$ to give $P_1$, and $P_1$ is rotated by $90^\circ$ counterclockwise about $B$ to give $P_2$. Show that the midpoint of $PP_2$ does not depend on $P$, and describe it.
\end{example}

\begin{proof}[Solution]
Work in $\C$ with $A=a$, $B=b$, $P=z$. The first rotation gives $p_1=a+i(z-a)$ and the second
\[
p_2=b+i(p_1-b)=b+i\bigl(a+i(z-a)-b\bigr)=-z+a+b+i(a-b).
\]
So $p_2=-z+2c$ with $c=\tfrac12\bigl(a+b+i(a-b)\bigr)$, which does not involve $z$: the composite map is the half-turn about $c$, as the composition of two quarter-turns must be, and the midpoint of $PP_2$ is $\tfrac12(z+p_2)=c$ for every $P$.

To describe $c$: $c-a=\tfrac12(b-a)(1-i)$ and $c-b=\tfrac12(a-b)(1+i)$, so $|c-a|=|c-b|=|a-b|/\sqrt2$, and $(c-a)/(c-b)=-(1-i)/(1+i)=i$, so the angle at $c$ is a right angle. $C$ is the apex of the isosceles right triangle with hypotenuse $AB$. Which of the two apexes: with $a=0$, $b=1$ one gets $c=\tfrac12(1-i)$, below the real axis, so $C$ lies to the right of the directed segment from $A$ to $B$.

Check with $P=A$: $P_1=A$, and $P_2=b+i(a-b)$; the midpoint of $A$ and $P_2$ is $\tfrac12(a+b+i(a-b))=c$. The general fact behind this, that two rotations compose to a rotation by the sum of the angles about a centre determined by the two centres and angles alone, is the whole content; the complex numbers merely locate the centre.
\end{proof}

\begin{example}
Let $N$ be a positive integer that is neither a perfect square nor twice a perfect square. Show that the number of lattice points $(x,y)$ with $x^2+y^2=N$ is divisible by $8$.
\end{example}

\begin{proof}[Solution]
The eight maps $(x,y)\mapsto(\pm x,\pm y)$ and $(x,y)\mapsto(\pm y,\pm x)$ preserve $x^2+y^2$, so they permute the solutions, and they form a group (the symmetries of the square). The solutions split into orbits under this group, and the orbit of $(x,y)$ has exactly $8$ elements unless two of the eight images coincide, which happens only when $x=0$, or $y=0$, or $|x|=|y|$.

If $x=0$ or $y=0$ then $N=y^2$ or $N=x^2$ is a perfect square, excluded. If $|x|=|y|$ then $N=2x^2$ is twice a square, excluded. So under the hypothesis every orbit has exactly $8$ elements, and the number of solutions, being a disjoint union of orbits, is a multiple of $8$.

Check: $N=5$ has the eight solutions $(\pm1,\pm2)$ and $(\pm2,\pm1)$; $N=25$, a square, has twelve, namely $(\pm5,0)$, $(0,\pm5)$ and the eight from $(3,4)$; $N=50$, twice a square, has twelve, from $(\pm5,\pm5)$ and $(1,7)$. For $N\equiv3\pmod4$ the count is $0$, which is also a multiple of $8$; the parity of squares modulo $4$ and the symmetry argument are both at work, and the symmetry argument is the one that generalises: any set of lattice points invariant under a group of $8$ symmetries has size $\equiv$ (number of points with a nontrivial stabiliser) modulo $8$.
\end{proof}

\begin{trap}
A configuration drawn in a special position that the problem did not promise, so a property that holds in the drawing is read as general; redraw with the free points moved. A sign lost in a coordinate or an orientation lost in a rotation; the two apexes in the first example differ by the sign of $i$, and the problem's ``counterclockwise'' picks one. Dividing by a length that may be zero in a degenerate case the problem allows. Using the inscribed angle theorem across the wrong arc, where the angle is supplementary rather than equal. Treating ``lattice point'' as ``integer coordinates'' and forgetting that a circle or line through lattice points can carry non-lattice points too. Picking coordinates that make the algebra heavy when a different origin would make half the terms vanish.
\end{trap}

\begin{ladder}
\item \putnam{2019 A2}. A triangle in which a line through two centres is parallel to a side; an angle.
\item \putnam{2017 B1}. Two distinct lines in the plane; a point of one scaled onto the other.
\item \putnam{2019 B1}. Lattice points on a family of circles; count the squares among them.
\item \putnam{2021 B1}. A random square on a grid; the chance it avoids certain points.
\item \putnam{2024 B2}. A chain of quadrilaterals each built from the last; infinitely many non-congruent.
\item \putnam{2025 B1}. A two-colouring of the plane closed under circumcentres; a monochromatic configuration.
\item \putnam{2021 A3}. A sphere passing through the vertices of an integer regular tetrahedron; which radii.
\item \putnam{2016 B3}. A finite set of points in the plane; one bounded triangle covers it.
\end{ladder}

\begin{hides}
Inside counting, where the objects counted are lattice points or configurations in the plane (\putnam{2019 B1}, \putnam{2025 B1}, \putnam{2016 A4}). Inside probability, where the chance is a ratio of areas (\putnam{2021 B1}, \putnam{2024 A5}). Inside number theory, where a Diophantine equation is the equation of a circle or a sphere (\putnam{2021 A3}, \putnam{2018 A6}). Inside linear algebra, where lines and planes are vectors and a Gram matrix encodes the distances (\putnam{2017 B1}, \putnam{2018 A6}). Inside inequalities, where unit complex numbers or centroids of regions are compared (\putnam{2020 B5}, \putnam{2025 B2}). Inside pigeonhole and extremal arguments, where the plane is cut into small pieces or the farthest pair is taken (\putnam{2016 B3}, \putnam{2017 B5}).
\end{hides}

\section{Games}
\label{sec:games}

Two players alternate moves on a finite structure and the question is who wins, or what one player can force. The Putnam asks this rarely in the early slots and often enough later that the two standard arguments should be in hand: backward induction from the end, and symmetry.

\begin{recognise}
Two players, alternating turns, complete information, no chance. A question of the form ``who has a winning strategy'', ``show the second player can always'', ``what is the largest number the first player can guarantee''. Also, disguised: a process in which an adversary chooses the worst case at each step, which is a game against the adversary.
\end{recognise}

\subsection*{The model}

Consider a game with finitely many positions, in which the players alternate moves, the game always ends, and the player who cannot move loses. (Any game with a different ending rule can be rewritten in this form by adjusting the positions.) Positions in which the game has ended are terminal.

\begin{definition}
A position is a \emph{losing position} ($P$-position, the previous player wins) if the player about to move loses with best play on both sides, and a \emph{winning position} ($N$-position, the next player wins) otherwise.
\end{definition}

\begin{theorem}[Backward induction]
Every position of such a game is either winning or losing, and the two classes are characterised by: a terminal position is losing; a position is winning if and only if some move from it leads to a losing position; a position is losing if and only if every move from it leads to a winning position.
\end{theorem}

\begin{proof}
Induction on the length of the longest sequence of moves from the position, which is finite because the game always ends. Terminal positions, with length $0$, are losing by the rule that the player to move loses. For a position of length $\ell \ge 1$, every position reached by a move has length less than $\ell$ and so is already classified. If some move reaches a losing position, the player to move makes it and wins. If every move reaches a winning position, then whatever the player does, the opponent is in a winning position and wins. So the position is classified, and by exactly the stated rule.
\end{proof}

The theorem is also an algorithm: compute from the end. For a game on a number $n$, list the positions $0, 1, 2, \dots$ and mark each losing or winning from the marks of the positions it can move to. The pattern usually becomes periodic within a dozen positions, and the proof is then an induction that the pattern persists.

\begin{proposition}[The invariant view]
A set $L$ of positions is exactly the set of losing positions if and only if: every terminal position is in $L$; no move leads from a position in $L$ to a position in $L$; and from every position not in $L$ some move leads into $L$.
\end{proposition}

\begin{proof}
If $L$ has the three properties, then the player who faces a position not in $L$ can move into $L$, and the opponent must then move out of $L$; repeating, the first player always faces a position outside $L$ and never runs out of moves, so the opponent eventually reaches a terminal position inside $L$ and loses. Hence positions outside $L$ are winning and positions inside $L$ are losing. The converse is the theorem above.
\end{proof}

This is Section~\ref{sec:invariants} seen from the second player's side. A guessed description of the losing positions is proved by checking the three properties, and the winning player's strategy is ``move into $L$''. One never needs to show the guess was found by backward induction.

\begin{theorem}[Pairing strategy]
Suppose the moves available in a game can be partitioned into pairs $\{m, m'\}$ such that whenever one player makes $m$, the move $m'$ is available to the opponent and, after both, the pairing of the remaining moves is still valid. Then the second player wins by always answering $m$ with $m'$. If a single move $m_0$ is excluded from the pairing and all the other properties hold after $m_0$, then the first player wins by playing $m_0$ first and then answering in pairs.
\end{theorem}

\begin{proof}
The player following the pairing always has a move available, because their opponent's last move came from a pair whose other half is still available by hypothesis. A player who always has a move never loses in a game that ends.
\end{proof}

The usual pairing is a geometric symmetry: a reflection or a half-turn of the board, with the first player taking the centre when there is one. The hypothesis that must be checked, and that problems are designed to break, is that the mirror move is actually legal after the opponent's move. The second example below checks it.

\begin{firstmoves}
Compute the losing positions from the end for the smallest board, by hand, for at least ten positions or until the pattern repeats twice. Guess the description of the losing set, and prove it by the three properties, not by recomputing. For a game on a board, look for a symmetry; for a game with a central cell, let the first player take it. For a game on integers, look at residues of the position modulo a small number. For ``what can the first player guarantee'', find a strategy giving a value and a counter-strategy for the opponent capping it, which is two arguments.
\end{firstmoves}

\subsection*{Worked examples}

\begin{example}
Three piles of stones hold $a$, $b$ and $c$ stones. Two players alternate; a move chooses one, two or all three of the piles, each of which must be nonempty, and removes exactly one stone from each chosen pile. The player who cannot move loses. For which starting triples does the first player win?
\end{example}

\begin{proof}[Solution]
Claim: the losing positions are exactly the triples in which $a$, $b$ and $c$ are all even. We check the three properties of the invariant view.

The only terminal position is $(0, 0, 0)$, which is all even.

From an all-even position, every move removes one stone from at least one pile, making that pile odd. So no move leads from all-even to all-even.

From a position that is not all-even, let the odd piles be some nonempty subset of the three. Removing one stone from each odd pile is a legal move (each odd pile is nonempty) and makes every pile even. So from every position outside the set, some move leads into it.

Hence the first player wins exactly when at least one of $a$, $b$, $c$ is odd, with the strategy of always emptying the parities: take one stone from each odd pile.
\end{proof}

Backward induction on a few small triples, $(1,0,0)$, $(1,1,0)$, $(2,0,0)$, $(1,1,1)$, $(2,1,0)$, $(2,2,0)$, produces the guess in a minute; the proof then never mentions induction. The same argument works for any number of piles.

\begin{example}
A row of $n$ cells is empty, with $n$ odd. Two players alternate placing a stone in an empty cell, subject to the rule that no two stones may ever occupy adjacent cells. The player who cannot move loses. Show that the first player has a winning strategy.
\end{example}

\begin{proof}[Solution]
Number the cells $1$ to $n$ and write $c = (n+1)/2$ for the centre, which is a cell because $n$ is odd. Let $\sigma$ be the reflection $k \mapsto n + 1 - k$, which fixes $c$ and swaps the two halves.

The first player places a stone at $c$. Thereafter, whenever the second player places a stone at $k$, the first player places a stone at $\sigma(k)$. We show by induction on the number of rounds that this reply is always legal and that after each reply the set of occupied cells is symmetric under $\sigma$.

After the first move the occupied set is $\{c\}$, which is symmetric. Suppose the occupied set $S$ is symmetric and the second player legally places at $k$; then $k \notin S$, and $k \ne c$ since $c \in S$, so $\sigma(k) \ne k$. Legality of the reply at $\sigma(k)$ requires three things. First, $\sigma(k)$ is empty: $\sigma(k) \notin S$ because $k \notin S$ and $S$ is symmetric, and $\sigma(k) \ne k$. Second, no stone of $S$ is adjacent to $\sigma(k)$: a stone at $\sigma(k) \pm 1$ in $S$ would give, by symmetry, a stone at $k \mp 1$ in $S$, contradicting the legality of the second player's move. Third, the new stone at $k$ is not adjacent to $\sigma(k)$: adjacency would mean $|k - \sigma(k)| = |2k - n - 1| = 1$, so $2k = n$ or $2k = n + 2$, impossible since $n$ is odd and $2k$ is even. After the reply the occupied set is $S \cup \{k, \sigma(k)\}$, again symmetric.

So the first player always has a legal reply. The game ends, since each move fills a cell, so the second player is the one who eventually cannot move.
\end{proof}

The third check is the one that fails for even $n$: with no central cell, a second player stone next to the axis of symmetry would be adjacent to its mirror image, and the strategy breaks. The even case is a genuinely different problem, and nothing above says who wins it.

\begin{trap}
A pairing broken by a central element, or by two cells that are each other's mirror and adjacent. Check the mirror move's legality, every condition, in writing.

A strategy that assumes the game ends. The pairing theorem gives a player who never lacks a move; it is the finiteness of the game that converts this into a win.

Confusing ``the last player to move wins'' with ``the last player to move loses''. The misère version of a game has different losing positions, and the backward induction starts from a different base.

A losing-set guess verified only on the examples that produced it. The three properties must be proved for every position, and the hardest one is usually ``from outside the set there is always a move into it''.
\end{trap}

\begin{ladder}
\item \putnam{2020 B2}. Pegs moved along a line by two players; for which starting positions the first player wins. Backward induction from the end, then a residue pattern.
\item \putnam{2022 A5}. Dominoes placed alternately on a strip; how many cells one player can keep uncovered. A strategy and a counter-strategy give matching bounds.
\item \putnam{2023 A6}. Integers chosen alternately, with a parity goal; who wins for each starting size. The winning condition is a congruence.
\item \putnam{2025 A3}. A game on digit strings in which repeats are forbidden; the winner. Count the available moves and their parity.
\item \putnam{2017 A5}. A card-drawing game with three players; show each is sometimes favoured. Games meet probability.
\end{ladder}

Seminar: \seminar{PS 6 \#24} to \seminar{\#28} are the probability games, where the opponent is partly chance; \seminar{PS 10 \#8, \#13, \#14, \#18} are configuration games on an array, on stones with a divisibility rule, on the faces of a polyhedron, and on commuting matrices.

\begin{hides}
Invariants: the losing set is the level set of an invariant that the winning player restores at each turn, which is the proposition above read the other way. Counting: ``the first player can guarantee at least $k$'' is an extremal question with a construction and a bound. Probability: a game with a random element is analysed by conditioning on the first move, as in the section on probability and expectation. Any adversarial ``for every choice of the opponent'' statement, including worst-case bounds on algorithms, is a game whether or not the problem says so.
\end{hides}


\section{Sources}

Two sources carry the whole plan, and neither is reproduced here.

\begin{sloppypar}
\paragraph{The seminar.} MIT 18.A34 Mathematical Problem Solving (Putnam Seminar), Fall 2018, taught by Yufei Zhao, problem sets 1 to 11, at \url{https://ocw.mit.edu/courses/18-a34-mathematical-problem-solving-putnam-seminar-fall-2018/pages/assignments/}, published under CC BY-NC-SA 4.0. Four further sets come from the Fall 2007 edition of the same seminar, 18.S34, by Richard Stanley: pigeonhole, greatest integer, roots of polynomials, and limits. This document refers to problems by set and number only, as in \seminar{PS 9 \#14} or \seminar{Roots \#6}, and paraphrases the sets' own prefaces where they say something about difficulty. The statements are in the sets, not here.

\paragraph{The archive.} Kiran Kedlaya's Putnam archive at \url{https://kskedlaya.org/putnam-archive/}, with every paper from 1985 and solutions from 1995. A year's paper is \texttt{2024.pdf} and its solutions \texttt{2024s.pdf}. The solutions are compiled by Manjul Bhargava, Kiran Kedlaya and Lenhard Ng, who ask that they be linked to rather than copied; this document links, and quotes nothing. Problems are named as \putnam{2021 A1}.
\end{sloppypar}

\paragraph{The format.} The exam's shape, timetable and room rules are from the MAA's \emph{Policies, Rules, and Guidelines} for the 87th William Lowell Putnam Mathematical Competition, 5 December 2026, in the version updated 2026-08-31, which Adit supplied as a PDF on 2026-11-09. Its rules are restated here, not quoted. The policy says the MAA may change it and that its own website governs, so the supervisor's instructions on the day come first.

\paragraph{The tags.} The topic tags on the 120 Putnam problems of 2016 to 2025, the tag counts in Section~1, and the ladders in \secref{sec:ladders}{the four ladders} come from a model's reading of the problem statements and of the fifteen seminar sets. They were assigned by statement, not by solving the problems, so a tag is a guess at the technique and can be wrong. When a call is settled blind because the solution used a technique the tag did not suggest, the tag was wrong too, and that is worth a mark on the section.

\paragraph{The book.} \emph{Putnam and Beyond} by R\u{a}zvan Gelca and Titu Andreescu was used in one way only: to check that the technique sections above cover its parts (methods of proof, algebra, real analysis, geometry and trigonometry, number theory, combinatorics and probability). No text, problem or solution from it is reproduced. A copy of the book that circulates from a WordPress site is not the publisher's.

\paragraph{The examples.} Every worked example in this document was invented for it. None is a competition or textbook problem with its numbers changed. The standard theorems and their standard proofs are the common property of the subject.



\end{document}
